\documentclass[11pt,reqno]{amsart}

\usepackage{amsmath,amssymb,amsthm,mathtools}
\usepackage{booktabs,array,multirow,longtable,graphicx}
\usepackage{xcolor}
\usepackage[margin=1.08in]{geometry}
\usepackage[hypertexnames=false,colorlinks=true,linkcolor=blue!60!black,citecolor=blue!60!black,urlcolor=blue!60!black]{hyperref}
\urlstyle{tt}
\makeatletter
\g@addto@macro{\UrlBreaks}{\UrlOrds\do\_\do\-\do\.\do\0\do\1\do\2\do\3\do\4\do\5\do\6\do\7\do\8\do\9\do\a\do\b\do\c\do\d\do\e\do\f\do\i\do\k\do\l\do\m\do\n\do\o\do\p\do\r\do\s\do\t\do\u\do\x\do\y}
\makeatother
\usepackage{microtype}
\setlength{\emergencystretch}{1.6em}
\allowdisplaybreaks

\newtheorem{theorem}{Theorem}[section]
\newtheorem{proposition}[theorem]{Proposition}
\newtheorem{lemma}[theorem]{Lemma}
\newtheorem{corollary}[theorem]{Corollary}
\theoremstyle{definition}
\newtheorem{definition}[theorem]{Definition}
\theoremstyle{remark}
\newtheorem{remark}[theorem]{Remark}

\numberwithin{equation}{section}
\numberwithin{table}{section}

\DeclareMathOperator{\tr}{tr}
\DeclareMathOperator{\rk}{rk}
\DeclareMathOperator{\adj}{adj}
\DeclareMathOperator{\Pf}{Pf}
\DeclareMathOperator{\SRG}{SRG}
\DeclareMathOperator{\Deck}{Deck}
\DeclareMathOperator{\mult}{mult}
\DeclareMathOperator{\diag}{diag}
\DeclareMathOperator{\Gal}{Gal}
\DeclareMathOperator{\Tr}{Tr}
\DeclareMathOperator{\Stab}{Stab}
\DeclareMathOperator{\Aut}{Aut}
\DeclareMathOperator{\Sym}{Sym}
\DeclareMathOperator{\SNF}{SNF}
\DeclareMathOperator{\lcm}{lcm}
\DeclareMathOperator{\spn}{span}

\title[Equiangular lines in dimensions 18 and 19]{Equiangular lines in dimensions \texorpdfstring{$18$}{18} and \texorpdfstring{$19$}{19}: \texorpdfstring{$2$}{2}-adic Pfaffian congruences, Galois orthogonality, and exact spectral certificates}
\author{Mehardeep Singh}
\address{Sri Chaitanya Educational Institutions, India}
\email{mehardeep.sim@gmail.com}
\date{November 2026}
\subjclass[2020]{Primary 52C35, 05C50; Secondary 11C20, 05E30, 68V15}
\keywords{Equiangular lines, Seidel matrix, $2$-adic congruences, Pfaffian, Eulerian switching, Galois orthogonality, Jacobi complementary minor identity, strongly regular graphs, exact rational certificates, Lean 4}

\begin{document}

\begin{abstract}
Let $N(d)$ denote the maximum cardinality of an equiangular line system in $\mathbb{R}^d$. Prior to this work, the smallest dimensions in which $N(d)$ remained open were $d = 18$, $19$, and $20$, with known intervals $N(18) \in \{57, 58, 59\}$ and $N(19) \in \{72, 73, 74\}$. We establish two main bounds:
\begin{enumerate}
  \item Unconditionally (certified by finite exact integer and rational certificates), no system of $59$ equiangular lines exists in $\mathbb{R}^{18}$; hence $N(18) \le 58$ and $N(18) \in \{57, 58\}$.
  \item Unconditionally, a system of $74$ equiangular lines in $\mathbb{R}^{19}$ exists if and only if there exists a strongly regular graph $\SRG(75, 32, 10, 16)$. Combined with the theorem of Azarija and Marc (2018) ruling out $\SRG(75, 32, 10, 16)$, this implies $N(19) \le 73$ and hence $N(19) \in \{72, 73\}$.
\end{enumerate}
The remaining boundary cases $(18, 58)$ and $(19, 73)$ stay open.

To overcome the degree-$17$ odd-order enumeration barrier at $(18, 59)$ and the multi-eigenvalue interlacing barriers in both dimensions, we prove an exact mod-$2^{n+1}$ Pfaffian congruence $\chi_S(2z - 1) \equiv 2^{n-1}(2z - n)Q_0(z)^2 \pmod{2^{n+1}}$ for symmetric Seidel matrices of odd order $n$, derive universal mod-$4$ matrix polynomial laws in the strictly unique odd-order Eulerian switching gauge, establish a family of $\binom{e+1}{2}$ rational linear Galois orthogonality equations over degree-$e$ simple number-field factors, strengthen submatrix extraction to multiplicity-preserving vertex deletions, and prove a subsumption theorem showing that $2$-adic Hensel unit-eigenvector mod-$8$ congruences are automatically implied by Strict Type-2 divisibility. Remaining candidate spectra are ruled out via block-commutator simultaneous diagonalization, spectral projector rank and Parseval frame obstructions, De Caen--Jacobi resolvent pole cancellation, and an exact rational Moore--Penrose, $LDL^T$, and squarefree-radical decomposition incorporating a $\le 2$-copy multiplicity bound on identical core projections. Every computational step uses exact integer, rational, or algebraic-number arithmetic corroborated by an independent C++20 and Python implementation, and the algebraic theory together with all $147$ one-vertex Farkas infeasibility certificates for $(19, 74)$ is formally verified in Lean~4 (\texttt{v4.36.0}, $20{,}765$ lines, $127$ audited declarations on standard axioms).
\end{abstract}

\maketitle
\setcounter{tocdepth}{2}
\tableofcontents

% ============================================================================
\section{Introduction, main results, and literature review}\label{sec:intro}
% ============================================================================

\subsection{Problem statement and main theorems}
A set of $n$ lines through the origin in Euclidean space $\mathbb{R}^d$, spanned by unit vectors $v_1, \dots, v_n \in \mathbb{R}^d$, is called an \emph{equiangular line system} with common angle $\arccos\alpha$ ($0 \le \alpha < 1$) if $|\langle v_i, v_j \rangle| = \alpha$ for all $1 \le i < j \le n$. Writing $N_\alpha(d)$ for the maximum cardinality of an equiangular system in $\mathbb{R}^d$ with common inner product $\alpha$, the fundamental quantity is
\[
  N(d) := \max_{0 \le \alpha < 1} N_\alpha(d).
\]
Initiated by Haantjes~\cite{Haantjes1948}, van Lint and Seidel~\cite{vanLintSeidel1966}, and Lemmens and Seidel~\cite{LemmensSeidel1973}, the study of $N(d)$ splits into two regimes. In high dimensions with fixed $\alpha$, breakthroughs of Bukh~\cite{Bukh2016}, Balla, Dr\"axler, Keevash, and Sudakov~\cite{BDKS2018}, and Jiang, Tidor, Yao, Zhang, and Zhao~\cite{JTYZZ2021} determined the exact asymptotic slope $\lim_{d\to\infty} N_\alpha(d)/d$ in terms of spectral radii of bounded-degree graphs. In low dimensions ($d \le 23$), by contrast, $N(d)$ is governed by arithmetic, $2$-adic, and algebraic-number-theoretic obstructions on the characteristic polynomial of the associated \emph{Seidel matrix} $S = \alpha^{-1}(G - I_n) \in \mathbb{Z}^{n \times n}$, where $G = (\langle v_i, v_j \rangle)_{i,j=1}^n$ is the Gram matrix.

Prior to the present work (and its companion short communication~\cite{CompanionMRL}), exact values of $N(d)$ were known for all $d \le 17$ and $21 \le d \le 23$, leaving $d \in \{18, 19, 20\}$ as the lowest open dimensions:
\[
  N(18) \in \{57, 58, 59\}, \qquad N(19) \in \{72, 73, 74\}, \qquad N(20) \in \{90, 91, 92, 93, 94\}.
\]
We establish the following four main theorems.

\begin{theorem}[Unconditional upper bound in $\mathbb{R}^{18}$]\label{thm:main18}
There does not exist an equiangular line system of $59$ lines in $\mathbb{R}^{18}$. Consequently, $N(18) \le 58$, and since $N(18) \ge 57$~\cite{GSY2023}, we have $N(18) \in \{57, 58\}$.
\end{theorem}

\begin{theorem}[Bidirectional $\SRG(75, 32, 10, 16)$ equivalence and upper bound in $\mathbb{R}^{19}$]\label{thm:main19}
An equiangular line system of $74$ lines in $\mathbb{R}^{19}$ exists if and only if there exists a strongly regular graph $\SRG(75, 32, 10, 16)$. More precisely:
\begin{enumerate}
  \item Every hypothetical $74 \times 74$ Seidel matrix $S$ with $5I_{74} + S \succeq 0$ and $\rk(5I_{74} + S) \le 19$ must have characteristic polynomial
  \[
    \chi_S(x) = (x + 5)^{55}(x - 15)^{17}(x - 11)(x - 9)
  \]
  and is switching-equivalent to a $74 \times 74$ vertex-deleted principal submatrix of the Seidel matrix $J_{75} - I_{75} - 2A_{75}$ of an $\SRG(75, 32, 10, 16)$. Conversely, deleting any vertex from an $\SRG(75, 32, 10, 16)$ yields $74$ equiangular lines in $\mathbb{R}^{19}$ with this spectrum.
  \item Combining part~(1) with the theorem of Azarija and Marc~\cite{AzarijaMarc2018} that no $\SRG(75, 32, 10, 16)$ exists, we obtain $N(19) \le 73$, and hence $N(19) \in \{72, 73\}$.
\end{enumerate}
\end{theorem}

\begin{remark}[Explicit external dependency and open boundary pairs]\label{rem:external_open}
Theorem~\ref{thm:main18} depends solely on first-principles proofs and finite exact integer and rational certificates archived with this paper. In Theorem~\ref{thm:main19}, part~(1) is proved unconditionally here (and formalized bidirectionally in Lean~4), whereas part~(2) explicitly relies on Azarija and Marc~\cite{AzarijaMarc2018} (\emph{Linear Algebra Appl.} \textbf{557} (2018), 62--83), who ruled out $\SRG(75, 32, 10, 16)$ via star-complement clique enumeration ($\approx 1.5 \times 10^5$ CPU hours); we do not re-execute that graph census. After Theorems~\ref{thm:main18} and~\ref{thm:main19}, the single-step boundary pairs $(18, 58)$ and $(19, 73)$ remain open (see~\S\ref{sec:compute} for exact polynomial prefix benchmarks and bounded constructive search evidence).
\end{remark}

\begin{theorem}[Odd-order mod-$2^{n+1}$ Pfaffian congruence]\label{thm:pfaffian}
Let $S \in \mathbb{Z}^{n \times n}$ be a symmetric Seidel matrix of odd order $n \ge 1$. Then $2^{n-1} \mid \chi_S(2z - 1)$ in $\mathbb{Z}[z]$, and there exists a monic polynomial $Q_0(z) \in \mathbb{Z}[z]$ of degree $(n - 1)/2$ such that
\begin{equation}\label{eq:intro_pfaffian}
  \frac{1}{2^{n-1}}\,\chi_S(2z - 1) \equiv (2z - n)\,Q_0(z)^2 \pmod{4}.
\end{equation}
\end{theorem}

\begin{theorem}[Generalized Galois orthogonality in degree-$e$ number fields]\label{thm:galois_intro}
Let $S \in \mathbb{Z}^{n \times n}$ be a symmetric Seidel matrix and let $q(x) \in \mathbb{Z}[x]$ be a monic irreducible factor of $\chi_S(x)$ of multiplicity $1$ and degree $e \ge 2$, with $K := \mathbb{Q}[x]/(q(x))$ and $\bar{\lambda} := x \bmod q$. Write $\chi_S(x) = (x+5)^m f(x)$ and $\chi_{S_u}(x) = (x+5)^{\max(0, m-1)} f_u(x)$ ($m \ge 0$). Suppose $u_0 \in \{1, \dots, n\}$ satisfies $f_{u_0}(\bar{\lambda}) \neq 0$ in $K$, and for each $u \in \{1, \dots, n\}$ let $r_u \in K$ satisfy $r_u^2 = f_u(\bar{\lambda})/f_{u_0}(\bar{\lambda})$ ($r_u = 0$ if $q \mid f_u$). Then for all $0 \le a \le b \le e - 1$,
\begin{equation}\label{eq:intro_galois}
  \sum_{u=1}^n \Big(\Tr_{K/\mathbb{Q}}(\bar{\lambda}^a r_u)\,\Tr_{K/\mathbb{Q}}(\bar{\lambda}^b r_u) - \Tr_{K/\mathbb{Q}}(\bar{\lambda}^{a+b} r_u^2)\Big) = 0.
\end{equation}
\end{theorem}

\subsection{Primary-source literature review, conventions, and errata}\label{subsec:lit_review}
Table~\ref{tab:frontier_dims} summarizes the known bounds on $N(d)$ for $14 \le d \le 23$, incorporating Theorems~\ref{thm:main18} and~\ref{thm:main19}.

\begin{table}[htbp]
\centering
\caption{Maximum cardinality $N(d)$ of equiangular lines in $\mathbb{R}^d$ for $14 \le d \le 23$.}\label{tab:frontier_dims}
\footnotesize
\setlength{\tabcolsep}{4.5pt}
\begin{tabular}{@{}cccll@{}}
\toprule
$d$ & $N(d)$ & Optimal $\alpha$ & Lower bound origin & Upper bound origin \\
\midrule
$14$ & $28$ & $1/3, 1/5$ & Lemmens--Seidel~\cite{LemmensSeidel1973} & Greaves--Syatriadi--Yatsyna~\cite{GSY2021} \\
$15$ & $36$ & $1/5$ & Lemmens--Seidel~\cite{LemmensSeidel1973} & Relative bound~\cite{vanLintSeidel1966,LemmensSeidel1973} \\
$16$ & $40$ & $1/5$ & Witt design / Lemmens--Seidel~\cite{LemmensSeidel1973} & Greaves--Syatriadi--Yatsyna~\cite{GSY2021} \\
$17$ & $48$ & $1/5$ & Netto / Lemmens--Seidel~\cite{LemmensSeidel1973,GKMS2016} & Greaves--Syatriadi--Yatsyna~\cite{GSY2023} \\
$\mathbf{18}$ & $\mathbf{57\text{--}58}$ & $1/5$ & GSY23~\cite{GSY2023}; LMTY25~\cite{LMTY2025}; MSY26~\cite{MSY2026} & \textbf{Theorem~\ref{thm:main18}} (prior $\le 59$: GS24~\cite{GS2024}) \\
$\mathbf{19}$ & $\mathbf{72\text{--}73}$ & $1/5$ & Asche / Taylor ($S(5,8,24)$)~\cite{Taylor1977,LinYu2020} & \textbf{Theorem~\ref{thm:main19}} +~\cite{AzarijaMarc2018} (prior $\le 74$:~\cite{GSY2021}) \\
$20$ & $90\text{--}94$ & $1/5$ & Asche / Taylor ($S(5,8,24)$)~\cite{Taylor1977,LinYu2020} & Greaves--Syatriadi--Yatsyna~\cite{GSY2021} \\
$21$ & $126$ & $1/5$ & Taylor / Lemmens--Seidel~\cite{LemmensSeidel1973,Taylor1977} & Relative bound~\cite{LemmensSeidel1973} \\
$22$ & $176$ & $1/5$ & Higman--Sims / Lemmens--Seidel~\cite{LemmensSeidel1973} & Relative bound~\cite{LemmensSeidel1973} \\
$23$ & $276$ & $1/5$ & Conway / Leech lattice~\cite{LemmensSeidel1973} & Gerzon bound $23 \cdot 24 / 2 = 276$~\cite{LemmensSeidel1973} \\
\bottomrule
\end{tabular}
\end{table}

We record the progression of upper-bound and constructive results in dimensions $17$--$20$, together with two literature conventions and errata verified directly from primary sources:

\begin{enumerate}
  \item \emph{Spectral-interlacing upper bounds:}
  Greaves, Koolen, Munemasa, and Sz\"oll\H{o}si~\cite{GKMS2016} introduced Eulerian determinant congruences modulo~$4$ and~$8$ together with almost-equality AM--GM bounds, proving $N(14) \le 29$ and $N(16) \le 41$. Yu~\cite{Yu2015} established $N(19) \le 75$ and $N(20) \le 95$, and Greaves~\cite{Greaves2018} proved $N(18) \le 60$ while reducing $3$-eigenvalue even-order Seidel matrices with $|n - 1 + \lambda\mu| < 8$ to regular graphs. Greaves and Yatsyna~\cite{GY2019} established Type-2 and weakly Type-2 $2$-adic shift divisibilities, proved the Burnside walk-orbit bound $|\mathcal{P}_{n,e}| \le 2^{\binom{e-2}{2} + (n \bmod 2)}$ on characteristic polynomials modulo $2^e$, and combined McKee--Smyth real-rooted polynomial enumeration~\cite{McKeeSmyth2004} with one-vertex angle compatibility to show $N(17) \le 49$.

  In~\cite{GSY2021}, Greaves, Syatriadi, and Yatsyna proved the general shifted multiplicity bound at the odd integer $\kappa$ closest to $5(n-d)/d$, establishing $N(14)=28$, $N(16)=40$, $N(19) \le 74$ (ruling out $(19, 75)$ across $8$ admissible characteristic polynomials: $6$ Farkas-infeasible, $1$ warranted-infeasible, and $1$ equivalent to $\SRG(75, 32, 10, 16)$~\cite{AzarijaMarc2018}), and $N(20) \le 94$ (via $\SRG(95, 40, 12, 20)$~\cite{AzarijaMarc2020}). In~\cite{GSY2023}, the same authors proved $N(17) = 48$ by ruling out all $194$ weakly Type-2 candidate polynomials for $(17, 49)$ (using Bussemaker--Haemers--Mathon--Wilbrink~\cite{BHMW1989} on $\SRG(49, 16, 3, 6)$, Farkas certificates, Seidel compatibility, and $8$ two-vertex and Eulerian arguments). In~\cite{GS2024}, Greaves and Syatriadi ruled out $(18, 60)$ ($N(18) \le 59$) across $44$ candidate polynomials: $39$ by one-vertex Farkas duality, $4$ by submatrix extraction and warranted compatibility, and the $44$th candidate $(x+5)^{42}(x-11)^{15}(x-15)^3$ by regular-graph equivalence~\cite{Greaves2018} and De Caen--Jacobi complementary minor analysis~\cite{DeCaen2000} forcing $3$ blocks of $5C_4$ and a clique bound $\omega(G_u) = 7 < 9$. (We note in passing that sequence A002853 in the OEIS still lists $57 \le N(18) \le 60$ despite citing~\cite{GS2024}.)

  \item \emph{Structural obstacles beyond $N(18) \le 59$ and $N(19) \le 74$:}
  Four algebraic obstacles separated $(18, 60)$ from $(18, 59)$: (a) forced multiplicity at $\kappa = 11$ drops from $6$ to $1$, raising the cofactor degree from $12$ to $17$ and $\mathcal{T}$ from $48$ to $65$; (b) order parity flips from even ($60$, Strict Type~2: $2^k \mid b_k$) to odd ($59$, Weakly Type~2: $2^{k-1} \mid b_k$ on odd $k$), prompting~\cite[\S 9]{GSY2023} to remark that $(18, 59)$ appeared out of reach on a personal computer; (c) Greaves and Syatriadi's Fisk-cone linear programming method~\cite{GS2025} accelerates vertex-deletion deck generation by $85$--$90\times$ once $\chi_S$ is fixed, but does not apply to characteristic polynomial enumeration where no fixed interlacing root grid exists; and (d) none of $(18, 59)$, $(18, 58)$, $(19, 74)$, or $(19, 73)$ admits an integer $3$-eigenvalue Seidel spectrum, while hundreds of interlacing-compatible candidates possess non-rational simple factors of degree $e \in \{2, 3, 4, 5\}$ where rational angle compatibility breaks down. By contrast, $(19, 74)$ has even order $74$ (Strict Type~2) and $m_{15} \ge 6$ ($\deg\phi = 13, \mathcal{T} = 52$), making polynomial enumeration tractable ($169$ candidates) while concentrating the mathematical difficulty in the structural analysis of $p_{13}$ and $p_{136}$.

  \item \emph{Constructive classifications in $(18, 57)$ and $(19, 72)$, the LMTY25 sign convention, and the GSY23 $F_4$ erratum:}
  In $\mathbb{R}^{19}$, $72$ equiangular lines arise from Asche's selection of $72$ octads of the Witt design $S(5, 8, 24)$~\cite{Taylor1977,LinYu2020}, with $\chi_S(x) = (x+5)^{53}(x-13)^{16}(x-19)^3$ (containing saturated subsets of size $54$~\cite{Szollosi2019} and $56$~\cite{LinYu2020} in $\mathbb{R}^{18}$). In $\mathbb{R}^{18}$, Greaves, Syatriadi, and Yatsyna~\cite[\S 2]{GSY2023} discovered four $18 \times 57$ integer matrices $F_1, F_2, F_3, F_4$ realizing $57$ lines (with $S_3, S_4$ not embeddable in Asche's $72$ lines; see also Martinet~\cite{Martinet2024}). Using switching roots~\cite{CKMY2021,CKLY2022}, Lin, Munemasa, Taniguchi, and Yoshino~\cite{LMTY2025} constructed $\ge 246{,}896$ non-isometric strongly maximal $57$-line sets in $\mathbb{R}^{18}$ inside an integral overlattice of $A_9 \oplus A_9 \oplus A_1$ via $1$-factorizations of $K_{10}$, recovering $F_1$--$F_4$. Subsequently, Munemasa, Sz\"oll\H{o}si, and Yoshino~\cite{MSY2026} classified maximum affine equiangular sets of size $57$ in an overlattice of $A_5^3 \oplus A_1^4$ via order-$6$ Latin squares and Pasch configurations, obtaining the first \emph{non-strongly-maximal} $(18, 57)$ sets (which lift to at most $57 + 12 = 69 < 72$ lines in $\mathbb{R}^{19}$) and the first $5$-eigenvalue $(18, 57)$ spectrum $(x+5)^{39}(x-13)^8(x-9)^8(x^2-19x+72)$.

  We document two essential primary-source reconciliation notes:
  \begin{itemize}
    \item \emph{Negated Seidel convention in~\cite{LMTY2025}:} Whereas~\cite{GKMS2016,GSY2021,GSY2023,GS2024,MSY2026} and the present paper define $S_{\mathrm{std}} := \alpha^{-1}(G - I_n)$ (putting the degenerate eigenvalue at $-5$ so $5I_n + S_{\mathrm{std}} \succeq 0$), Lin et al.~\cite[\S 2]{LMTY2025} define $S_{\mathrm{LMTY}} := \alpha^{-1}(I_n - G) = -S_{\mathrm{std}}$, negating all eigenvalues ($x \mapsto -x$). Under $S_{\mathrm{std}}$, their two cyclic $6$-eigenvalue spectra~\cite[Ex.~9.2]{LMTY2025} convert from $(x-5)^{39}(x+9)^2\allowbreak(x+11)^{11}\allowbreak(x+13)^3\allowbreak(x^2+17x+36)$ and $(x-5)^{39}(x+9)\allowbreak(x+11)^{13}\allowbreak(x+13)^2\allowbreak(x^2+17x+32)$ to
    \begin{gather*}
      (x+5)^{39}(x-9)^2(x-11)^{11}(x-13)^3(x^2-17x+36), \\
      (x+5)^{39}(x-9)(x-11)^{13}(x-13)^2(x^2-17x+32).
    \end{gather*}
    \item \emph{Printed characteristic polynomial erratum on $F_4$ in~\cite{GSY2023}:} In~\cite[p.~1876]{GSY2023} (arXiv:\allowbreak 2104.04330v2), the characteristic polynomial of $S_4 = \frac{1}{2}F_4^T F_4 - 5I_{57}$ (from Fig.~4) is printed as
    \[
      (x+5)^{39}(x-11)^{13}(x^2-18x+57)(x^3-34x^2+361x-1136).
    \]
    Direct exact integer evaluation of $F_4$ from~\cite[Fig.~4]{GSY2023} and independent reconstruction from $\mathcal{I}_4$ in~\cite[Appendix]{LMTY2025} both yield the true characteristic polynomial
    \[
      \chi_{S_4}(x) = (x+5)^{39}(x-11)^{13}(x^2 - 20x + 87)(x^3 - 32x^2 + 303x - 716)
    \]
    (both expressions share $\tr(S_4) = 0$, $\tr(S_4^2) = 3192$, $7$ distinct eigenvalues, and non-divisibility by $x - 13$, so all mathematical conclusions of~\cite{GSY2023} remain unaffected).
  \end{itemize}

  \item \emph{Inapplicability of $k$-point SDP hierarchies below $d = 23$:}
  Semidefinite programming hierarchies---from Barg and Yu's $3$-point Bachoc--Vallentin bound ($\Delta_3(G)^*$, improving Gerzon's bound on $N(d)$ for $24 \le d \le 136$)~\cite{BargYu2014} and Yu's SDP relative bound ($23 \le d \le 59$)~\cite{Yu2017} to $k$-point hierarchies ($k \le 6$, establishing $N_{1/5}(d) \le 276$ for $23 \le d \le 70$ and sharpening $\Delta_3(G)^*$ for $65 \le d \le 230$) of de~Laat, Machado, de~Oliveira~Filho, and Vallentin~\cite{deLaat2022} and symbolic $4$-point bounds of Kao and Yu~\cite{KaoYu2022}---yield, as noted explicitly in~\cite[\S 5]{deLaat2022}, no improvement below $d = 23$ over the integer relative bound $\lfloor 24d/(25 - d) \rfloor \in \{61, 76, 96\}$ in dimensions $d = 18, 19, 20$.
\end{enumerate}

\subsection{Classification summary and spectral eliminations}
Throughout this paper, $p_0,\allowbreak \dots,\allowbreak p_{168}$ denote the $169$ admissible characteristic polynomials $(x+5)^{55}\allowbreak Q(x)$ for $(19, 74)$ (listed in Appendix~\ref{app:census_19_74}), and $q_0,\allowbreak \dots,\allowbreak q_{1213}$ denote the $1{,}214$ admissible characteristic polynomials $(x+5)^{41}\allowbreak Q(x)$ for $(18, 59)$ (summarized in Appendix~\ref{app:census_18_59}), ordered lexicographically as in the archived tables (\S\ref{sec:compute}). Table~\ref{tab:master_funnel} summarizes the spectral eliminations for $(19, 74)$ and $(18, 59)$ alongside the three calibration benchmarks $(19, 75)$, $(18, 60)$, and $(17, 49)$. Every elimination count is backed by exact integer and rational certificates.

\begin{table}[htbp]
\centering
\caption{Summary of spectral eliminations across $(19, 74)$, $(18, 59)$, and literature calibration benchmarks.}\label{tab:master_funnel}
\resizebox{\textwidth}{!}{%
\begin{tabular}{@{}lcrrrrrrl@{}}
\toprule
Case $(d,n)$ & $\kappa$ ($m_\kappa$) & Admissible & LP (Farkas) & Pfaffian / $K^\times$ / Extr. & Rem. & Eulerian / Gal. & Higher-order & Final outcome \\
\midrule
$(19, 74)$ & $15$ ($\ge 6$) & $169$ & $-147$ & $-15$ ($8+6+1$) & $7$ & $-2$ ($p_{74},p_{134}$) & $-4$ ($p_6,p_{24},p_{47},p_{13}$) & $1$: $p_{136} \Leftrightarrow \SRG(75,32,10,16)$ \\
$(18, 59)$ & $11$ ($\ge 1$) & $1{,}214$ & $-572$ & $-556$ ($0+552+4$) & $86$ & $-66$ & $-20$ ($q_{861} + 19$ $2$-vert.) & $0$ ($N(18) \le 58$ unconditional) \\
\midrule
$(19, 75)$ & $15$ ($\ge 12$) & $8^{\text{a}}$ & $-6$ & $-1$ (warranted) & $1$ & --- & --- & $\SRG(75,32,10,16)$ (matches~\cite{GSY2021}) \\
$(18, 60)$ & $11$ ($\ge 6$) & $44$ & $-39$ & $-3$ ($1+1+1$) & $2$ ($11,34$) & $-1$ (cand.~$11$) & $-1$ (cand.~$34$) & $0$ (matches~\cite{GS2024}) \\
$(17, 49)$ & $9$ ($\ge 4$) & $109^{\text{b}}$ & $-84$ & $-16$ ($0+16+0$) & $9$ & $-5$ & $-4$ & $0$ (matches~\cite{GSY2023,BHMW1989}) \\
\bottomrule
\end{tabular}%
}
\par\vspace{3pt}
\raggedright\scriptsize
$^{\text{a}}$Unfiltered real-rooted polynomial enumeration produces $14$ polynomials, reduced to $8$ by $\mathcal{P}_{75,7}$ and Theorem~\ref{thm:pfaffian}, matching~\cite[Prop.~3.3]{GSY2021} (see~\S\ref{sec:compute} for exact file records).\\
$^{\text{b}}$Filtration basis on $(17, 49)$: weakly Type-2 enumeration gives $2{,}291$ unfiltered polynomials and $194$ after $\mathcal{P}_{49,7}$, matching~\cite[Lem.~3.6]{GSY2023}. Enforcing exact full-degree odd-$n$ parity ($2^{2r}\mid e_{2r}$, Theorem~\ref{thm:twoadic_val}(3)) gives $704$ unfiltered $\to 188$ after $\mathcal{P}_{49,7}$ ($6$ of the $194$ fail $2^8\mid e_8$) $\to 109$ after Theorem~\ref{thm:pfaffian}.
\end{table}

% ============================================================================
\section{Angle determination and spectral factorization}\label{sec:phase1}
% ============================================================================

\subsection{Equiangular systems and Seidel--Gram equivalence}
Let $d, n \in \mathbb{N}$ and $\alpha \in \mathbb{R}$. A tuple $v = (v_1, \dots, v_n) \in (\mathbb{R}^d)^n$ is an \emph{equiangular system} of parameters $(d, n, \alpha)$ if $0 \le \alpha < 1$, $\|v_i\|_2 = 1$ for all $i$, and $|\langle v_i, v_j \rangle| = \alpha$ for all $i \neq j$. By Cauchy--Schwarz, $|\langle v_i, v_j \rangle| \le 1$ with equality if and only if $v_i = \pm v_j$, so for $n \ge 2$ the strict inequality $\alpha < 1$ is equivalent to pairwise distinctness of the lines $\mathbb{R}v_1, \dots, \mathbb{R}v_n$. Moreover, if $\alpha = 0$, the vectors $v_1, \dots, v_n$ are orthonormal and hence linearly independent in $\mathbb{R}^d$, forcing $n \le d$; contrapositively, $n > d \implies \alpha > 0$.

A matrix $S \in \mathbb{Z}^{n \times n}$ is a \emph{Seidel matrix} if $S^T = S$, $S_{ii} = 0$ for all $i$, and $S_{ij} \in \{+1, -1\}$ for all $i \neq j$. For any $D = \diag(\delta_1, \dots, \delta_n)$ with $\delta_i \in \{\pm 1\}$, the \emph{switched} matrix $S' = D S D$ is again a Seidel matrix with $\chi_{S'}(x) = \chi_S(x)$ and $\rk(\alpha^{-1}I_n + S') = \rk(\alpha^{-1}I_n + S)$, corresponding geometrically to $v_i \mapsto \delta_i v_i$.

\begin{theorem}[Exact Seidel--Gram equivalence]\label{thm:seidel_gram}
Let $d, n \in \mathbb{N}$ and $\alpha \in (0, 1)$. An equiangular system $v \in (\mathbb{R}^d)^n$ with common inner product $\alpha$ exists if and only if there exists a symmetric $n \times n$ Seidel matrix $S \in \mathbb{Z}^{n \times n}$ such that
\[
  M := \frac{1}{\alpha} I_n + S \succeq 0 \qquad \text{and} \qquad \rk_{\mathbb{R}}(M) \le d.
\]
Furthermore, every $n \times n$ Seidel matrix satisfies $\tr(S) = 0$ and $\tr(S^2) = n(n - 1)$.
\end{theorem}

\begin{proof}
($\Rightarrow$) Let $V \in \mathbb{R}^{d \times n}$ have columns $v_1, \dots, v_n$, and set $G := V^T V$. Since $\alpha > 0$, setting $S_{ii} := 0$ and $S_{ij} := \alpha^{-1}\langle v_i, v_j \rangle \in \{\pm 1\}$ ($i \neq j$) defines a symmetric Seidel matrix with $M = \alpha^{-1}I_n + S = \alpha^{-1} V^T V$. For any $x \in \mathbb{R}^n$, $x^T M x = \alpha^{-1}\|Vx\|_2^2 \ge 0$, and $\rk(M) = \rk(V^T V) = \rk(V) \le d$.

($\Leftarrow$) Conversely, let $G := \alpha M = I_n + \alpha S \succeq 0$ have rank $r := \rk(G) = \rk(M) \le d$. By the real spectral theorem, $G = \sum_{k=1}^n \mu_k q_k q_k^T$ for an orthonormal eigenbasis $q_1, \dots, q_n \in \mathbb{R}^n$ with eigenvalues $\mu_1 \ge \dots \ge \mu_r > 0 = \mu_{r+1} = \dots = \mu_n$. Define $V \in \mathbb{R}^{d \times n}$ by setting row $k$ equal to $\sqrt{\mu_k}\,q_k^T$ for $1 \le k \le r$ and $0$ for $r < k \le d$. Its columns $v_1, \dots, v_n \in \mathbb{R}^d$ satisfy $\langle v_i, v_j \rangle = (V^T V)_{ij} = G_{ij} = \delta_{ij} + \alpha S_{ij}$, whence $\|v_i\|_2 = 1$ and $|\langle v_i, v_j \rangle| = \alpha$ ($i \neq j$). Finally, $\tr(S) = \sum_i S_{ii} = 0$ and $(S^2)_{ii} = \sum_{j \neq i} (\pm 1)^2 = n - 1$, giving $\tr(S^2) = n(n - 1)$.
\end{proof}

\subsection{Classical bounds, Neumann's theorem, and the Lemmens--Seidel theorem}
For completeness, we give concise, self-contained proofs of the four preliminary bounds determining $\alpha$.

\begin{proposition}[Gerzon and relative bounds~\cite{vanLintSeidel1966,LemmensSeidel1973}]\label{prop:gerzon_rel}
Let $v_1, \dots, v_n \in \mathbb{R}^d$ ($d \ge 1$) be an equiangular system with parameter $\alpha \in [0, 1)$.
\begin{enumerate}
  \item \emph{(Gerzon's bound)} $n \le d(d + 1)/2$.
  \item \emph{(Relative bound)} $n^2 \le d\big(n + n(n - 1)\alpha^2\big)$; equivalently, $(n - d)\alpha^{-2} \le d(n - 1)$ whenever $n > d$, and $n \le \frac{d(1 - \alpha^2)}{1 - d\alpha^2}$ whenever $d\alpha^2 < 1$.
\end{enumerate}
\end{proposition}

\begin{proof}
Set $P_i := v_i v_i^T \in \Sym_d(\mathbb{R})$, equipped with $\langle A, B \rangle_F = \tr(AB)$, so $\langle P_i, P_j \rangle_F = \langle v_i, v_j \rangle^2$ equals $1$ for $i = j$ and $\alpha^2$ for $i \neq j$.
(1) If $\sum_{i=1}^n c_i P_i = 0$, then $0 = \|\sum_i c_i P_i\|_F^2 = (1 - \alpha^2)\sum_i c_i^2 + \alpha^2(\sum_i c_i)^2 \ge (1 - \alpha^2)\sum_i c_i^2$; since $\alpha < 1$, all $c_i = 0$, so $n \le \dim\Sym_d(\mathbb{R}) = d(d+1)/2$.
(2) Setting $W := \sum_{i=1}^n P_i$, we have $\tr(W) = n$ and $\tr(W^2) = n + n(n-1)\alpha^2$, so $0 \le \|W - \frac{n}{d}I_d\|_F^2 = \tr(W^2) - \frac{n^2}{d}$, which rearranges to both stated inequalities.
\end{proof}

\begin{theorem}[Neumann~\cite{LemmensSeidel1973}]\label{thm:neumann}
If $d \ge 1$, $n > 2d$, and an equiangular system of $n$ lines in $\mathbb{R}^d$ has parameter $\alpha$, then $\alpha = 1/k$ for an odd integer $k \ge 3$.
\end{theorem}

\begin{proof}
Since $n > 2d \ge d$, we have $\alpha > 0$ and $n \ge 3$. Let $S$ be the Seidel matrix from Theorem~\ref{thm:seidel_gram} with $\gamma := \alpha^{-1} > 1$. Then $-\gamma$ is a root of $\chi_S(x) \in \mathbb{Z}[x]$ of multiplicity $m = \dim\ker(\gamma I_n + S) = n - \rk(\gamma I_n + S) \ge n - d$, satisfying $2m \ge 2(n - d) > n$ and $m \ge d + 1 \ge 2$. Let $p(x) \in \mathbb{Z}[x]$ be the monic minimal polynomial of $-\gamma$ over $\mathbb{Q}$ (integral by Gauss's lemma). Since $\mathbb{Q}$ has characteristic $0$, $p$ is separable and $\Gal(\overline{\mathbb{Q}}/\mathbb{Q})$ preserves root multiplicities of $\chi_S$, so $p(x)^m \mid \chi_S(x)$ and $n = \deg\chi_S \ge m \deg p$. As $2m > n$, $\deg p = 1$, so $k := \gamma \in \mathbb{Z}_{\ge 2}$. If $k$ were even, then $S + kI_n \equiv J_n - I_n \pmod 2$. Over $\mathbb{F}_2$, $(J_n - I_n)x = 0 \iff x_i = \sum_j x_j$ for all $i$, which forces $x = 0$ when $n$ is even ($\rk_{\mathbb{F}_2} = n$) and $\dim_{\mathbb{F}_2}\ker(J_n - I_n) = 1$ when $n$ is odd ($\rk_{\mathbb{F}_2} = n - 1$). Thus $\rk_{\mathbb{R}}(S + kI_n) \ge \rk_{\mathbb{F}_2}(J_n - I_n) \ge n - 1$, giving $m \le 1$, contradicting $m \ge 2$. Hence $k \ge 3$ is odd.
\end{proof}

\begin{theorem}[Lemmens--Seidel~\cite{LemmensSeidel1973}]\label{thm:lemmens_seidel}
For every $d \ge 15$, $N_{1/3}(d) = 2(d - 1)$. In particular, $N_{1/3}(18) = 34$ and $N_{1/3}(19) = 36$.
\end{theorem}

\begin{proof}
The vectors $\frac{1}{\sqrt{3}}(e_d \pm \sqrt{2} e_i)$ ($1 \le i \le d - 1$) realize $2(d - 1)$ lines at $\alpha = 1/3$. Conversely, let $u_i := \sqrt{3}\,v_i \in \mathbb{R}^d$, so $\|u_i\|_2^2 = 3$ and $\langle u_i, u_j \rangle \in \{\pm 1\}$ ($i \neq j$).

\emph{Case 1 (A zero-sum $4$-tuple exists).} After reindexing and sign flips, $u_1 + u_2 + u_3 + u_4 = 0$. Taking inner products with $u_r$ ($1 \le r \le 4$) forces $\langle u_r, u_s \rangle = -1$ ($r \neq s$), so $G[\{1,2,3\}] = 4I_3 - J_3$ ($\det = 16$) and $W := \spn(u_1, \dots, u_4)$ has dimension $3$, with orthonormal basis $p_1 = \frac{1}{4}(u_1+u_2-u_3-u_4)$, $p_2 = \frac{1}{4}(u_1-u_2+u_3-u_4)$, $p_3 = \frac{1}{4}(u_1-u_2-u_3+u_4)$. Every $k \ge 5$ satisfies $\sum_{r=1}^4 \langle u_k, u_r \rangle = 0$ with terms in $\{\pm 1\}$; switching $u_k$ so $\langle u_k, u_1 \rangle = +1$ assigns $k$ to one of three pillars $C_1, C_2, C_3$ according to the second $+1$, giving $u_k = p_a + r_k$ ($k \in C_a$) with $r_k \in W^\perp$ and $\|r_k\|_2^2 = 2$. Within $C_a$, $\langle r_k, r_\ell \rangle \in \{0, -2\}$ ($k \neq \ell$), and $\langle r_k, r_\ell \rangle = -2 \implies \|r_k + r_\ell\|_2^2 = 0 \implies r_\ell = -r_k$; selecting one vector per antipodal pair gives an orthogonal set $B_a \subseteq W^\perp$ with $|C_a| \le 2|B_a| \le 2(d - 3)$. If a second pillar $C_b$ ($b \neq a$) is nonempty, pick $j_0 \in C_b$; then $\langle r_{j_0}, r_k/\sqrt{2} \rangle^2 = \frac{1}{2}$ for all $k \in B_a$, so Bessel's inequality gives $2 = \|r_{j_0}\|_2^2 \ge \frac{1}{2}|B_a| \implies |C_a| \le 8$. Hence $n \le \max(4 + 24, 4 + 2(d - 3)) = \max(28, 2(d - 1)) = 2(d - 1)$ for $d \ge 15$.

\emph{Case 2 (No zero-sum $4$-tuple exists).} Assume $n \ge 2$ and switch so $\langle u_1, u_j \rangle = +1$ for $j \ge 2$. Projecting $w_j := \sqrt{\frac{3}{2}}(u_j - \frac{1}{3}u_1) \in u_1^\perp$ ($2 \le j \le n$) yields Gram matrix $M = 3I_{n-1} + J_{n-1} - 3A_\Gamma \succeq 0$ of rank $\le d - 1$, where $j \sim_\Gamma \ell \iff \langle u_j, u_\ell \rangle = -1$. If $\{a,b,c\}$ formed $K_3$ in $\Gamma$, $\|u_a + u_b + u_c - u_1\|_2^2 = 0$; if $\{a,b\}$ and $\{c,e\}$ formed an induced $2K_2$, $\|u_a + u_b - u_c - u_e\|_2^2 = 0$; both contradict the hypothesis, so $\Gamma$ is $(K_3, 2K_2)$-free.
\begin{itemize}
  \item If $\Gamma$ has an odd cycle, a shortest odd cycle has length $5$ (length $3$ is $K_3$; length $\ge 7$ contains $2K_2$), and $(K_3, 2K_2)$-freeness forces every non-isolated vertex outside $C_5$ to have neighborhood $\{i-1, i+1\}$ on $C_5$, making $\Gamma$ a complete $C_5$-blowup $V_1 \sqcup \dots \sqcup V_5$ plus isolates $Z_0$. On the subspace $W_0$ of vectors summing to $0$ on each of $V_1, \dots, V_5, Z_0$, $A_\Gamma x = J x = 0 \implies Mx = 3x$, so $d - 1 \ge \rk(M) \ge \dim W_0 \ge (n - 1) - 6$, giving $n \le d + 6$.
  \item Otherwise $\Gamma$ is bipartite on $X \sqcup Y$ (plus isolates $Z_0$). By $2K_2$-freeness, neighborhoods of $X$ are nested, partitioning non-isolates into a chain graph $X_1 \sqcup \dots \sqcup X_k$ and $Y_1 \sqcup \dots \sqcup Y_k$ ($X_i \sim Y_j \iff i + j \le k + 1$). If $k \ge 5$, picking $x_i \in X_i, y_i \in Y_i$ ($1 \le i \le 5$) and setting $z = (-2, 0, 1, 1, 2)$ on $(x_1, \dots, x_5)$, $(2, 0, -1, -1, -2)$ on $(y_1, \dots, y_5)$, and $0$ elsewhere gives $\mathbf{1}^T z = 0$, $\|z\|_2^2 = 20$, and $z^T A_\Gamma z = 22$, whence $z^T M z = 3(20 - 22) = -6 < 0$, contradicting $M \succeq 0$. Thus $k \le 4$, and the zero-sum subspace across $X_1, \dots, X_k, Y_1, \dots, Y_k, Z_0$ gives $d - 1 \ge \rk(M) \ge (n - 1) - 9$, i.e.\ $n \le d + 9 \le 2(d - 1)$ for $d \ge 15$.\qedhere
\end{itemize}
\end{proof}

\begin{corollary}[Uniqueness of $\alpha = 1/5$ for $(18, n)$ and $(19, n)$]\label{cor:one_fifth}
For $(d, n) \in \{(18, 58),\allowbreak (18, 59),\allowbreak (19, 73),\allowbreak (19, 74)\}$, every equiangular system has $\alpha = 1/5$, yielding a symmetric $n \times n$ Seidel matrix $S$ with $5I_n + S \succeq 0$ and $\rk(5I_n + S) \le d$.
\end{corollary}

\begin{proof}
Here $n > 2d$, so $\alpha = 1/k$ ($k \ge 3$ odd) by Theorem~\ref{thm:neumann}. Theorem~\ref{thm:lemmens_seidel} rules out $k = 3$ ($34 < 58$, $36 < 73$), while for $k \ge 7$ Proposition~\ref{prop:gerzon_rel}(2) requires $49(n - d) \le d(n - 1)$, failing for all four pairs ($1960 > 1026$, $2009 > 1044$, $2646 > 1368$, $2695 > 1387$).
\end{proof}

\subsection{Spectral factorization and multiplicity bounds (\texorpdfstring{$\rk \le d$}{rk <= d})}
\begin{proposition}[Spectral factorization and shifted moments]\label{prop:spectral_fac}
Let $S \in \mathbb{Z}^{n \times n}$ be a symmetric Seidel matrix with $G := 5I_n + S \succeq 0$ and $r := \rk(G) \le d \le n$.
\begin{enumerate}
  \item Since $\dim\ker(5I_n + S) = n - r \ge n - d$ and $(x + 5)^{n - d}$ is monic in $\mathbb{Z}[x]$, Euclidean division yields a unique monic cofactor $Q(x) \in \mathbb{Z}[x]$ of degree $d$ such that
  \begin{equation}\label{eq:cofactor_Q}
    \chi_S(x) = (x + 5)^{n - d} Q(x),
  \end{equation}
  whose roots $\lambda_1 \ge \dots \ge \lambda_d$ are real and satisfy $\lambda_i \ge -5$ (with $(x + 5)^{d - r} \mid Q(x)$, i.e.\ $Q(-5) = 0$, whenever $r < d$).
  \item Setting $h(y) := Q(y - 5) = \prod_{i=1}^d (y - \mu_i) = \sum_{k=0}^d (-1)^k e_k(\mu) y^{d - k} \in \mathbb{Z}[y]$ ($\mu_i = \lambda_i + 5 \ge 0$), we have $\sum_{i=1}^d \lambda_i = 5(n - d)$, $\sum_{i=1}^d \lambda_i^2 = n(n - 1) - 25(n - d)$, $e_1(\mu) = 5n$, and $e_2(\mu) = 12n(n - 1)$.
\end{enumerate}
\end{proposition}

\begin{proof}
Part~(1) follows from real symmetry of $S$ and Gauss's lemma. For part~(2), $\tr(G) = 5n$ and $\tr(G^2) = \tr(25I_n + 10S + S^2) = 25n + n(n - 1) = n(n + 24)$, so $e_2(\mu) = \frac{1}{2}((5n)^2 - n(n + 24)) = 12n(n - 1)$. (Notice $Q(y - 5)$ has nonnegative roots $\mu_i \ge 0$, whereas $(-1)^d Q(5 - y)$ has negative roots whenever $\lambda_1 > 5$.)
\end{proof}

\begin{proposition}[Shifted eigenvalue multiplicity bound~\cite{GSY2021}]\label{prop:gsy_mult}
Let $S$ and $Q$ be as in Proposition~\ref{prop:spectral_fac}, let $\gamma(n) := n \bmod 2$, and assume $2^{-(d - \gamma(n))} Q(2z - 1) \in \mathbb{Z}[z]$ (proved in Theorem~\ref{thm:twoadic_val} below). For any odd integer $\kappa \in 2\mathbb{Z} + 1$, define
\begin{gather*}
  \mathcal{T} := \sum_{i=1}^d (\lambda_i - \kappa)^2 = n(n - 1) - 25(n - d) - 10\kappa(n - d) + d\kappa^2, \\
  \theta := \min\!\big\{\eta \in \mathbb{Z}_{\ge 1} : \eta^\eta\,4^{\eta - \gamma(n)} > \mathcal{T}^\eta\big\}.
\end{gather*}
Then $(x - \kappa)^{\max(0,\, d + 1 - \theta)} \mid Q(x)$ in $\mathbb{Z}[x]$.
\end{proposition}

\begin{proof}
Since $\kappa$ is odd, $z = w + \frac{\kappa + 1}{2}$ shifts $2^{\gamma(n)} F(w) := 2^{-(d - \gamma(n))} Q(2w + \kappa) \in \mathbb{Z}[w]$ to a polynomial of leading coefficient $2^{\gamma(n)}$ with roots $w_i = (\lambda_i - \kappa)/2$. Factor $2^{\gamma(n)} F(w) = w^m \prod_{a=1}^s q_a(w)$ into primitive irreducibles in $\mathbb{Z}[w]$ with $q_a(0) \neq 0$, $\eta_a := \deg q_a \ge 1$, and $\eta := \sum_{a=1}^s \eta_a = d - m$. By Gauss's lemma, $\prod_{a=1}^s \mathrm{LC}(q_a) = 2^{\gamma(n)}$.
If $\gamma(n) = 0$, every $q_a$ is monic, so $\prod_{w_j \in Z(q_a)} w_j^2 = q_a(0)^2 \ge 1$, and AM--GM gives $4\sum_{w_j \in Z(q_a)} w_j^2 \ge 4\eta_a$, summing to $\mathcal{T} \ge 4\eta$.
If $\gamma(n) = 1$, at most one factor $q_0$ has $\mathrm{LC}(q_0) = 2$ ($|\prod_{Z(q_0)} w_j| = |q_0(0)|/2 \ge \frac{1}{2}$), so AM--GM gives $4\sum_{Z(q_0)} w_j^2 \ge 4\eta_0 (1/4)^{1/\eta_0} = \eta_0 4^{(\eta_0 - 1)/\eta_0}$. Since $t \mapsto 4t(1 - 4^{-1/t})$ has derivative $4 - 4^{1 - 1/t}(1 + \frac{\ln 4}{t}) > 0$ on $[1, \infty)$ (as $1 + u < e^u$ for $u = \frac{\ln 4}{t} > 0$), we obtain $\mathcal{T} \ge 4(\eta - \eta_0) + \eta_0 4^{(\eta_0 - 1)/\eta_0} \ge \eta 4^{(\eta - 1)/\eta}$. Thus $\eta^\eta 4^{\eta - \gamma(n)} \le \mathcal{T}^\eta$ in both parities, forcing $\eta \le \theta - 1$ and $m = d - \eta \ge d + 1 - \theta$.
\end{proof}

Table~\ref{tab:gsy_params} displays the exact integer parameters for all calibration and target pairs; note in particular that $(18, 59)$ forces only $m_{11} \ge 1$ ($17^{17} \cdot 4^{16} < 65^{17} < 18^{18} \cdot 4^{17}$), whereas $(18, 58)$ and $(19, 73)$ have $\theta = 22 > d$ and yield only the trivial lower bound $m_\kappa \ge 0$ (no guaranteed root at $\kappa$).

\begin{table}[htbp]
\centering
\caption{Exact multiplicity parameters $\kappa, \mathcal{T}, \theta$ and leading coefficients of $\phi(x) = Q(x)/(x-\kappa)^{m_\kappa}$.}\label{tab:gsy_params}
\footnotesize
\setlength{\tabcolsep}{4.5pt}
\begin{tabular}{@{}lccccrrll@{}}
\toprule
Case $(d,n)$ & Parity & $\kappa$ & $\mathcal{T}$ & $\theta$ & Bound $m_\kappa \ge$ & $\deg\phi$ & Top $3$ of $\phi(x)$ & Top $3$ of $\phi(x-1)$ \\
\midrule
$(19, 75)$ & Odd & $15$ & $25$ & $8$ & $12$ & $7$ & $(1, -100, 4275)$ & $(1, -107, 4896)$ \\
$(17, 49)$ & Odd & $9$ & $49$ & $14$ & $4$ & $13$ & $(1, -124, 7074)$ & $(1, -137, 8640)$ \\
$(18, 60)$ & Even & $11$ & $48$ & $13$ & $6$ & $12$ & $(1, -144, 9486)$ & $(1, -156, 11136)$ \\
$\mathbf{(19, 74)}$ & \textbf{Even} & $\mathbf{15}$ & $\mathbf{52}$ & $\mathbf{14}$ & $\mathbf{6}$ & $\mathbf{13}$ & $(1, -185, 15774)$ & $(1, -198, 18072)$ \\
$\mathbf{(18, 59)}$ & \textbf{Odd} & $\mathbf{11}$ & $\mathbf{65}$ & $\mathbf{18}$ & $\mathbf{1}$ & $\mathbf{17}$ & $(1, -194, 17680)$ & $(1, -211, 20920)$ \\
$(18, 58)$ & Even & $11$ & $84$ & $22$ & $0$ & $18$ & $(1, -200, 18847)$ & $(1, -218, 22400)$ \\
$(19, 73)$ & Odd & $15$ & $81$ & $22$ & $0$ & $19$ & $(1, -270, 34497)$ & $(1, -289, 39528)$ \\
\bottomrule
\end{tabular}
\end{table}

% ============================================================================
\section{Exact \texorpdfstring{$2$}{2}-adic Seidel congruences, the Pfaffian theorem, and Hensel obstructions}\label{sec:twoadic}
% ============================================================================

\subsection{Even- and odd-order \texorpdfstring{$2$}{2}-adic coefficient valuations}
Write any symmetric Seidel matrix $S \in \mathbb{Z}^{n \times n}$ as $S = J_n - I_n - 2A$ with symmetric zero-diagonal $A \in \{0, 1\}^{n \times n}$. Using the rank-$1$ determinant update $\det(B - uv^T) = \det B - v^T \adj(B) u$ on $y := x + 1$, we have
\begin{equation}\label{eq:rank1_charpoly}
  \chi_S(y - 1) = \det(yI_n + 2A - \mathbf{1}\mathbf{1}^T) = \det(yI_n + 2A) - \mathbf{1}^T \adj(yI_n + 2A) \mathbf{1} = \sum_{k=0}^n b_k y^{n - k}.
\end{equation}

\begin{theorem}[$2$-adic valuations of $\chi_S(y - 1)$ and $Q(x)$]\label{thm:twoadic_val}
Let $S \in \mathbb{Z}^{n \times n}$ ($n \ge 1$) be a symmetric Seidel matrix, and when $(x + 5)^{n - d} \mid \chi_S(x)$ write $Q(x) = \chi_S(x)/(x + 5)^{n - d}$ and $h(y) = Q(y - 5) = \sum_{k=0}^d (-1)^k e_k y^{d - k}$.
\begin{enumerate}
  \item $b_0 = 1$, $b_1 = -n$, and $b_2 = 0$.
  \item If $n$ is \emph{even}, $2^k \mid b_k$ for all $0 \le k \le n$ (\emph{Strict Type~2}); equivalently, $2^n \mid \chi_S(2z - 1)$ and $2^d \mid Q(2z - 1)$ in $\mathbb{Z}[z]$, $2^k \mid e_k$ ($0 \le k \le d$), and $\chi_S(x) \equiv (x+1)^{n-1}(x - n + 1) \pmod 8$.
  \item If $n$ is \emph{odd}, $2^{2\lfloor k/2 \rfloor} \mid b_k$ for all $k \ge 3$ ($2^k \mid b_k$ for even $k$, $2^{k-1} \mid b_k$ for odd $k$); equivalently, $2^{n-1} \mid \chi_S(2z - 1)$ and $2^{d-1} \mid Q(2z - 1)$ in $\mathbb{Z}[z]$, $2^{k - (k \bmod 2)} \mid e_k$ ($0 \le k \le d$), and $\chi_S(x) \equiv (x+1)^{n-1}(x - n + 1) \pmod 4$.
\end{enumerate}
\end{theorem}

\begin{proof}
Write $\det(yI_n + 2A) = \sum_{k=0}^n c_k y^{n - k}$ with $c_k = 2^k \sum_{|I|=k} \det(A[I])$. For any subset $I$ of odd cardinality, $\det(A[I]) = \sum_{\sigma \in \mathfrak{S}_I} \mathrm{sgn}(\sigma) \prod_{i \in I} A_{i,\sigma(i)}$: permutations with fixed points vanish ($A_{ii} = 0$), and fixed-point-free permutations pair into $\{\sigma, \sigma^{-1}\}$ with $\sigma \neq \sigma^{-1}$ (no fixed-point-free involution exists on an odd set) and equal terms ($A^T = A$). Hence $2 \mid \det(A[I])$ for odd $|I|$, giving $2^{2\lceil k/2 \rceil} \mid c_k$ for all $k \ge 0$ ($c_0 = 1, c_1 = 0$). Next, by symmetry of $\adj(yI_n + 2A)$,
\[
  \mathbf{1}^T \adj(yI_n + 2A) \mathbf{1} = \sum_{i=1}^n \det(yI_{n-1} + 2A_{\hat{i},\hat{i}}) + 2 \sum_{1 \le i < j \le n} \adj(yI_n + 2A)_{ij}.
\]
The diagonal sum equals $\frac{d}{dy}\det(yI_n + 2A) = \sum_{k=0}^{n-1} (n - k)c_k y^{n-1-k}$, while for $i < j$, $[y^{n-k}]\adj(yI_n + 2A)_{ij}$ is a signed sum of $(k-1)\times(k-1)$ minors of $2A$ (divisible by $2^{k-1}$ for $k \ge 2$ and $0$ for $k = 1$), so doubling makes the off-diagonal contribution to $b_k$ divisible by $2^k$. Thus $b_1 = -n$, $b_2 = \frac{1}{2}(\tr(S+I)^2 - \tr((S+I)^2)) = \frac{1}{2}(n^2 - (n + n(n-1))) = 0$, and for $k \ge 3$, $b_k \equiv c_k - (n - k + 1)c_{k-1} \pmod{2^k}$:
\begin{itemize}
  \item If $k$ is even, $k - 1$ is odd, so $2^k \mid c_k$ and $2^k \mid c_{k-1}$, whence $2^k \mid b_k$ for all $n$.
  \item If $k$ is odd, $2^{k+1} \mid c_k$ and $2^{k-1} \mid c_{k-1}$, giving $2^k \mid b_k$ when $n$ is even ($2 \mid n - k + 1$) and $2^{k-1} \mid b_k$ when $n$ is odd.
\end{itemize}
Substituting $y = 2z$ ($x = 2z - 1$, $x + 5 = 2(z + 2)$) and canceling monic $(z + 2)^{n - d}$ by Gauss's lemma yields the divisibilities on $Q(2z - 1)$ and $h(2u) = Q(2u - 5)$ (noting for odd $n$ that $2^{-(n-1)}\chi_S(2z-1) \equiv z^{n-d}\,2^{-(d-1)}Q(2z-1) \pmod 2$ transfers even-index parity to $Q$ and $h$).
\end{proof}

\subsection{Proof of the odd-order \texorpdfstring{mod-$2^{n+1}$}{mod-2^(n+1)} Pfaffian congruence}
\begin{proof}[Proof of Theorem~\ref{thm:pfaffian}]
For $n = 1$, $S = (0)$ and $2^0 \chi_S(2z - 1) = 2z - 1 = (2z - 1)Q_0(z)^2$ holds with monic $Q_0(z) = 1$ of degree $(1 - 1)/2 = 0$. Suppose now that $n = 2m + 1 \ge 3$.

\emph{Step 1 (Switching to an Eulerian graph).} Write $S = J_n - I_n - 2A$. Let $U \subseteq \{1, \dots, n\}$ be the set of odd-degree vertices of $A$. By the handshaking lemma, $|U|$ is even and $|U^c| = n - |U|$ is odd. Switching $S$ on $U$ preserves $\chi_S(x)$ and alters vertex degrees of $A$ modulo $2$ by $|U^c| \equiv 1$ on $U$ and $|U| \equiv 0$ on $U^c$, yielding an Eulerian graph $A$ satisfying $A\mathbf{1} = 2w$ ($w \in \mathbb{Z}^n$, $\mathbf{1}^T w = |E(A)|$).

\emph{Step 2 (Unimodular-modulo-$2^{n+1}$ basis reduction).} Let $P := [\mathbf{1}, e_2 - e_1, \dots, e_n - e_1] \in \mathbb{Z}^{n \times n}$ and $V := [e_2 - e_1, \dots, e_n - e_1] \in \mathbb{Z}^{n \times 2m}$. Adding rows $2, \dots, n$ to row $1$ converts $P$ to an upper-triangular matrix with diagonal $(n, 1, \dots, 1)$, so $\det P = n$. Since $V^T \mathbf{1} = 0$, conjugating $M(z) := (2z - 1)I_n - S = 2(zI_n + A) - J_n$ by $P$ produces
\[
  P^T M(z) P = \begin{pmatrix} n(2z - n) + 4|E(A)| & 4 w^T V \\ 4 V^T w & 2 B(z) \end{pmatrix}, \qquad B(z) := V^T(zI_n + A)V \in \mathbb{Z}[z]^{2m \times 2m}.
\]
Expand $\det(P^T M(z) P)$ along row~$1$ and column~$1$: every term selecting $4(w^T V)_j$ in row~$1$ and $4(V^T w)_i$ in column~$1$ multiplies a $(2m - 1) \times (2m - 1)$ minor of $2B(z)$, which is divisible by $16 \cdot 2^{2m-1} = 2^{n+2}$. Meanwhile, $4|E(A)| \cdot \det(2B(z)) = 2^{n+1}|E(A)|\det B(z) \equiv 0 \pmod{2^{n+1}}$. Hence
\begin{equation}\label{eq:full_laplace}
  n^2 \chi_S(2z - 1) = \det(P^T M(z) P) \equiv 2^{n-1} n(2z - n) \det B(z) \pmod{2^{n+1}}.
\end{equation}

\emph{Step 3 (Pfaffian identity for $\det B(z) \pmod 4$).} Observe that $B(z)$ is symmetric of even order $2m$, and its diagonal entries are even:
\[
  B(z)_{ii} = (e_{i+1} - e_1)^T(zI_n + A)(e_{i+1} - e_1) = 2z - 2A_{1, i+1} \in 2\mathbb{Z}[z].
\]
Define the skew-symmetric matrix $K(z) \in \mathbb{Z}[z]^{2m \times 2m}$ by $K_{ii} := 0$, $K_{ij} := B_{ij}$ for $i < j$, and $K_{ij} := -B_{ij}$ for $i > j$. Then $B(z) = K(z) + 2L(z)$, where $L_{ij} = 0$ ($i < j$), $L_{ii} = B_{ii}/2$, and $L_{ij} = B_{ij} \equiv K_{ij} \pmod 2$ ($i > j$). By multilinearity of the determinant modulo $4$,
\[
  \det B(z) \equiv \det K(z) + 2\tr\!\big(\adj(K(z)) L(z)\big) \pmod{4}.
\]
Since $K(z)$ is skew-symmetric of even order $2m$, $\det K(z) = \Pf(K(z))^2$. For $i < j$, $B(z)_{ij} = z + A_{i+1, j+1} - A_{1, i+1} - A_{1, j+1} + A_{1, 1} = z + O(1)$, so $K(z) = z K_\infty + O(1)$ where $K_\infty$ is the $2m \times 2m$ skew-symmetric matrix with $(K_\infty)_{ij} = +1$ for all $i < j$. By standard Pfaffian Laplace expansion along row~$1$, $\Pf(K_\infty) = 1$ (inductively, deleting $\{1, k\}$ flips $k - 2$ signs and leaves $K_\infty$ of order $2m - 2$, so $\Pf(K_\infty) = \sum_{k=2}^{2m} (-1)^k \cdot 1 = 1$). Therefore $Q_0(z) := \Pf(K(z)) \in \mathbb{Z}[z]$ is monic of degree $m = (n - 1)/2$. Moreover, $\adj(K)^T = -\adj(K)$ and $\adj(K)_{ii} = 0$ (odd-order skew-symmetric determinants vanish), so in $\mathbb{Z}[z]$:
\[
  2m \det K = \tr(\adj(K) K) = \sum_{i > j} \adj(K)_{ji} K_{ij} + \sum_{i < j} \adj(K)_{ji} K_{ij} = 2 \sum_{i > j} \adj(K)_{ji} K_{ij}.
\]
Dividing by $2$ in $\mathbb{Z}[z]$ and using $L_{ij} \equiv K_{ij} \pmod 2$ for $i > j$ and $L_{ij} = 0$ for $i < j$ gives
\[
  \tr(\adj(K) L) = \sum_{i > j} \adj(K)_{ji} L_{ij} \equiv m \det K \pmod{2}.
\]
Consequently, $\det B(z) \equiv (1 + 2m)\det K(z) = n\,Q_0(z)^2 \pmod 4$. Substituting into~\eqref{eq:full_laplace} yields $n^2 \chi_S(2z - 1) \equiv n^2 \cdot 2^{n-1}(2z - n)Q_0(z)^2 \pmod{2^{n+1}}$; dividing by $2^{n-1}$ and canceling the odd unit $n^2 \equiv 1 \pmod 8$ completes the proof.
\end{proof}

\begin{corollary}[Shifted mod-$4$ coefficient constraints]\label{cor:cor43_full}
Write $\chi_S(x) = (x + 5)^{n - d}(x - \kappa)^{m_\kappa} \phi(x)$ for odd $n$ and $2F(w) := 2^{-(\deg\phi - 1)} \phi(2w + \kappa) \in \mathbb{Z}[w]$:
\begin{enumerate}
  \item For $(18, 59)$ ($\kappa = 11, m_\kappa = 1, \deg\phi = 17, z = w + 6$): $(z + 2)^{41}(z - 6) \equiv z^{42} \pmod 4$ forces $z^{21} \mid Q_0(z) \pmod 2$, yielding $2F(w) \equiv (2w + 1)\tilde{Q}(w)^2 \pmod 4$ ($\deg\tilde{Q} = 8$).
  \item For $(19, 73)$ ($\kappa = 15, m_\kappa = 0, \deg Q = 19, z = w + 8$): $(z + 2)^{54} \equiv z^{54} \pmod 4$ yields $2F(w) \equiv (2w - 1)\tilde{Q}(w)^2 \pmod 4$ ($\deg\tilde{Q} = 9$).
  \item For $(17, 49)$ ($\kappa = 9, m_\kappa = 4, \deg\phi = 13, z = w + 5$): $(z + 2)^{32}(z - 5)^4 \equiv z^{32}(z - 1)^4 \pmod 4$ yields $2F(w) \equiv (2w + 1)\tilde{Q}(w)^2 \pmod 4$ ($\deg\tilde{Q} = 6$).
\end{enumerate}
In particular, writing $2F(w) = \sum_{k=0}^{2r+1} c_k w^{2r+1-k}$, the coefficients $q_1, \dots, q_r \in \{0, 1\}$ of $\tilde{Q}(w) \bmod 2$ are fixed by depth $k = r + 1$, uniquely determining $c_k \bmod 4$ at every lower coefficient $r + 2 \le k \le 2r + 1$.
\end{corollary}

\begin{proposition}[Vertex-deletion $2$-adic lifting dichotomy]\label{prop:vdel_dichotomy}
Let $S \in \mathbb{Z}^{n \times n}$ be a Seidel matrix and $S_v := S[V \setminus \{v\}]$ any one-vertex deletion of order $n - 1$:
\begin{enumerate}
  \item \emph{Even $n \to$ odd $n - 1$ (e.g.\ $(19, 74) \to 73$ and $(18, 60) \to 59$):} $S_v$ has odd order $n - 1$, so Theorem~\ref{thm:pfaffian} applies directly to $\chi_{S_v}$, enforcing both $2^{n-2} \mid \chi_{S_v}(2z - 1)$ and the mod-$4$ square congruence $2^{-(n-2)}\chi_{S_v}(2z - 1) \equiv (2z - n + 1)Q_0(z)^2 \pmod 4$ (eliminating $2{,}744$ of $5{,}069$ vertex-deletion polynomials on $(19, 74)$), and activating the odd-order Eulerian gauge (Theorem~\ref{thm:eulerian}) on $S_v$.
  \item \emph{Odd $n \to$ even $n - 1$ (e.g.\ $(18, 59) \to 58$ and $(17, 49) \to 48$):} $S_v$ has even order $n - 1$, so Theorem~\ref{thm:twoadic_val}(2) enforces full Strict Type-2 divisibility $2^{n-1} \mid \chi_{S_v}(2z - 1)$ in $\mathbb{Z}[z]$, whereas odd-order Eulerian gauge congruences do \emph{not} apply to $S_v$.
\end{enumerate}
\end{proposition}

\subsection{Integral lattice discriminants and the Hensel \texorpdfstring{mod-$8$}{mod-8} subsumption theorem}
\begin{proposition}[Integral lattice discriminant and even-eigenvalue bounds]\label{prop:lattice_gkms}
Let $G = 5I_n + S = U^T U \succeq 0$ have rank $r \le d$, and let $\Lambda := \sum_{i=1}^n \mathbb{Z} u_i \subset \mathbb{R}^r$ be the integral lattice spanned by columns of $U$, with basis Gram matrix $Q_\Lambda = B^T B \in \mathbb{Z}^{r \times r}$ and discriminant $\det(\Lambda) = \det Q_\Lambda > 0$. Let $D := \prod_{i=1}^r \mu_i \in \mathbb{Z}_{>0}$ be the product of nonzero eigenvalues of $G$.
\begin{enumerate}
  \item $\det(\Lambda) \mid D$, $2^{r-1} \mid \det(\Lambda)$, and $\tilde{\Lambda}_0 := \frac{1}{\sqrt{2}}\{\sum c_i u_i : 2 \mid \sum c_i\}$ is an even integral lattice of rank $r$ with $\det(\tilde{\Lambda}_0) = \det(\Lambda)/2^{r-2} \in 2\mathbb{Z}_{>0}$. Moreover, $2^{2\lfloor r/2 \rfloor} \mid D$ (so $2^{18} \mid |Q(-5)|$ when $r = 18$).
  \item \emph{(Greaves--Koolen--Munemasa--Sz\"oll\H{o}si~\cite{GKMS2016})} For any even integer $k \in 2\mathbb{Z}$, $\mult_k(S) = 0$ if $n$ is even, and $\mult_k(S) \le 1$ if $n$ is odd.
\end{enumerate}
\end{proposition}

\begin{proof}
Writing $U = B M$ ($M \in \mathbb{Z}^{r \times n}$ primitive), Cauchy--Binet gives $D = \sum_{|I|=r} \det(G[I]) = \det(\Lambda)\sum_{|I|=r}(\det M_I)^2$. Since $G \equiv J_n \pmod 2$ has $\mathbb{F}_2$-rank $1$, $Q_\Lambda$ has $r - 1$ even invariant factors, so $2^{r-1} \mid \det(\Lambda)$, and Theorem~\ref{thm:twoadic_val} at $z = -2$ on each $r \times r$ principal submatrix gives $2^{2\lfloor r/2 \rfloor} \mid \det(G[I]) \mid D$. Part~(2) follows immediately from $\rk_{\mathbb{F}_2}(S - kI_n) = \rk_{\mathbb{F}_2}(J_n - I_n) = n - (n \bmod 2)$ as in Theorem~\ref{thm:neumann}.
\end{proof}

Finally, we examine a natural $2$-adic Hensel congruence on vertex-deleted characteristic polynomials and prove why it is automatically implied by Strict Type-2 divisibility.

\begin{theorem}[Hensel $2$-adic unit-eigenvector congruence and Strict Type-2 subsumption]\label{thm:weapon_b}
Let $n \ge 3$ be odd and let $S \in \mathbb{Z}^{n \times n}$ be a symmetric Seidel matrix.
\begin{enumerate}
  \item $\chi_S(x) \equiv x(x + 1)^{n - 1} \pmod 2$ lifts by Hensel's lemma to a unique monic irreducible factor $q_0 \mid \chi_S$ of multiplicity $1$ with $q_0(0)$ even, $q_0'(0)$ odd, and a unique root $\lambda_0 \in 2\mathbb{Z}_2$ satisfying $\lambda_0 \equiv n - 1 \pmod 4$ and $v_2(\chi_S'(\lambda_0)) = 0$. Every primitive $2$-adic eigenvector $v \in \mathbb{Z}_2^n$ of $S$ for $\lambda_0$ has $v_u \in \mathbb{Z}_2^\times$ for all $u$, forcing $v_2(\chi_{S_u}(\lambda_0)) = 0$ and
  \begin{equation}\label{eq:weapon_b}
    \chi_{S_u}(\lambda_0) \equiv \chi_S'(\lambda_0) \cdot n^{-1} \pmod{8} \qquad \forall u \in \{1, \dots, n\}.
  \end{equation}
  \item \emph{(Subsumption by Strict Type~2)} Let $m := n - 1$ (even) and let $g(x) = x^m + c_1 x^{m-1} + c_2 x^{m-2} + \dots + c_m \in \mathbb{Z}[x]$ be \emph{any} monic polynomial satisfying $c_1 = 0$, $c_2 = -\binom{m}{2}$, and Strict Type~2 ($2^m \mid g(2t - 1)$ in $\mathbb{Z}[t]$). Then for every $\lambda_0 \in 2\mathbb{Z}_2$,
  \begin{equation}\label{eq:strict_type2_mod8}
    g(\lambda_0) \equiv (\lambda_0 + 1)^{n - 2}(\lambda_0 + 2 - n) \pmod{8}.
  \end{equation}
  Consequently, $g(\lambda_0) \bmod 8$ is independent of the choice of $g$ and automatically satisfies~\eqref{eq:weapon_b} whenever $\sum_u \chi_{S_u} = \chi_S'$.
\end{enumerate}
\end{theorem}

\begin{proof}
(1) Reducing $(J_n - I_n)v \equiv S v = \lambda_0 v \equiv 0 \pmod 2$ gives $(\sum_j v_j)\mathbf{1} \equiv v \pmod 2$. Since $v$ is primitive ($v \not\equiv 0 \pmod 2$), $\sum_j v_j \equiv 1 \pmod 2$ and every $v_u \equiv 1 \pmod 2$, so $v_u^2 \equiv 1 \pmod 8$ and $\|v\|_2^2 \equiv n \pmod 8$. Over $\mathbb{Q}_2$, $\chi_{S_u}(\lambda_0) = \chi_S'(\lambda_0) v_u^2 / \|v\|_2^2 \equiv \chi_S'(\lambda_0) n^{-1} \pmod 8$. Moreover, in the Eulerian gauge $S_\varepsilon \mathbf{1} \equiv (n-1)\mathbf{1} \pmod 4$, writing $v = \mathbf{1} + 2w$ gives $\lambda_0 \equiv \lambda_0 n \equiv \mathbf{1}^T S_\varepsilon v \equiv (n-1)n \equiv n - 1 \pmod 4$.

(2) Write $P(t) := 2^{-m} g(2t - 1) = t^m + \sum_{j=1}^m a_j t^{m - j} \in \mathbb{Z}[t]$, so $g(x) = (x+1)^m + 2a_1(x+1)^{m-1} + 4a_2(x+1)^{m-2} + \sum_{j=3}^m 2^j a_j (x+1)^{m-j}$. Matching $c_1 = m + 2a_1 = 0$ gives $a_1 = -m/2 \in \mathbb{Z}$, and matching $c_2 = \binom{m}{2} + 2a_1(m - 1) + 4a_2 = -\binom{m}{2} + 4a_2 = -\binom{m}{2}$ gives $a_2 = 0$. Since $2^j \equiv 0 \pmod 8$ for $j \ge 3$, evaluating at $\lambda_0 \in 2\mathbb{Z}_2$ yields $g(\lambda_0) \equiv (\lambda_0 + 1)^m - m(\lambda_0 + 1)^{m-1} = (\lambda_0 + 1)^{n-2}(\lambda_0 + 2 - n) \pmod 8$. Summing over $n$ vertices when $\sum_u g_u = \chi_S'$ gives $n g(\lambda_0) \equiv \chi_S'(\lambda_0) \pmod 8$ (corroborated computationally across $126$ random odd-$n$ Seidel matrices, all $9$ calibrated $(18,57)$ matrices, and all $24$ remaining interlacing-compatible candidates of $(18,59)$ and $(17,49)$; see~\S\ref{sec:compute}).
\end{proof}

% ============================================================================
\section{Jacobi identities, Eulerian gauge laws, Galois orthogonality, and \texorpdfstring{$\SRG(75,32,10,16)$}{SRG(75,32,10,16)}}\label{sec:jacobi_galois}
% ============================================================================

\subsection{Ring-theoretic Jacobi complementary minor identities}
\begin{theorem}[Jacobi's identity for complementary principal submatrices]\label{thm:jacobi_ring}
Let $A \in R^{n \times n}$ be a square matrix over a commutative ring $R$, let $I \subseteq \{1, \dots, n\}$ have cardinality $k$ ($1 \le k \le n$), and set $I^c := \{1, \dots, n\} \setminus I$. Then in $R$:
\begin{equation}\label{eq:jacobi_ring}
  \det\!\big(\adj(A)[I]\big) = (\det A)^{k - 1} \det\!\big(A[I^c]\big)
\end{equation}
(with $\det(\emptyset) := 1$ for $k = n$). Over a field $K$ when $\det A \neq 0$, $\det(A[I^c]) = (\det A)\det(A^{-1}[I])$.
\end{theorem}

\begin{proof}
Conjugating $A$ and $B := \adj(A)$ by a common permutation matrix placing $I$ before $I^c$ preserves determinants and $AB = (\det A)I_n$. In block form,
\[
  \begin{pmatrix} A[I] & A[I, I^c] \\ A[I^c, I] & A[I^c] \end{pmatrix}
  \begin{pmatrix} B[I] & 0 \\ B[I^c, I] & I_{n-k} \end{pmatrix}
  = \begin{pmatrix} (\det A)I_k & A[I, I^c] \\ 0 & A[I^c] \end{pmatrix}.
\]
Taking determinants over the universal polynomial domain $\mathbb{Z}[a_{ij}]_{1 \le i,j \le n}$ gives $(\det A)\det(B[I]) = (\det A)^k \det(A[I^c])$; canceling the nonzero polynomial $\det A$ and specializing to $R$ proves~\eqref{eq:jacobi_ring}.
\end{proof}

\begin{corollary}[Reduced $1$-, $2$-, and $k$-vertex spectral identities]\label{cor:jacobi_seidel}
Let $S \in \mathbb{Z}^{n \times n}$ be symmetric with distinct eigenvalues $\lambda_1, \dots, \lambda_m$ and orthogonal eigenprojectors $E_1, \dots, E_m$ ($S = \sum_r \lambda_r E_r$, $(xI_n - S)^{-1} = \sum_r (x - \lambda_r)^{-1} E_r$).
\begin{enumerate}
  \item For $k = 1$ ($I = \{u\}$), $\chi_{S_u}(x) = \adj(xI_n - S)_{uu} = \chi_S(x)\sum_{r=1}^m \frac{(E_r)_{uu}}{x - \lambda_r}$ and $\sum_{u=1}^n \chi_{S_u}(x) = \chi_S'(x)$.
  \item For $k = 2$ ($I = \{u, v\}$, $u \neq v$), $\chi_S(x)\chi_{S_{uv}}(x) = \chi_{S_u}(x)\chi_{S_v}(x) - \big(\adj(xI_n - S)_{uv}\big)^2$, where $\adj(xI_n - S)_{uv} \in \mathbb{Z}[x]$ has degree $n - 2$, leading coefficient $+S_{uv} \in \{\pm 1\}$, and $[x^{n-3}]$ coefficient $(S^2)_{uv}$ (since $(xI_n - S)^{-1} = x^{-1}I_n + x^{-2}S + x^{-3}S^2 + \dots$).
  \item When $\chi_S(x) = (x+5)^{n-d} Q(x)$ ($1 \le k \le n - d$, allowing $\rk(5I_n + S) \le d$), $-5$ is at most a simple pole of $(xI_n - S)^{-1}$, so $(x+5)^{n-d-1}$ divides every entry of $\adj(xI_n - S)$ in $\mathbb{Z}[x]$. Writing $\chi_{S[I^c]} = (x+5)^{n-d-k} Q_I$ and $\adj(xI_n - S)[I] = (x+5)^{n-d-1} G_I(x)$ gives $Q(x)^{k-1} Q_I(x) = \det G_I(x)$; in particular, for $k = 2$, $Q_u(x)Q_v(x) - Q(x)Q_{uv}(x) = g_{uv}(x)^2$ in $\mathbb{Z}[x]$ with $\deg g_{uv} = d - 1$ and $\mathrm{LC}(g_{uv}) = S_{uv}$.
\end{enumerate}
\end{corollary}

\subsection{Multiplicity-preserving submatrix extraction (\texorpdfstring{$\mult_\lambda(S_v) \ge m$}{mult >= m})}
We strengthen vertex-extraction~\cite[Lem.~2.10]{GS2024} to retain all vertices whose deletion preserves or increases the multiplicity of $\lambda$:

\begin{proposition}[Multiplicity-preserving submatrix extraction]\label{prop:prop61}
Let $\lambda$ be an eigenvalue of $S \in \mathbb{Z}^{n \times n}$ of multiplicity $m = \mult_\lambda(S)$, and let $I := \{v \in \{1, \dots, n\} : \mult_\lambda(S_v) \ge m\}$ have cardinality $k < n$. Then the $(n - k) \times (n - k)$ principal submatrix $T := S[I^c]$ satisfies
\[
  \mult_\lambda(T) \ge m \qquad \text{and} \qquad \mult_\mu(T) \ge \max(0,\, \mult_\mu(S) - k) \quad \forall \mu \neq \lambda.
\]
Consequently, $\sum_{\nu \in \mathrm{Spec}(S)} \mult_\nu^{\min}(T)\,\nu^2 \le \tr(T^2) = (n - k)(n - k - 1)$ (further tightened by Cauchy--Schwarz $(\sum_{\mathrm{rem}}\nu_j)^2 / r_{\mathrm{rem}}$ on the remaining $r_{\mathrm{rem}}$ eigenvalues of $T$).
\end{proposition}

\begin{proof}
Set $A := S - \lambda I_n$, of rank $R = n - m$. If some $z \in \ker(A)$ had $z_v \neq 0$ for $v \in I$, then column $v$ of $A$ would be a linear combination of the other $n - 1$ columns; by symmetry $A^T = A$, row $v$ would likewise be a linear combination of the other $n - 1$ rows, forcing $\rk(A_{\hat{v},\hat{v}}) = \rk(A) = R$ and $\mult_\lambda(S_v) = (n - 1) - R = m - 1 < m$, contradicting $v \in I$. Hence every $z \in \ker(A)$ vanishes identically on $I$, so restriction to $I^c$ injects $\ker(A)$ into $\ker(T - \lambda I_{n-k})$, proving $\mult_\lambda(T) \ge m$. (Equivalently, $\mathrm{Res}_{x=\lambda}(xI_n - S)^{-1}_{vv} = (E_\lambda)_{vv} \ge 0$ vanishes iff $\mult_\lambda(S_v) \ge m$.) Iterating Cauchy's interlacing theorem $k$ times yields $\mult_\mu(T) \ge \max(0, \mult_\mu(S) - k)$ for all $\mu \neq \lambda$.
\end{proof}

\subsection{Unique odd-order Eulerian switching gauge and \texorpdfstring{mod-$4$}{mod-4} matrix polynomials}
\begin{theorem}[Uniqueness and mod-$4$ polynomial laws of the odd-$n$ Eulerian gauge]\label{thm:eulerian}
Let $S \in \mathbb{Z}^{n \times n}$ be a symmetric Seidel matrix of \emph{odd} order $n$.
\begin{enumerate}
  \item \emph{(Strict uniqueness)} There exists a \emph{unique} Seidel matrix $S_\varepsilon = D S D$ in the switching class of $S$ (where $D = \diag(\delta_1, \dots, \delta_n) \in \{\pm 1\}^n$ is unique up to $D \mapsto -D$) such that $S_\varepsilon \mathbf{1} \equiv (n - 1)\mathbf{1} \pmod 4$, equivalently $S_\varepsilon = J_n - I_n - 2A_\varepsilon$ with every vertex of $A_\varepsilon$ having even degree.
  \item \emph{(Universal mod-$4$ power law)} In this Eulerian gauge, $S_\varepsilon^2 \equiv I_n + (n - 2)J_n \pmod 4$ (so $(S_\varepsilon^2)_{uv} \equiv n - 2 \pmod 4$ for all $u \neq v$), and for every integer $m \ge 1$:
  \[
    S_\varepsilon^{2m} \equiv I_n + (n - 2)J_n \pmod{4}, \qquad S_\varepsilon^{2m+1} \equiv S_\varepsilon + (n - 2)(n - 1)J_n \pmod{4}.
  \]
  \item \emph{(Matrix polynomial mod-$4$ congruence)} For any $P(x) \in \mathbb{Z}[x]$, set $p_{\mathrm{ev}} := \frac{P(1) + P(-1)}{2} \in \mathbb{Z}$ and $p_{\mathrm{od}} := \frac{P(1) - P(-1)}{2} \in \mathbb{Z}$. Then
  \begin{equation}\label{eq:poly_mod4}
    P(S_\varepsilon) \equiv c_I I_n + c_S S_\varepsilon + c_J J_n \pmod{4},
  \end{equation}
  where $c_I = p_{\mathrm{ev}}$, $c_S = p_{\mathrm{od}}$, and $c_J = (n - 2)\big[(p_{\mathrm{ev}} - P(0)) + (n - 1)(p_{\mathrm{od}} - P'(0))\big]$.
\end{enumerate}
\end{theorem}

\begin{proof}
Existence of $A_\varepsilon \mathbf{1} = 2w$ was proved in Step~1 of Theorem~\ref{thm:pfaffian}. Suppose $D' = \diag(\delta_1, \dots, \delta_n)$ switches $S_\varepsilon$ across a cut $(T, T^c)$ with $\emptyset \subsetneq T \subsetneq \{1, \dots, n\}$. Vertex degrees modulo $2$ shift by $|T^c| \bmod 2$ on $T$ and by $|T| \bmod 2$ on $T^c$. Because $|T| + |T^c| = n$ is odd, exactly one of $|T|, |T^c|$ is odd, creating odd-degree vertices on one side of the cut; thus $D' = \pm I_n$, leaving $S_\varepsilon$ strictly unique.

Next, $S_\varepsilon \mathbf{1} = (n - 1)\mathbf{1} - 4w \equiv (n - 1)\mathbf{1} \pmod 4$, and expanding $S_\varepsilon = (J_n - I_n) - 2A_\varepsilon$ gives
\[
  S_\varepsilon^2 = I_n + (n - 2)J_n - 4(\mathbf{1}w^T + w\mathbf{1}^T - A_\varepsilon - A_\varepsilon^2) \equiv I_n + (n - 2)J_n \pmod{4}.
\]
Since $n$ is odd, $(n - 1)^2 \equiv 0 \pmod 4$, $J_n^2 = nJ_n$, and $n^2 - 2n + 2 = (n - 1)^2 + 1 \equiv 1 \pmod 4$, so $(I_n + (n-2)J_n)^2 \equiv I_n + (n-2)(n^2 - 2n + 2)J_n \equiv I_n + (n-2)J_n \pmod 4$, and multiplying by $S_\varepsilon$ (using $S_\varepsilon J_n \equiv (n-1)J_n \pmod 4$) proves~(2). Splitting $P(x) = P(0) + P'(0)x + \sum_{m \ge 1} a_{2m} x^{2m} + \sum_{m \ge 1} a_{2m+1} x^{2m+1}$ and collecting coefficients of $I_n$, $S_\varepsilon$, and $J_n$ yields~(3).
\end{proof}

\begin{remark}[Strict odd-$n$ domain guard]\label{rem:odd_guard}
The odd-order hypothesis in Theorem~\ref{thm:eulerian} is indispensable. For even $n$, every cut $(T, T^c)$ satisfies $|T| \equiv |T^c| \pmod 2$, so switching flips all $n$ degree parities simultaneously: an even-order Seidel matrix either admits \emph{no} Eulerian representative (mixed degree parities) or admits $2^{n-2}$ distinct Eulerian switchings. In particular, when descending from $(18, 59)$ to one-vertex deletions $S_v$ of even order $n_v = 58$, Theorem~\ref{thm:eulerian} does \emph{not} apply to $S_v$ and is never invoked on even-order submatrices (see~\S\ref{sec:compute}).
\end{remark}

\paragraph{Application to $M, N$ block constraints, row-sum vectors $k_u$, and global $K = K^T$ compatibility.}
For odd $n$, factor $\chi_S(x) = \mathrm{Gcd}_x(x)\,\mathrm{Sim}_x(x)\,\mathrm{Mult}_x(x)$ where $\mathrm{Gcd}_x = \gcd(\chi_S, \chi_S')$ and $\mathrm{Sim}_x$ is the product of multiplicity-$1$ irreducible factors. Define
\[
  M := \mathrm{Mult}_x(S_\varepsilon) = \sum_{\lambda \in Z(\mathrm{Sim}_x)} \mathrm{Mult}_x(\lambda) E_\lambda \in \mathbb{Z}^{n \times n}, \qquad N := S_\varepsilon M = (x\,\mathrm{Mult}_x)(S_\varepsilon) \in \mathbb{Z}^{n \times n}.
\]
By Theorem~\ref{thm:eulerian}(3), diagonal entries satisfy $M_{uu} \equiv c_I(\mathrm{Mult}_x) + c_J(\mathrm{Mult}_x) \pmod 4$, while off-diagonal entries ($u \neq v$) satisfy $M_{uv} \equiv c_J(\mathrm{Mult}_x) \pm c_S(\mathrm{Mult}_x) \pmod 4$ (and analogously for $N$). Partitioning vertices into one-vertex spectral classes $C_0, \dots, C_{m-1}$ of sizes $c_0, \dots, c_{m-1}$, each $u \in C_a$ has Eulerian row-sum vector $k_u \in \mathbb{Z}^m$ with $k_u[b] := \sum_{v \in C_b \setminus \{u\}} (S_\varepsilon)_{uv}$, subject to $|k_u[b]| \le c_b - \delta_{ab}$, $k_u[b] \equiv c_b - \delta_{ab} \pmod 2$, and $\sum_{b=0}^{m-1} k_u[b] \equiv n - 1 \pmod 4$. Whenever mod-$4$ congruences uniquely determine the Galois-trace signs of $(M_{bc}, N_{ac})$ across all $b$ (and $M_{cc}^{\mathrm{off}} = M_{cc}^{\mathrm{diag}}$ if $c_c > 1$), $(S_\varepsilon M)_{uw} = N_{uw}$ enforces the exact linear equation $\sum_{b=0}^{m-1} M_{bc} k_u[b] = N_{ac}$. Moreover, symmetry $S_\varepsilon = S_\varepsilon^T$ enforces $\sum_i y_{a,i} = c_a$ and $\sum_i k^{(a,i)}[b]\,y_{a,i} = \sum_j k^{(b,j)}[a]\,y_{b,j}$ ($y_{a,i} \in \mathbb{Z}_{\ge 0}$) over admissible $k_u$-profiles, tested by exact rational Farkas duality.

\subsection{Generalized Galois orthogonality in \texorpdfstring{degree-$e$}{degree-e} number fields}
\begin{theorem}[Generalized Galois orthogonality]\label{thm:galois}
Under the hypotheses of Theorem~\ref{thm:galois_intro} (specifically: $q(x) \in \mathbb{Z}[x]$ is monic irreducible of multiplicity $1$ in $\chi_S(x)$ and degree $e \ge 2$, $K = \mathbb{Q}[x]/(q(x))$, $f_{u_0}(\bar{\lambda}) \neq 0$, and $r_u \in K$ satisfies $r_u^2 = f_u(\bar{\lambda})/f_{u_0}(\bar{\lambda})$ with $r_u = 0$ when $q \mid f_u$), the $\binom{e + 1}{2}$ rational linear equations~\eqref{eq:intro_galois} hold identically:
\[
  \sum_{u=1}^n \Big(\Tr_{K/\mathbb{Q}}(\bar{\lambda}^a r_u)\,\Tr_{K/\mathbb{Q}}(\bar{\lambda}^b r_u) - \Tr_{K/\mathbb{Q}}(\bar{\lambda}^{a+b} r_u^2)\Big) = 0 \qquad (0 \le a \le b \le e - 1).
\]
\end{theorem}

\begin{proof}
Because $S$ is real symmetric and $\mult_q(\chi_S) = 1$, $q(x)$ has $e$ distinct real roots $\lambda_1, \dots, \lambda_e \in \mathbb{R}$ ($\lambda_k \neq -5$ since $q$ has multiplicity $1$ and if $m \ge 2$ then $-5$ is repeated; if $m \le 1$, $\chi_S = (x+5)^m f$ with $q \mid f$ and $\mult_q(\chi_S) = 1$ gives $\chi_S'(\lambda_k) \neq 0$). Let $E_{\lambda_k} = v^{(k)}(v^{(k)})^T$ be the rank-$1$ orthogonal eigenprojector onto $\ker(S - \lambda_k I_n)$, so $\langle v^{(k)}, v^{(\ell)} \rangle = \delta_{k\ell}$ and $(v^{(k)}_u)^2 = \chi_{S_u}(\lambda_k)/\chi_S'(\lambda_k)$. Since $f_{u_0}(\bar{\lambda}) \neq 0$ in $K$, $(v^{(k)}_{u_0})^2 \neq 0$ for every $1 \le k \le e$.

By Corollary~\ref{cor:jacobi_seidel}(2)--(3), for each $u \in \{1, \dots, n\}$ there exists $g_{u, u_0}(x) \in \mathbb{Z}[x]$ such that $(E_{\lambda_k})_{u, u_0} = (\lambda_k + 5)^{\max(0, m-1)} g_{u, u_0}(\lambda_k) / \chi_S'(\lambda_k)$ and $g_{u, u_0}(x)^2 \equiv f_u(x)f_{u_0}(x) \pmod{q(x)}$. Working in the field $K = \mathbb{Q}[x]/(q(x))$, we have
\[
  \left(\frac{g_{u, u_0}(\bar{\lambda})}{f_{u_0}(\bar{\lambda})}\right)^2 = \frac{f_u(\bar{\lambda})}{f_{u_0}(\bar{\lambda})} = r_u(\bar{\lambda})^2.
\]
Because $K$ is a field, $g_{u, u_0}(\bar{\lambda})/f_{u_0}(\bar{\lambda}) = \eta_u r_u(\bar{\lambda})$ for a sign $\eta_u \in \{+1, -1\} \subset \mathbb{Q}$ (setting $\eta_u = +1$ if $r_u = 0$). Applying the $e$ ring embeddings $\sigma_k : K \hookrightarrow \mathbb{R}$ ($\bar{\lambda} \mapsto \lambda_k$) fixes $\eta_u \in \mathbb{Q}$, giving $g_{u, u_0}(\lambda_k)/f_{u_0}(\lambda_k) = \eta_u r_u(\lambda_k)$ and hence
\[
  v^{(k)}_u v^{(k)}_{u_0} = (E_{\lambda_k})_{u, u_0} = \frac{g_{u, u_0}(\lambda_k)}{f_{u_0}(\lambda_k)}\,(v^{(k)}_{u_0})^2 = \eta_u r_u(\lambda_k)\,(v^{(k)}_{u_0})^2 \implies v^{(k)}_u = \eta_u r_u(\lambda_k)\,v^{(k)}_{u_0}
\]
for all $1 \le k \le e$. For any $k \neq \ell$, dividing $\langle v^{(k)}, v^{(\ell)} \rangle = 0$ by $v^{(k)}_{u_0} v^{(\ell)}_{u_0} \neq 0$ and using $\eta_u^2 = 1$ cancels all vertex signs:
\[
  0 = \sum_{u=1}^n \eta_u^2\,r_u(\lambda_k) r_u(\lambda_\ell) = \sum_{u=1}^n r_u(\lambda_k) r_u(\lambda_\ell).
\]
Expanding $\Tr_{K/\mathbb{Q}}(\alpha) = \sum_{k=1}^e \sigma_k(\alpha)$ over all pairs $(k, \ell) \in \{1, \dots, e\}^2$ gives
\begin{align*}
  \sum_{u=1}^n \Tr_{K/\mathbb{Q}}(\bar{\lambda}^a r_u)\Tr_{K/\mathbb{Q}}(\bar{\lambda}^b r_u)
  &= \sum_{k,\ell=1}^e \lambda_k^a \lambda_\ell^b \sum_{u=1}^n r_u(\lambda_k) r_u(\lambda_\ell) \\
  &= \sum_{k=1}^e \lambda_k^{a+b} \sum_{u=1}^n r_u(\lambda_k)^2 = \sum_{u=1}^n \Tr_{K/\mathbb{Q}}(\bar{\lambda}^{a+b} r_u^2).\qedhere
\end{align*}
\end{proof}

\subsection{Closed-form \texorpdfstring{$2$}{2}-vertex Jacobi equations and two-vertex interlacing descent}
Let $\chi_S = \mathrm{Gcd}_x \cdot \mathrm{Sim}_x \cdot \mathrm{Mult}_{\mathrm{rad}}$ ($\deg(\mathrm{Sim}_x \cdot \mathrm{Mult}_{\mathrm{rad}}) = s$, $e_{\mathrm{sim}} := \deg\mathrm{Sim}_x$) and $\mathrm{Gcd}_2 := \gcd(\mathrm{Gcd}_x, \mathrm{Gcd}_x')$. For any pair of vertex-deletion polynomials $\chi_{S_u} = \mathrm{Gcd}_x \cdot g_a$ and $\chi_{S_v} = \mathrm{Gcd}_x \cdot g_b$ ($\deg g_a = s - 1$), Corollary~\ref{cor:jacobi_seidel}(2) reduces to
\begin{equation}\label{eq:reduced_2v}
  g_a(x) g_b(x) - \mathrm{Sim}_x(x)\,q_{uv}(x) = h_{uv}(x)^2, \qquad q_{uv} := \frac{\chi_{S_{uv}}}{\mathrm{Gcd}_2}, \qquad h_{uv} := \frac{\adj(xI_n - S)_{uv}}{\mathrm{Gcd}_x} \in \mathbb{Z}[x].
\end{equation}
Writing $h_{uv}(x) = Q_{uv}(x)\mathrm{Sim}_x(x) + R_{uv}(x)$ with $\deg R_{uv} < e_{\mathrm{sim}}$ and $\mathrm{diff} := \deg Q_{uv} = (s - 2) - e_{\mathrm{sim}} = \deg\mathrm{Mult}_{\mathrm{rad}} - 2 \in \{0, 1, 2, 3\}$:
\begin{itemize}
  \item $R_{uv}(x) \in \mathbb{Z}[x]$ is solved algebraically via CRT square roots $R_{uv}^2 \equiv g_a g_b \pmod{p_i}$ across the irreducible factors $p_i \mid \mathrm{Sim}_x$ (retaining only integral CRT lifts).
  \item Writing $\mathrm{Gcd}_x(x) = x^{\deg\mathrm{Gcd}_x} + \gamma_1 x^{\deg\mathrm{Gcd}_x - 1} + \dots$ and $S_{1,\mathrm{sim}} := -[x^{e_{\mathrm{sim}}-1}]\mathrm{Sim}_x$, Corollary~\ref{cor:jacobi_seidel}(2) fixes the top two coefficients of $h_{uv}$: $[x^{s-2}]h_{uv} = s_{uv} \in \{\pm 1\}$ and $[x^{s-3}]h_{uv} = (S^2)_{uv} - \gamma_1 s_{uv}$, so $Q_{uv}(x)$ equals $s_{uv}$ ($\mathrm{diff}=0$), $s_{uv}x + ((S^2)_{uv} - \gamma_1 s_{uv} + s_{uv}S_{1,\mathrm{sim}})$ ($\mathrm{diff}=1$), or has remaining low coefficient(s) bounded in exact $\mathbb{Q}$ from $h_{uv}(x_0)^2 \le g_a(x_0)g_b(x_0)$ at $x_0 = -5$ and $x_0 \ge x_{\mathrm{hi}} > \lambda_{\max}(S)$ where $\mathrm{Sim}_x(x_0)q_{uv}(x_0) \ge 0$.
  \item Every candidate $q_{uv} = (g_a g_b - h_{uv}^2)/\mathrm{Sim}_x \in \mathbb{Z}[x]$ is filtered without floating-point arithmetic by monicity, order-$(n-2)$ $2$-adic congruences (Theorem~\ref{thm:twoadic_val}(2) for even $n-2$; Theorem~\ref{thm:pfaffian} for odd $n-2$), exact Sturm real-rootedness in $[-5, \infty)$, and exact algebraic-number Cauchy interlacing against $\chi_{S_u}$ and $\chi_{S_v}$.
  \item On odd $n$, $(S_\varepsilon^2)_{uv} \equiv n - 2 \pmod 4$ and blockwise Cauchy--Schwarz $\sum_c (|tku_c + tkv_c| - n_c) \le (S_\varepsilon^2)_{uv} \le \sum_c (n_c - |tku_c - tkv_c|)$ ($n_c = c_c - \delta_{ac} - \delta_{bc}$, $tku_c = k_u[c] - \delta_{bc}s_{uv}$) mutually prune admissible row-sum profiles $k_u$ and $2$-vertex pairs to a fixed point; two-vertex interlacing descent at $v \in C_a$ further requires $\sum_{b} \sum_{q \in \mathcal{Q}_{ab}} x_{b,q}\,q(x) = \chi_{S_v}'(x)/\mathrm{Gcd}_2(x)$ with $\sum_q x_{b,q} = c_b - \delta_{ab}$ ($x_{b,q} \in \mathbb{Z}_{\ge 0}$) and $S_v$-Galois square-class consistency.
\end{itemize}

\subsection{Bidirectional equivalence of \texorpdfstring{$p_{136}$}{p\_136} with \texorpdfstring{$\SRG(75, 32, 10, 16)$}{SRG(75, 32, 10, 16)}}
\begin{theorem}[Bidirectional $\SRG(75, 32, 10, 16)$ equivalence]\label{thm:srg136}
A $74 \times 74$ symmetric Seidel matrix $S$ has characteristic polynomial
\[
  p_{136}(x) := (x + 5)^{55}(x - 15)^{17}(x - 11)(x - 9)
\]
if and only if $S$ is switching-equivalent to a $74 \times 74$ vertex-deleted principal submatrix of the Seidel matrix $S_{75} = J_{75} - I_{75} - 2A_{75}$ of a strongly regular graph $\SRG(75, 32, 10, 16)$.
\end{theorem}

\begin{proof}
($\Rightarrow$) Suppose $S \in \mathbb{Z}^{74 \times 74}$ is a Seidel matrix with $\chi_S = p_{136}$ and orthogonal eigenprojectors $E_{-5}, E_{15}, E_{11}, E_9$. Consider the matrix polynomial
\[
  R := \frac{1}{2}\big(75I_{74} + 10S - S^2\big).
\]
Since $p(x) = \frac{1}{2}(75 + 10x - x^2)$ takes values $p(-5) = 0$, $p(15) = 0$, $p(11) = 32$, and $p(9) = 42$, spectral decomposition yields $R = 32 E_{11} + 42 E_9 \succeq 0$ of rank $2$ with nonzero eigenvalues $32$ and $42$. Moreover, $R$ has integer entries: $R_{ii} = \frac{1}{2}(75 - 73) = 1$, and for $i \neq j$, $(S^2)_{ij} = \sum_{k \neq i,j} S_{ik}S_{kj}$ is a sum of $72$ products in $\{\pm 1\}$, hence even, so $R_{ij} \in \mathbb{Z}$. Positive semidefiniteness of every $2 \times 2$ principal minor $\left(\begin{smallmatrix} 1 & R_{ij} \\ R_{ij} & 1 \end{smallmatrix}\right) \succeq 0$ forces $R_{ij} \in \{-1, 0, 1\}$. Writing $R = U^T U$ with unit vectors $u_1, \dots, u_{74} \in \mathbb{R}^2$ satisfying $\langle u_i, u_j \rangle \in \{-1, 0, 1\}$ partitions $\{1, \dots, 74\} = V_1 \sqcup V_2$ into two nonempty orthogonal collinearity classes. Switching $S' = DSD$ ($D = \diag(\pm 1)$) so $u_i = +e_1$ on $V_1$ and $u_i = +e_2$ on $V_2$ gives $R' = D R D = J_{|V_1|} \oplus J_{|V_2|}$; matching nonzero eigenvalues $\{|V_1|, |V_2|\} = \{32, 42\}$ gives $|V_1| = 32$ and $|V_2| = 42$ after relabeling.

In block form $S' = \left(\begin{smallmatrix} B_{11} & B_{12} \\ B_{12}^T & B_{22} \end{smallmatrix}\right)$, $S'$ commutes with $R' = 32 E'_{11} + 42 E'_9$, so the eigenvectors $w_1 = (\mathbf{1}_{32}, \mathbf{0}_{42})^T$ (eigenvalue $32$ of $R'$) and $w_2 = (\mathbf{0}_{32}, \mathbf{1}_{42})^T$ (eigenvalue $42$ of $R'$) span $\mathrm{im}(E'_{11})$ and $\mathrm{im}(E'_9)$, forcing $S' w_1 = 11 w_1$ and $S' w_2 = 9 w_2$, i.e.,
\[
  B_{11}\mathbf{1}_{32} = 11\mathbf{1}_{32}, \qquad B_{12}^T\mathbf{1}_{32} = \mathbf{0}_{42}, \qquad B_{12}\mathbf{1}_{42} = \mathbf{0}_{32}, \qquad B_{22}\mathbf{1}_{42} = 9\mathbf{1}_{42}.
\]
Adjoin an apex vertex $\infty$ with border vector $w := (-\mathbf{1}_{32}, \mathbf{1}_{42})^T$:
\[
  S_{75} := \begin{pmatrix} 0 & w^T \\ w & S' \end{pmatrix}.
\]
Then $S_{75}\mathbf{1}_{75} = (-32 + 42,\; -\mathbf{1}_{32} + 11\mathbf{1}_{32},\; \mathbf{1}_{42} + 9\mathbf{1}_{42})^T = 10\mathbf{1}_{75}$. Using $(S')^2 = 75I_{74} + 10S' - 2(J_{32} \oplus J_{42})$ and $w w^T = 2(J_{32} \oplus J_{42}) - J_{74}$, we compute:
\begin{align*}
  (S_{75}^2)_{\infty, \infty} &= w^T w = 74 = 75 + 10(0) - 1, \\
  (S_{75}^2)_{\infty, V} &= w^T S' = (-11\mathbf{1}_{32}^T, 9\mathbf{1}_{42}^T) = 10 w^T - \mathbf{1}_{74}^T, \\
  (S_{75}^2)_{V, V} &= w w^T + (S')^2 = 75I_{74} + 10S' - J_{74},
\end{align*}
whence $S_{75}^2 = 75I_{75} + 10S_{75} - J_{75}$. Defining $A_{75} := \frac{1}{2}(J_{75} - I_{75} - S_{75}) \in \{0, 1\}^{75 \times 75}$ (symmetric, zero diagonal) gives $A_{75}\mathbf{1}_{75} = \frac{74 - 10}{2}\mathbf{1}_{75} = 32\mathbf{1}_{75}$ and
\[
  4A_{75}^2 = (J_{75} - I_{75} - S_{75})^2 = 64I_{75} - 24A_{75} + 64J_{75} \implies A_{75}^2 = 32I_{75} + 10A_{75} + 16(J_{75} - I_{75} - A_{75}),
\]
so $A_{75}$ is an $\SRG(75, 32, 10, 16)$.

($\Leftarrow$) Conversely, let $A_{75}$ be an $\SRG(75, 32, 10, 16)$ and $S_{75} := J_{75} - I_{75} - 2A_{75}$, satisfying $S_{75}\mathbf{1}_{75} = 10\mathbf{1}_{75}$ and $S_{75}^2 = 75I_{75} + 10S_{75} - J_{75}$. Deleting any vertex $\infty$ (with $32$ neighbors $V_1$, $42$ non-neighbors $V_2$, and border $w = (-\mathbf{1}_{32}, \mathbf{1}_{42})^T$) yields a $74 \times 74$ Seidel submatrix $S$ satisfying $S^2 = 75I_{74} + 10S - J_{74} - w w^T$, so $R := \frac{1}{2}(75I_{74} + 10S - S^2) = J_{32} \oplus J_{42}$, $Sw_1 = 11w_1$, $Sw_2 = 9w_2$, and $(S + 5I)(S - 15I) = 0$ on $\ker(J_{32} \oplus J_{42})$ (dimension $72$). From $\tr(S) = 0 = 11 + 9 + 15m_{15} - 5m_{-5}$ and $m_{15} + m_{-5} = 72$, we obtain $m_{15} = 17$ and $m_{-5} = 55$, giving $\chi_S = p_{136}$.
\end{proof}

% ============================================================================
\section{Higher-order spectral and structural obstructions for \texorpdfstring{$(19, 74)$}{(19, 74)} and \texorpdfstring{$(18, 59)$}{(18, 59)}}\label{sec:phase3_structural}
% ============================================================================

\subsection{Shared block-commutator simultaneous diagonalization}
\begin{lemma}[Block-commutator spectral quadratic reduction]\label{lem:block_comm}
Let $M = \left(\begin{smallmatrix} A & B \\ B^T & C \end{smallmatrix}\right)$ ($A \in \mathbb{R}^{n_1 \times n_1}$, $C \in \mathbb{R}^{n_2 \times n_2}$) be real symmetric and satisfy $AB + BC = c B$ ($c \in \mathbb{R}$).
\begin{enumerate}
  \item $(B^T B) C = C (B^T B)$, so every eigenspace $\ker(B^T B - \sigma^2 I_{n_2})$ is $C$-invariant and admits an orthonormal basis of simultaneous eigenvectors for $B^T B$ and $C$.
  \item If $e \in \mathbb{R}^{n_2} \setminus \{0\}$ satisfies $B^T B e = \sigma^2 e$ with $\sigma^2 > 0$ and $C e = \mu e$, then $M$ admits two nonzero eigenvectors $z_\pm = (Be,\, (\lambda_\pm - (c - \mu))e)^T$ with distinct eigenvalues $\lambda_+ > \lambda_-$ satisfying
  \[
    \lambda_+ + \lambda_- = c, \qquad \lambda_+ \lambda_- = (c - \mu)\mu - \sigma^2 \quad \Longleftrightarrow \quad \sigma^2 = (c - \mu)\mu - \lambda_+(c - \lambda_+).
  \]
  Moreover, pairwise orthogonal common eigenvectors $e_1 \perp e_2$ with $\sigma_1, \sigma_2 > 0$ produce pairwise orthogonal eigenvectors $z_{+,1} \perp z_{+,2}$ of $M$.
\end{enumerate}
\end{lemma}

\begin{proof}
(1) Substituting $BC = cB - AB$ gives $(B^T B)C = c(B^T B) - B^T A B$, which is symmetric since $A^T = A$; hence $(B^T B)C = ((B^T B)C)^T = C(B^T B)$, and restricting symmetric $C$ to each eigenspace of $B^T B$ yields a simultaneous orthonormal eigenbasis.
(2) Let $u_1 := Be$. Since $\|u_1\|_2^2 = e^T B^T B e = \sigma^2 \|e\|_2^2 > 0$, $u_1 \neq 0$, $A u_1 = (cB - BC)e = (c - \mu)u_1$, and $B^T u_1 = \sigma^2 e$. For $\lambda_\pm := \frac{1}{2}(c \pm \sqrt{(c - 2\mu)^2 + 4\sigma^2})$ ($\lambda_+ > \lambda_-$ since $\sigma^2 > 0$), we have $\sigma^2 + \mu(\lambda_\pm - (c - \mu)) = \lambda_\pm(\lambda_\pm - (c - \mu))$, so $M z_\pm = \lambda_\pm z_\pm$. If $e_1 \perp e_2$, then $\langle Be_1, Be_2 \rangle = \langle e_1, B^T B e_2 \rangle = \sigma_2^2 \langle e_1, e_2 \rangle = 0$, so $\langle z_{+,1}, z_{+,2} \rangle = 0$.
\end{proof}

\subsection{Elimination of \texorpdfstring{$p_{47}$}{p\_47} for \texorpdfstring{$(19, 74)$}{(19, 74)}}
\begin{proposition}[Unconditional elimination of $p_{47}$]\label{prop:kill47}
No $74 \times 74$ Seidel matrix $S$ realizes $p_{47}(x) = (x + 5)^{55}(x - 15)^{12}(x - 11)(x^2 - 28x + 191)^3$.
\end{proposition}

\begin{proof}
\emph{Step 1 (Block structure from interlacing).}
Here $\mathrm{Gcd}_x(x) = (x+5)^{54}(x-15)^{11}(x^2-28x+191)^2$ and $H_x(x) = \chi_S / \mathrm{Gcd}_x = (x+5)(x-15)(x-11)(x^2-28x+191)$ ($s = 5$, $\mathrm{Sim}_x = x - 11$, $\mathrm{diff} = 2$). Filtering $\Deck(p_{47})$ by Theorem~\ref{thm:pfaffian} at odd order $73$ leaves exactly $3$ reduced polynomials:
\begin{align*}
  g_0(x) &= (x - 11)(x^3 - 38x^2 + 469x - 1876), \\
  g_1(x) &= (x - 13)(x^3 - 36x^2 + 419x - 1576), \\
  g_2(x) &= (x - 15)(x - 11)(x^2 - 23x + 124).
\end{align*}
The $5 \times 3$ linear system $\sum_{j=0}^2 n_j g_j(x) = \chi_S'(x)/\mathrm{Gcd}_x(x)$ has full column rank $3$ and unique solution $(n_0, n_1, n_2) = (52, 16, 6)$, partitioning $V = C_0 \sqcup C_1 \sqcup C_2$ ($|C_0|=52, |C_1|=16, |C_2|=6$). Evaluating $(E_{11})_{uu} = \chi_{S_u}(11)/\chi_S'(11) = g_j(11)/H_x'(11)$ yields $g_0(11) = g_2(11) = 0$ and $(E_{11})_{uu} = \frac{1}{16}$ on $C_1$. Switching $S$ on $C_1$ so the unit $11$-eigenvector is $v_{11} = \frac{1}{4}(\mathbf{0}_{58}, \mathbf{1}_{16})^T$ places $S$ in block form $S = \left(\begin{smallmatrix} W & B \\ B^T & X \end{smallmatrix}\right)$ on $(C_0 \cup C_2) \sqcup C_1$ ($58 + 16$) with $X\mathbf{1}_{16} = 11\mathbf{1}_{16}$ and $B\mathbf{1}_{16} = \mathbf{0}_{58}$. Next, solving the exact $\mathrm{diff}=2$ $2$-vertex Jacobi equations~\eqref{eq:reduced_2v} ($x_{\mathrm{hi}} = 17 > 14 + \sqrt{5}$) yields $4$ valid solutions on pair $(0, 1)$ and $2$ valid solutions on pair $(1, 2)$; every solution satisfies $(S^2)_{wu} = 10 S_{wu}$, i.e.\ $(WB + BX)_{wu} = 10 B_{wu}$ for all $w \in C_0 \cup C_2$ and $u \in C_1$, so $WB + BX = 10B$.

\emph{Step 2 (Spectral contradiction).}
Apply Lemma~\ref{lem:block_comm} with $c = 10$. Because $X\mathbf{1}_{16} = 11\mathbf{1}_{16}$ and $B^T B \mathbf{1}_{16} = \mathbf{0}_{16}$, restrict to $\mathbf{1}_{16}^\perp$ (dimension $15$), choosing an orthonormal simultaneous eigenbasis $e_1, \dots, e_{15} \in \mathbf{1}_{16}^\perp$ with $B^T B e_i = \sigma_i^2 e_i$ and $X e_i = \mu_i e_i$.
\begin{itemize}
  \item Whenever $\sigma_i^2 > 0$, Lemma~\ref{lem:block_comm}(2) yields eigenvalues $\lambda_{+,i} + \lambda_{-,i} = 10$ in $\mathrm{Spec}(S) = \{-5, 11, 15, 14 \pm \sqrt{5}\}$. Direct pairwise inspection shows $\{\lambda_{+,i}, \lambda_{-,i}\} = \{15, -5\}$ uniquely (as $(14+\sqrt{5}) + (14-\sqrt{5}) = 28 \neq 10$ and $2(14 \pm \sqrt{5}) \notin \mathbb{Q}$), so $\lambda_{+,i}\lambda_{-,i} = -75$ and $\sigma_i^2 = 75 + 10\mu_i - \mu_i^2$.
  \item On $K_0 := \ker(B^T B) \cap \mathbf{1}_{16}^\perp$, every $e_i \in K_0$ satisfies $Be_i = 0$, so $(0, e_i)^T$ is an eigenvector of $S$ with eigenvalue $\mu_i \in \mathrm{Spec}(S)$. Since $X = J_{16} - I_{16} - 2A_1$ has $X\mathbf{1}_{16} = 11\mathbf{1}_{16}$, $A_1$ is $2$-regular, so $\mathrm{Spec}(X|_{\mathbf{1}^\perp}) = \{-1 - 2\theta_i\} \subseteq [-5, 3]$, intersecting $\mathrm{Spec}(S)$ solely at $\mu_i = -5$, where $75 + 10(-5) - (-5)^2 = 0 = \sigma_i^2$. (Equivalently, without invoking graph degrees, rationality of the orthogonal projector $P_{K_0} \in \mathbb{Q}^{16 \times 16}$ forces equal multiplicities of $14 \pm \sqrt{5}$ on $K_0$ via $\sqrt{5} \notin \mathbb{Q}$, while $75+10\mu-\mu^2 \ge 0$ at $-5, 11, 15$ and sums across conjugate pairs to $(75+10(14+\sqrt{5})-(14+\sqrt{5})^2) + (75+10(14-\sqrt{5})-(14-\sqrt{5})^2) = 28 > 0$, giving $\sum_{e_i \in K_0} \sigma_i^2 = 0 \le \sum_{e_i \in K_0} (75 + 10\mu_i - \mu_i^2)$.)
\end{itemize}
From $\tr(X) = 0 = 11 + \sum_{i=1}^{15} \mu_i$ and $\tr(X^2) = 16 \cdot 15 = 240 = 121 + \sum_{i=1}^{15} \mu_i^2$, we have $\sum_{i=1}^{15} \mu_i = -11$ and $\sum_{i=1}^{15} \mu_i^2 = 119$. Yet every entry of $B \in \{\pm 1\}^{58 \times 16}$ has $B_{wu}^2 = 1$, yielding the contradiction
\[
  928 = 58 \cdot 16 = \|B\|_F^2 = \tr(B^T B) = \sum_{i=1}^{15} \sigma_i^2 \le \sum_{i=1}^{15} (75 + 10\mu_i - \mu_i^2) = 15(75) + 10(-11) - 119 = 896.
\]
*(Independently, two-vertex odd-$73$ Galois Farkas duality at $v \in C_0, C_1, C_2$ also refutes $p_{47}$ with $0$ admissible cones.)*
\end{proof}

\subsection{Elimination of \texorpdfstring{$p_{134}, p_{74}, p_6, p_{24}$}{p\_134, p\_74, p\_6, p\_24} for \texorpdfstring{$(19, 74)$}{(19, 74)}}
\begin{proposition}[Elimination of $p_{134}, p_{74}, p_6, p_{24}$]\label{prop:other_19_74}
None of the four candidate polynomials $p_{134}, p_{74}, p_6, p_{24}$ of $(19, 74)$ is realizable:
\begin{enumerate}
  \item \emph{Candidate $p_{134}$, $\chi_S = (x+5)^{55}(x-17)(x-15)^{14}(x-13)^3(x-9)$:} Among $14$ $K^\times/(K^\times)^2$ cones over $\mathrm{Sim}_x = (x-17)(x-9)$, only cone $[1, 2, 16, 5, 0]$ is LP-feasible, with $81$ nonnegative integer solutions all satisfying $x_2 = 2 + 3x_5 + x_0/3 \ge 2$. Thus at least two vertices lie in spectral class~$2$, whereas exact $\mathrm{diff}=1$ $2$-vertex Jacobi on pair $(2, 2)$ ($2^{72}$ divisibility together with exact algebraic-number interlacing) admits $0$ valid $72$-vertex quotients.
  \item \emph{Candidate $p_{74}$, $\chi_S = (x+5)^{55}(x-15)^{13}(x-13)^4(x^2-28x+179)$:} Over $K = \mathbb{Q}(\sqrt{17})$, Theorem~\ref{thm:galois} rules out cone $[4, 5, 24, 0]$ by Farkas duality, while cone $[1, 3, 14, 15, 0]$ has $4$ integer solutions, all with $x_{15} \in \{3, 4, 5, 6\} > 0$. Deleting $v \in C_{15}$ yields an odd-$73$ submatrix $S_v$ whose Galois norm RREF over $72$-vertex quotients forces $2(x_1 + x_4) = 29$, impossible in $\mathbb{Z}$.
  \item \emph{Candidate $p_6$, $\chi_S = (x+5)^{55}(x-17)^3(x-15)^9(x-13)^6(x-11)$:} Admits $6$ one-vertex Galois solutions: the solution $[(1, 20), (2, 50), (6, 4)]$ has $0$ valid $2$-vertex Jacobi quotients on pair $(1, 2)$, while the five solutions $[(0, 20-t), (2, 42+3t), (4, 12-3t), (8, t)]$ ($t = 0, \dots, 4$) are all Farkas-infeasible under odd-$73$ submatrix Galois orthogonality upon deleting $v \in C_4$ ($t \le 3$) and $v \in C_2$ ($t = 4$).
  \item \emph{Candidate $p_{24}$, $\chi_S = (x+5)^{55}(x-17)(x-15)^9(x-13)^3(x^2-28x+191)^3$:} Admits $4$ one-vertex Galois solutions, all assigning positive count to deck polynomial $g_1(x) = (x-17)(x-16)(x-13)^2(x-9)$ (so $(x-13)^3 \mid \chi_{S_v}$ for $v \in C_1$). For any pair $(1, u)$, Cauchy interlacing forces $(x-13)^2 \mid q_{1u}$, so $g_1 g_u - (x-17)q_{1u} = \hat{h}_{1u}^2$ forces $(x-17)(x-13) \mid \hat{h}_{1u}$, reducing $\mathrm{diff} = 3$ to a single bounded integer parameter; odd-$73$ submatrix Galois Farkas duality at $v \in C_1$ ($\mathrm{Sim}_v = (x-17)(x-16)(x-9)$) then rules out all $4$ solutions (one solution also violates the $E_{17}$ triangle inequality $34 \le 32$).
\end{enumerate}
\end{proposition}

\subsection{Elimination of \texorpdfstring{$q_{861}$}{q\_861} for \texorpdfstring{$(18, 59)$}{(18, 59)}}
\begin{proposition}[Unconditional elimination of $q_{861}$]\label{prop:kill861}
No $59 \times 59$ Seidel matrix $S$ realizes $q_{861}(x) = (x + 5)^{41}(x - 13)^8(x - 11)^5(x - 9)^3(x^2 - 19x + 82)$.
\end{proposition}

\begin{proof}
\emph{Step 1 (Block structure from interlacing).}
Partitioning the $12$-element interlacing deck $\Deck(q_{861})$ into square classes over $\mathbb{Q}[x]/(x^2 - 19x + 82) \cong \mathbb{Q}(\sqrt{33})$ yields a unique nonnegative integer multiplicity vector: $30$ vertices in $C_0$ (deck type~$4$), $1$ vertex $u_1$ in $C_1$ (deck type~$5$), and $28$ vertices in $C_2$ (deck type~$9$). Applying Theorem~\ref{thm:eulerian} to $\mathrm{Mult}_x = \chi_S / (x^2 - 19x + 82)$ fixes $(M_{\mathrm{off}}, N_{\mathrm{off}}) \equiv (3, 2) \pmod 4$, uniquely determining all $3 \times 3$ block signs:
\[
  M_{\mathrm{blk}} = \begin{pmatrix} -17 & -29 & 11 \\ -29 & -57 & 15 \\ 11 & 15 & -9 \end{pmatrix}, \qquad
  N_{\mathrm{blk}} = \begin{pmatrix} -114 & -202 & 70 \\ -202 & -450 & 78 \\ 70 & 78 & -66 \end{pmatrix}.
\]
Solving $M_{\mathrm{blk}} k_u = N_{\mathrm{blk}, a}$ under parity, $|k_u[b]| \le c_b - \delta_{ab}$, and $\sum_b k_u[b] \equiv 58 \equiv 2 \pmod 4$ uniquely forces the Eulerian quotient matrix $K = \left(\begin{smallmatrix} 5 & 1 & 0 \\ 30 & 0 & 28 \\ 0 & 1 & 9 \end{smallmatrix}\right)$. In particular, $(S_\varepsilon)_{u_1, v} = +1$ for all $v \neq u_1$, and deleting $u_1$ yields $S_{u_1} = \left(\begin{smallmatrix} X & Y \\ Y^T & Z \end{smallmatrix}\right)$ ($30 + 28$) with $\chi_{S_{u_1}}(x) = (x+5)^{40}(x-5)(x-9)^4(x-11)^5(x-13)^8$. Next, exact $\mathrm{diff}=2$ $2$-vertex Jacobi ($x_{\mathrm{hi}} = 15$) on $(4, 9)$ combined with $(S_\varepsilon^2)_{uw} \equiv 57 \equiv 1 \pmod 4$ (Theorem~\ref{thm:eulerian}(2)) leaves only $(S_{uw}, (S_\varepsilon^2)_{uw}) \in \{(1, 9), (-1, -7)\}$, i.e.\ $(S_\varepsilon^2)_{uw} = 1 + 8(S_\varepsilon)_{uw}$ for all $u \in C_0, w \in C_2$. Since $(S_\varepsilon^2)_{uw} = (S_\varepsilon)_{u, u_1}(S_\varepsilon)_{u_1, w} + (S_{u_1}^2)_{uw} = 1 + (XY + YZ)_{uw}$, we obtain $XY + YZ = 8Y$.

\emph{Step 2 (Spectral contradiction).}
Apply Lemma~\ref{lem:block_comm} to $S_{u_1}$ with $c = 8$. Let $e_1, \dots, e_r \in \mathbb{R}^{28}$ ($r := \rk(Y) = \rk(Y^T Y)$) be an orthonormal basis of positive-eigenvalue eigenvectors of $Y^T Y$ simultaneously diagonalizing $Z$ ($Y^T Y e_i = \sigma_i^2 e_i$, $\sigma_i^2 > 0$, $Z e_i = \mu_i e_i$). By Lemma~\ref{lem:block_comm}(2), each $e_i$ induces eigenvalues $\lambda_{+,i} + \lambda_{-,i} = 8$ in $\mathrm{Spec}(S_{u_1}) = \{-5, 5, 9, 11, 13\}$. The unique pair summing to $8$ is $\{13, -5\}$, so $\lambda_{+,i} = 13$, $\lambda_{-,i} = -5$ ($\lambda_{+,i}\lambda_{-,i} = -65$), and the $r$ pairwise orthogonal vectors $z_{+,i}$ lie in $\ker(S_{u_1} - 13I_{58})$, forcing $r \le \mult_{13}(S_{u_1}) = 8$. Furthermore, for each $1 \le i \le r$,
\[
  \sigma_i^2 = (8 - \mu_i)\mu_i - (-65) = 81 - (\mu_i - 4)^2 \le 81.
\]
Since $Y \in \{\pm 1\}^{30 \times 28}$, $840 = 30 \cdot 28 = \|Y\|_F^2 = \sum_{i=1}^r \sigma_i^2 \le 8 \cdot 81 = 648$, contradiction.
\end{proof}

\subsection{Elimination of \texorpdfstring{$q_{488}$}{q\_488} for \texorpdfstring{$(18, 59)$}{(18, 59)}}
\begin{proposition}[Submatrix projector rank elimination of $q_{488}$]\label{prop:kill488}
No $59 \times 59$ Seidel matrix $S$ realizes $q_{488}(x) = (x + 5)^{41}\allowbreak (x - 15)\allowbreak (x - 13)^6\allowbreak (x - 11)^7\allowbreak (x - 10)\allowbreak (x - 9)^2\allowbreak (x - 7)$.
\end{proposition}

\begin{proof}
\emph{Step 1 (Projector structure from interlacing).}
Two-vertex interlacing descent ($59 \to 58 \to 57$) rules out $464$ of the $468$ one-vertex count vectors compatible with Proposition~\ref{prop:prop61} via linear Farkas duals, and $2$ further vectors fail $S_v$-Galois square-class compatibility, leaving two count vectors:
\begin{gather*}
  \mathbf{c}^{(0)} = [(3, 18), (5, 24), (41, 2), (43, 12), (154, 3)], \\
  \mathbf{c}^{(1)} = [(3, 20), (5, 24), (43, 10), (101, 2), (154, 3)].
\end{gather*}
Evaluating $(E_{11})_{uu} = \chi_{S_u}(11)/\chi_S'(11)$ at $\lambda = 11$ ($\rk E_{11} = \mult_{11}(S) = 7$) gives diagonal entries $\frac{1}{8}$ on types $\{3, 41\}$, $\frac{3}{16}$ on type~$5$, and $0$ on types $\{43, 101, 154\}$. Set $B := C_3 \cup C_{41}$ in $\mathbf{c}^{(0)}$ and $B := C_3$ in $\mathbf{c}^{(1)}$, so $|B| = 20$ and $|C_5| = 24$ ($B \cap C_5 = \emptyset$). Solving exact $2$-vertex Jacobi~\eqref{eq:reduced_2v} across all pairs shows $|(E_{11})_{uw}| = \frac{1}{16}$ for every $(u, w) \in B \times C_5$, $(E_{11})_{uz} = 0$ for every $u \in B$ and $z \notin B \cup C_5$, and $|(E_{11})_{uv}| \in \{0, \frac{1}{8}\}$ for $u \neq v$ in $B$.

\emph{Step 2 (Spectral contradiction).}
Let $T := E_{11}[B, B] \in \mathbb{R}^{20 \times 20}$. Since $T$ is a principal submatrix of $E_{11}$, $\rk(T) \le \rk(E_{11}) = 7$, while $\tr(T) = 20 \cdot \frac{1}{8} = \frac{5}{2}$. For each $u \in B$, idempotence $(E_{11}^2)_{uu} = (E_{11})_{uu} = \frac{1}{8}$ yields
\[
  \sum_{v \in B} T_{uv}^2 = \frac{1}{8} - \sum_{w \in C_5} (E_{11})_{uw}^2 - 0 = \frac{1}{8} - 24 \cdot \frac{1}{256} = \frac{1}{32},
\]
so $\tr(T^2) = \sum_{u \in B}\sum_{v \in B} T_{uv}^2 = 20 \cdot \frac{1}{32} = \frac{5}{8}$. (Note that with $T_{uu} = \frac{1}{8}$ and $T_{uv} \in \{0, \pm \frac{1}{8}\}$, $\sum_{v \neq u} T_{uv}^2 = \frac{1}{64}$ forces $T$ to decompose into $10$ disjoint $2 \times 2$ blocks $\frac{1}{8}\left(\begin{smallmatrix}1 & \pm 1 \\ \pm 1 & 1\end{smallmatrix}\right)$ of rank $10$.) Applying Cauchy--Schwarz $(\tr T)^2 \le \rk(T)\tr(T^2)$ gives $\frac{25}{4} = (\frac{5}{2})^2 \le 7 \cdot \frac{5}{8} = \frac{35}{8}$, contradiction.
\end{proof}

\subsection{Elimination of \texorpdfstring{$q_{441}$}{q\_441} for \texorpdfstring{$(18, 59)$}{(18, 59)}}
\begin{proposition}[Parseval $3$-clique projector orthogonality elimination of $q_{441}$]\label{prop:kill441}
No $59 \times 59$ Seidel matrix $S$ realizes $q_{441}(x) = (x + 5)^{41}(x - 15)^3(x - 11)^{13}(x - 10)(x - 7)$.
\end{proposition}

\begin{proof}
Two-vertex interlacing Farkas duality reduces all $23$ one-vertex Galois solutions to the single count vector $[(0, 49), (2, 2), (15, 8)]$. At $\lambda = 10$ ($\mult = 1$) and $\lambda = 15$ ($\mult = 3$), projector diagonals on types $(0, 2, 15)$ are $(E_{10})_{uu} = (\frac{1}{75}, \frac{3}{25}, \frac{1}{75})$ and $(E_{15})_{uu} = (\frac{1}{25}, \frac{3}{50}, \frac{23}{200})$. Because $(E_{10})_{uu} > 0$ for all $59$ vertices, switching $S$ so the unit $10$-eigenvector satisfies $v_{10} > 0$ fixes the sign $S_{uv} = \mathrm{sgn}((E_{10})_{uv}^{\mathrm{raw}})$ on every $2$-vertex Jacobi solution, restricting off-diagonal entries of $E_{15}$ in this gauge to:
\[
  E_{15}[C_0, C_0] \in \left\{-\frac{1}{100}, \frac{1}{25}\right\}, \quad
  E_{15}[C_0, C_2] \in \left\{-\frac{3}{100}, \frac{1}{50}\right\}, \quad
  E_{15}[C_0, C_{15}] \in \left\{-\frac{3}{50}, -\frac{1}{100}, \frac{1}{25}\right\}.
\]
On $C_0$ ($|C_0| = 49$), vectors in $\mathbb{R}^3$ of squared norm $\frac{1}{25}$ with inner products in $\{\frac{1}{25}, -\frac{1}{100}\}$ either coincide ($+\frac{1}{25}$) or meet at $\arccos(-1/4)$, partitioning $C_0$ into $k$ cliques of sizes $c_1 + \dots + c_k = 49$. Here $k = 1$ fails ($49/25 > 1$); $k = 2$ fails ($E_{15}^2 = E_{15} \succeq \mathrm{proj}$ requires $(25-c_1)(25-c_2) - \frac{1}{16}c_1 c_2 \ge 0$, which is strictly negative for all $1 \le c_1 \le 48$); and $k \ge 4$ fails ($\rk(\frac{5}{4}I_k - \frac{1}{4}J_k) = k > 3 = \rk E_{15}$). Hence $k = 3$ with clique Gram matrix $G_3 = \frac{5}{100}I_3 - \frac{1}{100}J_3$ ($\det G_3 > 0$). Every $u \in \{1, \dots, 59\}$ is determined in $\mathrm{im}(E_{15}) \cong \mathbb{R}^3$ by its $3$-clique inner-product vector $d_u \in \mathbb{Q}^3$ satisfying $d_u^T G_3^{-1} d_u = (E_{15})_{uu}$. Idempotence $E_{15}^2 = E_{15}$ requires the exact $3 \times 3$ rational Parseval identity $\sum_{u=1}^{59} d_u d_u^T = G_3$. Exhaustive rational enumeration over partitions $c_1 \ge c_2 \ge c_3 \ge 1$ ($c_1+c_2+c_3=49$), pairs in $d_2$, and $8$-multisets in $d_{15}$ yields a \emph{unique} Gram solution $(c_1, c_2, c_3) = (18, 16, 15)$; however, on this unique solution $\sqrt{75}\sum_{u=1}^{59} (v_{10})_u d_u \neq \mathbf{0}_3$, violating $E_{15} v_{10} = 0$.
\end{proof}

\subsection{Elimination of \texorpdfstring{$q_{692}$}{q\_692} for \texorpdfstring{$(18, 59)$}{(18, 59)}}
\begin{proposition}[De Caen--Jacobi $\mathbb{F}_2^4$ parity elimination of $q_{692}$]\label{prop:kill692}
No $59 \times 59$ Seidel matrix $S$ realizes $q_{692}(x) = (x + 5)^{41}(x - 15)^2(x - 11)^{14}(x^2 - 21x + 94)$.
\end{proposition}

\begin{proof}
One- and two-vertex interlacing compatibility reduces $q_{692}$ to the unique count vector $[(0, 17), (2, 40), (3, 2)]$. Let $r_{1,2} = \frac{1}{2}(21 \mp \sqrt{65})$ be the roots of $\mathrm{Sim}_x = x^2 - 21x + 94$. Direct exact evaluation in $\mathbb{Q}(\sqrt{65})$ shows $(E_{r_2})_{uu} > 0$ on all three types $\{0, 2, 3\}$. Switching $S$ into the $v_{r_2} > 0$ gauge uniquely fixes $(E_{r_1} + E_{r_2})[a, b]$ across $2$-vertex Jacobi pairs and determines row sums against $C_3$ ($|C_3| = 2$): $S[C_3, C_3] = J_2 - I_2$ ($k_{3,3} = -1$ forces $(S^2)_{w_1, w_2} \ge 17 > 9$, contradicting $2$-vertex Jacobi on $(3,3)$), $k_{3,0} = 15$, and $k_{3,2} = 16$.

Now examine $E_{11}$ ($\rk E_{11} = 14$) on the $19$-vertex set $V := C_0 \cup C_3$ ($V^c = C_2$, $|C_2| = 40$). On $V$, $(E_{11})_{uu}$ equals $\frac{1}{4}$ on $C_0$ and $\frac{3}{16}$ on $C_3$, $(E_{11})_{w_1, w_2} = -\frac{1}{16}$, and negative entries $-\frac{1}{8}$ define a graph $H$ on $V$ in which $w_1, w_2 \in C_3$ have degree $1$ (connected via a path $P_{m+2}$ of $m \ge 1$ internal $C_0$-vertices, as $m=0$ fails) and the remaining $17 - m$ vertices of $C_0$ form $c$ disjoint cycles $L_1, \dots, L_c$. Since $\rk(E_{11}[P_{m+2}]) = m + 1$ and each cycle drops rank by $1$, $14 \ge \rk(E_{11}[V]) = (m + 1) + \sum_{r=1}^c (|L_r| - 1) = 18 - c$, forcing $c \ge 4$. Moreover, for each cycle $L_r$, $\mathbf{1}_{L_r}^T E_{11} \mathbf{1}_{L_r} = 0 \implies (E_{11}\mathbf{1}_{L_r})|_{C_2} = \frac{1}{16} S[C_2, L_r]\mathbf{1}_{L_r} = 0$, so $|L_r| \ge 4$ is even. Thus $17 = m + \sum_{r=1}^c |L_r| \ge 1 + 4c \ge 17$, uniquely forcing $c = 4$, $|L_r| = 4$, and $m = 1$, i.e.\ $H = 4C_4 + P_3$.

Applying Jacobi's complementary minor identity (Theorem~\ref{thm:jacobi_ring}, interpolated at $42$ integer nodes $x = 20, \dots, 61$ in exact $\mathbb{Q}$) yields the exact characteristic polynomial of $S_{40} := S[C_2, C_2]$:
\[
  \chi_{S_{40}}(x) = (x - 15)^2 (x - 7)^{10} (x - 3)^5 (x + 5)^{23}.
\]
Let $M_{19} := S[C_2, V]S[V, C_2] \in \mathbb{Z}^{40 \times 40}$ and $w_r := S[C_2, L_r]\mathbf{1}_{L_r} \in \mathbb{Z}^{40}$ ($r = 0, 1, 2, 3$). Across all six $2$-vertex Jacobi solutions on $(2, 2)$ in the $v_{r_2} > 0$ gauge, $(15r_1 - 149)E_{r_1}[C_2, C_2] + (15r_2 - 149)E_{r_2}[C_2, C_2] = -J_{40}$, yielding the ambient block identity $(S^2 - 6S - 55I)|_{C_2} = S_{40}^2 + M_{19} - 6S_{40} - 55I_{40} = 80 E_{15}|_{C_2} - J_{40}$. Comparing this with the $S_{40}$ minimal-polynomial identity $S_{40}^2 - 2S_{40} - 35I_{40} = 160 E_{15}|_{C_2} - 32 E_3$ (where $E_{15}|_{C_2}$ persists at full rank $2$, and $E_3 = \frac{1}{16}\mathbf{1}_{++}\mathbf{1}_{++}^T + \frac{1}{64}\sum_{r=0}^3 w_r w_r^T$ on the $16$-vertex subset $C_{++} \subset C_2$ neighboring $C_3$ positively, with pairwise orthogonal cycle-projection vectors $w_0, \dots, w_3 \in \ker(S_{40} - 3I_{40})$ of norm $\|w_r\|_2^2 = 64$) gives $4S_{40} + 20I_{40} = 80 E_{15}|_{C_2} - 32 E_3 + J_{40} + M_{19}$. Reading diagonal entries and off-diagonal residues modulo $4$ assigns to each $u \in C_2$ a binary $4$-cycle imbalance indicator $W(u) \in \mathbb{F}_2^4$ (coordinate $r$ equals $1$ iff $u$ has nonzero sum $\pm 2$ on $L_r$, summing to $16$ per cycle across $C_2$) such that:
\begin{itemize}
  \item Every $u \in C_{++}$ ($|C_{++}| = 16$) has Hamming weight $\mathrm{wt}(W(u)) = 1$, and $W(u) \cdot W(v) \equiv 1 \pmod 2$ for $u \neq v$ in $C_{++}$; hence all $16$ vertices in $C_{++}$ share a single basis vector, say $e_0 = (1, 0, 0, 0)$.
  \item Every $v \in C_2 \setminus C_{++}$ ($24$ vertices) has $\mathrm{wt}(W(v)) = 2$, $W(u) \cdot W(v) \equiv 0 \pmod 2$ against $e_0$ (forcing $W(v)_0 = 0$), and $W(v) \cdot W(v') \equiv 0 \pmod 2$ pairwise. Yet any two distinct weight-$2$ vectors supported on $\{1, 2, 3\}$ intersect in $1$ coordinate ($W(v)\cdot W(v') \equiv 1 \pmod 2$), forcing all $24$ vertices to share a single vector (say $(0, 1, 1, 0)$).
\end{itemize}
Then coordinate $3$ has total weight $\sum_{u \in C_2} W(u)_3 = 0 \neq 16$, contradiction.
\end{proof}

\subsection{Elimination of \texorpdfstring{$q_{864}$}{q\_864} for \texorpdfstring{$(18, 59)$}{(18, 59)}}
\begin{proposition}[Projector clustering and Cauchy--Schwarz elimination of $q_{864}$]\label{prop:kill864}
No $59 \times 59$ Seidel matrix $S$ realizes $q_{864}(x) = (x + 5)^{41}(x - 13)^8(x - 11)^7(x - 9)^2(x - 6)$.
\end{proposition}

\begin{proof}
Two-vertex interlacing Farkas duality rules out $5$ of the $8$ one-vertex Galois solutions, leaving $3$ count vectors $(c_0, c_3, c_6, c_7) \in \{(34, 0, 21, 4),\, (33, 3, 18, 5),\, (32, 6, 15, 6)\}$ (note $c_3 + c_6 = 21$ in all three). Switching $S$ into the unique gauge where the simple $6$-eigenspace has $1155 E_6 = \gamma \gamma^T$ with $\gamma = (-3, 1, 5, 9)$ on types $(0, 3, 6, 7)$ determines exact $2$-vertex Jacobi candidates for $(21 E_9, 20 E_{11}, 63 E_{13})$.

On $(-C_3) \cup C_6$ (size $21$, flipping signs on $C_3$), $21 E_9$ ($\rk E_9 = 2$) has diagonal $2$ and off-diagonals in $\{-1, \frac{1}{2}, 2\}$ (unit circle vectors at angles $0^\circ, 120^\circ$, or $\arccos(1/4)$). If any pair realized $\frac{1}{2}$, all $21$ vectors would lie on at most two $\pm 120^\circ$-triads whose cross-Gram equations force $15k = 147$, impossible since $15 \nmid 147$. Thus $21 E_9 \in \{-1, 2\}$ on $(-C_3) \cup C_6$, partitioning it into $3$ clusters $K_m = A_m \sqcup B_m$ ($m = 0, 1, 2$, $A_m \subseteq -C_3$, $B_m \subseteq C_6$) of size $a_m + b_m = 7$. On each cluster $K_m$, the $2 \times 2$ Gram matrix of $(\mathbf{1}_{A_m}, \mathbf{1}_{B_m})$ under $63 E_{13} \succeq 0$ satisfies $q_{AA} = a_m(14 - 4a_m) \ge 0 \implies a_m \le 3$, and $q_{AA} q_{BB} - q_{AB}^2 \ge 0$ forces $e_{B_m} = 0$ and $k_{\mathrm{neg}}(A_m, B_m) = a_m b_m$ ($63 E_{13} = +1$ identically on $A_m \times B_m$). Similarly, $\mathbf{1}_{C_7}^T (63 E_{13}) \mathbf{1}_{C_7} = 14 c_7 - 2 c_7^2 - 14 e_7 \ge 0$ forces $e_7 \le 1$ ($c_7 \in \{4, 5\}$) and $e_7 = 0$ ($c_7 = 6$).

Now restrict $B := 9 E_{13} - 2I$ (eigenvalues $+7$ of multiplicity $\le 8 = \rk E_{13}$, and $-2$ on $\ker E_{13}$) to the orthogonal direct sum $W := \bigoplus_{m=0}^2 (\mathbf{1}_{A_m}^\perp \oplus \mathbf{1}_{B_m}^\perp) \oplus \mathbf{1}_{C_7}^\perp$. Since $\nu := \dim W - 8 > 0$ on every partition $(a_0, a_1, a_2)$, $B|_W$ has at least $\nu$ eigenvalues equal to $-2$; since $\tr(B|_W) = \frac{2e_7}{c_7} \ge 0$ ($0$ if $e_7 = 0$, $\frac{2}{c_7}$ if $e_7 = 1$), the remaining $\le 8$ eigenvalues sum to $\ge 2\nu$ and by Cauchy--Schwarz satisfy $\tr((B|_W)^2) \ge 4\nu + \frac{(2\nu)^2}{8} =: \mathrm{CS}_{\min}$. Bounding each off-diagonal block Frobenius norm of $B|_W$ via exact $2 \times 2$ positive-semidefinite minors (and $3 \times 3$ minors of $(\mathbf{1}_{A_m}, \mathbf{1}_{B_m}, \mathbf{1}_X)$ on $(2, 2, 2)$) yields $\tr((B|_W)^2) \le \mathrm{UB} < \mathrm{CS}_{\min}$ across all $7$ feasible partitions:
\begin{itemize}
  \item $(c_3, c_7) = (0, 4)$, $(0, 0, 0)$: $\mathrm{UB} = \frac{2675}{28} < \frac{273}{2}$;
  \item $(c_3, c_7) = (3, 5)$: $(0, 0, 3)$ gives $\frac{17058}{175} < \frac{273}{2}$; $(0, 1, 2)$ gives $\frac{102763}{1050} < 120$; $(1, 1, 1)$ gives $\frac{2234}{25} < \frac{209}{2}$;
  \item $(c_3, c_7) = (6, 6)$: $(0, 3, 3)$ gives $\frac{13133}{126} < \frac{273}{2}$; $(1, 2, 3)$ gives $\frac{6319}{60} < 120$; $(2, 2, 2)$ gives $\frac{2786}{25} < 120$.
\end{itemize}
Thus all $3$ remaining count vectors are impossible.
\end{proof}

% ============================================================================
\section{De Caen--Jacobi resolvent and exact rational elimination of \texorpdfstring{$p_{13}$}{p\_13}}\label{sec:idx13}
% ============================================================================

Among the candidate spectra of $(19, 74)$ other than $p_{136}$, the deepest structural analysis is required for candidate $p_{13}$,
\begin{equation}\label{eq:chi13}
  p_{13}(x) := (x + 5)^{55}(x - 17)(x - 15)^{15}(x - 11)^3.
\end{equation}
We present the first-principles reduction of $p_{13}$ from $74$ one-vertex multiplicity vectors to $3$ macroscopic block templates $S_1, S_2, S_3$, followed by the exact rational Schur-complement, $O(s_{\mathrm{rem}})$-anchor, and same-profile $\le 2$-copy multiplicity theory that rules out all three templates.

\subsection{Interlacing and triangle-inequality reduction \texorpdfstring{$74 \to 21 \to \{S_1, S_2, S_3\}$}{74 -> 21 -> {S1, S2, S3}}}
For $p_{13}$, $\mathrm{Gcd}_x(x) = (x+5)^{54}(x-15)^{14}(x-11)^2$ and $H_x(x) = (x+5)(x-17)(x-15)(x-11)$ ($s = 4$, $\mathrm{Sim}_x = x - 17$, $\mathrm{diff} = 1$). Filtering $\Deck(p_{13})$ by Theorem~\ref{thm:pfaffian} at odd order $73$ leaves $6$ reduced degree-$3$ polynomials:
\begin{align*}
  g_0(x) &= (x - 17)(x - 12)(x - 9), & g_1(x) &= x^3 - 38x^2 + 465x - 1832, \\
  g_2(x) &= (x - 13)(x^2 - 25x + 140), & g_3(x) &= x^3 - 38x^2 + 465x - 1816, \\
  g_4(x) &= x^3 - 38x^2 + 465x - 1804, & g_5(x) &= (x - 15)^2(x - 8).
\end{align*}
Evaluating at the simple eigenvalue $\lambda = 17$ ($H_x'(17) = 264$) gives $g_0(17) = 0$, $g_1(17) = 4$, $g_2(17) = 16$, $g_3(17) = 20$, $g_4(17) = 32$, and $g_5(17) = 36$. By Corollary~\ref{cor:jacobi_seidel}(3) over $K = \mathbb{Q}$, nonzero values $g_j(17)$ realized simultaneously in $S$ must lie in a single square class of $\mathbb{Q}^\times / (\mathbb{Q}^\times)^2$:
\begin{itemize}
  \item Square classes $5(\mathbb{Q}^\times)^2$ ($\{0, 3\}$) and $2(\mathbb{Q}^\times)^2$ ($\{0, 4\}$) fail $\sum_j c_j g_j(x) = p_{13}'(x)/\mathrm{Gcd}_x(x)$ immediately.
  \item Square class $1(\mathbb{Q}^\times)^2$ ($\{0, 1, 2, 5\}$, where $(E_{17})_{uu} = g_j(17)/264 \in \{0, \frac{1}{66}, \frac{4}{66}, \frac{9}{66}\}$ gives unit $17$-eigenvector magnitudes $|v_u| = \frac{r_u}{\sqrt{66}}$ with $r_u \in \{0, 1, 2, 3\}$) admits $74$ nonnegative integer solutions $(c_0, c_1, c_2, c_5)$.
\end{itemize}

\begin{lemma}[Reduction to $S_1, S_2, S_3$]\label{lem:idx13_74to3}
Any $74 \times 74$ Seidel realization of $p_{13}$ has one-vertex multiplicity vector $(c_0, c_1, c_2, c_5)$ in one of three templates:
\[
  S_1 = (24, 48, 0, 2), \qquad S_2 = (25, 45, 3, 1), \qquad S_3 = (26, 42, 6, 0).
\]
\end{lemma}

\begin{proof}
First, let $v \in \mathbb{R}^{74}$ be a unit $17$-eigenvector ($Sv = 17v$), with $|v_u| = r_u / \sqrt{66}$ ($r_u \in \{0, 1, 2, 3\}$ on $C_0, C_1, C_2, C_5$). For any vertex $u$ in a nonempty class $C_j$, $17 r_u = |\sum_{w \neq u} S_{uw} r_w| \le \sum_{w \neq u} r_w = (c_1 + 2c_2 + 3c_5) - r_u$, i.e.\ $18 r_u \le c_1 + 2c_2 + 3c_5$ (note $\sum_w r_w^2 = 66 \implies c_1 + 4c_2 + 9c_5 = 66$). This exact triangle inequality eliminates $53$ of the $74$ count vectors, leaving $21$ vectors (all having $c_0 \ge 12 > 0$). Second, deleting a vertex $u \in C_0$ ($\chi_{S_u} = \mathrm{Gcd}_x \cdot g_0$, $\mathrm{Sim}_u = (x-17)(x-12)(x-9)$) produces an odd-$73$ Seidel matrix $S_u$ whose $72$-vertex deck is restricted to admissible $\mathrm{diff}=1$ $2$-vertex Jacobi quotients $\mathcal{Q}_{0,j}$; applying Theorem~\ref{thm:pfaffian}, Theorem~\ref{thm:eulerian}, and Theorem~\ref{thm:galois} to $S_u$ refutes $18$ of the $21$ vectors via exact integer Farkas duals, leaving solely $S_1, S_2, S_3$.
\end{proof}

\subsection{De Caen--Jacobi resolvent, pole cancellation, and the \texorpdfstring{$120^\circ$-ray $3$-block}{120-degree-ray 3-block} split}
Let $m(x) := (x + 5)(x - 17)(x - 15)(x - 11) = x^4 - 38x^3 + 392x^2 + 230x - 14025$ be the minimal polynomial of $S$.

\begin{lemma}[Exact cubic resolvent and pole cancellation]\label{lem:resolvent_pole}
Define the matrix polynomial
\[
  R_{\mathrm{res}}(x, S) := S^3 + (x - 38)S^2 + (x^2 - 38x + 392)S + (x^3 - 38x^2 + 392x + 230)I_{74}.
\]
\begin{enumerate}
  \item $(xI_{74} - S)\,R_{\mathrm{res}}(x, S) = m(x)I_{74} - m(S) = m(x)I_{74}$ in $\mathrm{Mat}_{74}(\mathbb{Z}[x])$.
  \item For any subset $V \subseteq \{1, \dots, 74\}$ of cardinality $k := |V|$, Jacobi's complementary minor identity (Theorem~\ref{thm:jacobi_ring}) gives
  \begin{equation}\label{eq:idx13_complementary}
    \det\!\big(R_{\mathrm{res}}(x, S)[V, V]\big) = \frac{m(x)^k}{\chi_S(x)}\,\chi_{S[V^c, V^c]}(x)
    = \frac{(x+5)^{k-55}(x-17)^{k-1}(x-11)^{k-3}}{(x-15)^{15 - k}}\,\chi_{S[V^c, V^c]}(x).
  \end{equation}
  In particular, since $R_{\mathrm{res}}(x, S)[V, V] \in \mathrm{Mat}_k(\mathbb{Z}[x])$ has polynomial determinant, whenever $k > 15$ and $(x - 15) \nmid \chi_{S[V^c, V^c]}(x)$, the polynomial $(x - 15)^{k - 15}$ divides $\det(R_{\mathrm{res}}(x, S)[V, V])$ in $\mathbb{Z}[x]$, forcing $\rk_{\mathbb{R}}(R_{\mathrm{res}}(15, S)[V, V]) \le 15$.
\end{enumerate}
\end{lemma}

\begin{proof}
Part~(1) is the telescopic identity $m(x) - m(y) = (x - y)(y^3 + (x-38)y^2 + (x^2-38x+392)y + (x^3-38x^2+392x+230))$ evaluated at $y = S$, where $m(S) = 0$ since $p_{13}(x) \mid m(x)^{55}$ and $S$ is real symmetric. Part~(2) follows by taking $\adj(xI - S) = \frac{\chi_S(x)}{m(x)} R_{\mathrm{res}}(x, S)$ in Theorem~\ref{thm:jacobi_ring}: clearing denominators in~\eqref{eq:idx13_complementary} over the UFD $\mathbb{Z}[x]$ gives $(x - 15)^{k - 15} \mid \det(R_{\mathrm{res}}(x, S)[V, V])$, so at least $k - 15$ invariant factors vanish at $x = 15$.
\end{proof}

Next, write $5I_{74} + S = 16 E_{11} + 20 E_{15} + 22 E_{17}$ ($\rk E_{11} = 3$, $\rk E_{15} = 15$, $\rk E_{17} = 1$) and $R := \frac{1}{2}(75I_{74} + 10S - S^2) = 32 E_{11} - 22 E_{17} \in \mathbb{Z}^{74 \times 74}$. Fix the switching gauge where the unit $17$-eigenvector satisfies $v_{17} \ge 0$, so $66(E_{17})_{ab} = r_a r_b$ ($r = (0, 1, 2, 3)$ on $(C_0, C_1, C_2, C_5)$).

\begin{proposition}[$120^\circ$-ray $3$-block split of $C_1$ and $20E_{15}$ Gram projector]\label{prop:alg84}
In every template $S_\ell$ ($\ell \in \{1, 2, 3\}$, where $|C_1| = 3m_\ell$ with $m_1 = 16, m_2 = 15, m_3 = 14$):
\begin{enumerate}
  \item $E_{11}[C_0, C_1] = 0$, $\rk(E_{11}[C_0, C_0]) = 1$, and $\rk(E_{11}[C_1, C_1]) = 2$. Moreover, switching signs on $C_0$ sets $R[C_0, C_0] = J_{c_0}$, making the complement graph $A_0$ on $C_0$ $8$-regular in $S_1$ (and $(c_0 - 16)$-regular in general).
  \item On $C_1$, $(E_{11})_{uu} = \frac{1}{24}$ and $(E_{11})_{uv} \in \{+\frac{1}{24}, -\frac{1}{48}\}$, partitioning $C_1 = B_1 \sqcup B_2 \sqcup B_3$ into three blocks of equal size $|B_1| = |B_2| = |B_3| = m_\ell \in \{16, 15, 14\}$.
  \item Defining $M := 20 E_{15} = 5I_{74} + S - 16 E_{11} - 22 E_{17} \succeq 0$ ($\rk M = 15$), every full-rank factorization $M = Z Z^T$ ($Z \in \mathbb{R}^{74 \times 15}$) satisfies $Z^T Z = 20 I_{15}$.
\end{enumerate}
\end{proposition}

\begin{proof}
(1) Solving $\mathrm{diff}=1$ $2$-vertex Jacobi~\eqref{eq:reduced_2v} on pair $(0, 1)$ gives $(S^2)_{uv} = 10 S_{uv}$, i.e.\ $R_{uv} = 0$, for all $(u, v) \in C_0 \times C_1$. Since $E_{17}|_{C_0} = 0$, $0 = R[C_0, C_1] = 32 E_{11}[C_0, C_1]$. Evaluating $H_x'(11) = 16(-6)(-4) = 384$, $g_0(11) = (-6)(-1)(2) = 12$, $g_1(11) = 16$, and $g_5(11) = 48$ gives $(E_{11})_{uu} = \frac{1}{32}$ on $C_0$, $\frac{1}{24}$ on $C_1$, and $\frac{1}{8}$ on $C_5$ (summing on $S_1$ to $\frac{24}{32} + \frac{48}{24} + \frac{2}{8} = 3 = \rk E_{11}$). On $C_0$, $R_{uu} = 32(\frac{1}{32}) = 1$ and $R_{uv} \in \{-1, 0, 1\}$, while on $C_1$, $\tr(E_{11}[C_1, C_1]) = |C_1|/24 \in \{2, \frac{15}{8}, \frac{7}{4}\} > 1$. Since $E_{11}[C_0, C_1] = 0$ and $\rk E_{11} = 3$, $\rk(E_{11}[C_1, C_1]) \ge 2$ ($\tr > 1$ excludes rank $1$ as $E_{11} \preceq I$) and $\rk(E_{11}[C_0, C_0]) \ge 1$, forcing $\rk(E_{11}[C_0, C_0]) = 1$ and $\rk(E_{11}[C_1, C_1]) = 2$. Thus $32 E_{11}[C_0, C_0] = z z^T$ ($z \in \{\pm 1\}^{c_0}$); switching $C_0$ so $z = \mathbf{1}_{c_0}$ sets $32 E_{11}[C_0, C_0] = J_{c_0}$, so $(c_0 - 1 - 2\deg_{A_0}(u))$-regularity gives $\deg_{A_0}(u) = 8$ in $S_1$ ($c_0 = 24$).

(2) On $C_1 \times C_1$, $2$-vertex Jacobi~\eqref{eq:reduced_2v} restricts $R_{uv} = 32(E_{11})_{uv} - \frac{1}{3}$ to $\{+1, -1\}$: specifically, $R_{uv} = +1 \implies (E_{11})_{uv} = \frac{4/3}{32} = +\frac{1}{24}$, and $R_{uv} = -1 \implies (E_{11})_{uv} = \frac{-2/3}{32} = -\frac{1}{48}$ ($R_{uv} = 0$ would give $(E_{11})_{uv} = \frac{1}{96}$, excluded by $2$-vertex Jacobi on $(1,1)$). Since $\rk(E_{11}[C_1, C_1]) = 2$ and $v_{17} > 0$ is already fixed on $C_1$ (no further sign switching allowed on $C_1$), unit vectors in $\mathbb{R}^2$ with pairwise inner products in $\{+1, -\frac{1}{2}\}$ lie on three $120^\circ$ rays $e^{(1)}, e^{(2)}, e^{(3)} \in S^1$ ($e^{(1)} + e^{(2)} + e^{(3)} = 0$), partitioning $C_1 = B_1 \sqcup B_2 \sqcup B_3$ of sizes $b_1 + b_2 + b_3 = 3m_\ell$. Using $E_{11}^2 = E_{11}$ on $C_1$:
\begin{itemize}
  \item In $S_1$ ($3m_1 = 48$, $\tr(E_{11}[C_1]) = 2 = \rk(E_{11}[C_1])$), $E_{11}[C_1, V \setminus C_1] = 0$, so $E_{11}[C_1, C_1]^2 = E_{11}[C_1, C_1]$ forces $\frac{1}{24}b_a^2 + \frac{1}{96}\sum_{c \neq a} b_a b_c = b_a$---equivalently $\sum_{a=1}^3 b_a e^{(a)} (e^{(a)})^T = 24 I_2$---whose trace and determinant uniquely force $b_1 = b_2 = b_3 = 16$.
  \item In $S_2$ ($3m_2 = 45$) and $S_3$ ($3m_3 = 42$), matching eigenvalues of $\sum_{a=1}^3 \frac{b_a}{24} e^{(a)}(e^{(a)})^T \preceq I_2$ against the rank-$2$ tail on $C_2 \cup C_5$ forces $b_1 = b_2 = b_3 = 15$ and $b_1 = b_2 = b_3 = 14$, respectively.
\end{itemize}
(3) From $M^2 = 400 E_{15}^2 = 20 M$ and $M = Z Z^T$ with $\rk Z = 15$, left- and right-multiplying $(Z Z^T)(Z Z^T) = 20 Z Z^T$ by $(Z^T Z)^{-1} Z^T$ and $Z(Z^T Z)^{-1}$ yields $Z^T Z = 20 I_{15}$.
\end{proof}

\begin{proposition}[Even-cycle zero-sum lock]\label{prop:alg85}
In $S_1$, $16(E_{11})_{uv} + 22(E_{17})_{uv}$ equals $16(\frac{1}{24}) + 22(\frac{1}{66}) = 1$ within each $B_a$ ($a \in \{1,2,3\}$) and $16(-\frac{1}{48}) + 22(\frac{1}{66}) = 0$ between distinct blocks $B_a \neq B_b$, so
\begin{gather*}
  M[B_a, B_a] = 2(2I_{16} - A_a), \qquad M[B_a, B_b] = S[B_a, B_b] \in \{\pm 1\}^{16 \times 16}, \\
  M[C_0, B_a] = S[C_0, B_a] \in \{\pm 1\}^{24 \times 16},
\end{gather*}
where each $A_a$ ($S[B_a, B_a] = J_{16} - I_{16} - 2A_a$, row sum $11$) is a $2$-regular simple graph on $16$ vertices. For every cycle component $L \subseteq B_a$, $\mathbf{1}_L^T M[B_a, B_a] \mathbf{1}_L = 0$; since $M \succeq 0$, $M \mathbf{1}_L = 0$, forcing $\sum_{u \in L} S_{x, u} = 0$ for every $x \in C_0 \cup \bigcup_{b \neq a} B_b$. Because $S_{x, u} \in \{\pm 1\}$, $|L| \ge 4$ is even, restricting each $B_a$ in $S_1$ to the $7$ even partitions of $16$:
\[
  (16), \quad (12, 4), \quad (10, 6), \quad (8, 8), \quad (8, 4, 4), \quad (6, 6, 4), \quad (4, 4, 4, 4).
\]
Analogously, every cycle component inside the core block $K$ of $S_2$ ($|K| = 18 = 15 + 3$) and $S_3$ ($|K| = 20 = 14 + 6$) has even length $\ge 4$.
\end{proposition}

\begin{proof}
Directly from $A_a \mathbf{1}_L = 2\mathbf{1}_L \implies M[B_a, B_a]\mathbf{1}_L = 2(2I_{16} - A_a)\mathbf{1}_L = 0$, positive-semidefinite kernel vanishing $z^T M z = 0 \implies Mz = 0$, and parity of sums of $\pm 1$ entries.
\end{proof}

\subsection{Exact \texorpdfstring{$\mathbb{Q}$-algebraic coordinates, $O(s_{\mathrm{rem}})$ anchors, and the same-profile $\le 2$-copy cap}{Q-algebraic coordinates, O(s_rem) anchors, and the same-profile <= 2-copy cap}}
Let $K$ be the core vertex set ($K = B_1$ of size $m = 16$ in $S_1$; $K = B_1 \cup C_2$ of size $18$ in $S_2$; $K = B_1 \cup C_2$ of size $20$ in $S_3$), with rational core matrix $M_{KK} = M[K, K] \in \mathbb{Q}^{m \times m}$ of rank $r := \rk(M_{KK}) \le 15$ and complementary dimension $s_{\mathrm{rem}} := 15 - r \in \{0, 1, 2, 3\}$.

\begin{theorem}[Exact $\mathbb{Q}$-coordinate Schur-complement system and squarefree radical splitting]\label{thm:alg86_87}
Let $M_{KK}^+ \in \mathbb{Q}^{m \times m}$ be the exact rational Moore--Penrose pseudoinverse of $M_{KK}$, let $\Pi \subseteq \{1, \dots, m\}$ ($|\Pi| = r$) be the RREF pivot column set of $M_{KK}$, and set $W := 20 M_{KK} - M_{KK}^2 \in \mathbb{Q}^{m \times m}$ and $W_\Pi := W[\Pi, \Pi] \in \mathbb{Q}^{r \times r}$. For each outside vertex $x \notin K$, write $v_x := M[K, x] \in \mathbb{Q}^m$ ($v_x \in \mathrm{col}(M_{KK})$), $u_x := v_x[\Pi] \in \mathbb{Q}^r$, $g(x, y) := v_x^T M_{KK}^+ v_y \in \mathbb{Q}$, $Q(v_x) := g(x, x) \in \mathbb{Q}$, and $\mathrm{res}(x) := M_{xx} - Q(v_x) \in \mathbb{Q}_{\ge 0}$ ($M_{xx} = \frac{9}{2}$ on $C_0$, $M_{xx} = 4$ on $B_2 \cup B_3$).
\begin{enumerate}
  \item Writing $M_{KK} = U_r \Lambda U_r^T$ ($\Lambda \succ 0$) and rotating $Z \in \mathbb{R}^{74 \times 15}$ by $O(15)$ so $Z_K = (U_r \Lambda^{1/2},\, 0_{m \times s_{\mathrm{rem}}})$ splits each row $z_x = (p_x, r_x) \in \mathbb{R}^r \times \mathbb{R}^{s_{\mathrm{rem}}}$ ($x \notin K$) with $\tilde{T} p_x = u_x$ ($\tilde{T} := \Pi_{\mathrm{row}} U_r \Lambda^{1/2}$ invertible), $\langle p_x, p_y \rangle = g(x, y)$, and $\|r_x\|_2^2 = \mathrm{res}(x)$. In these rational coordinates, $Z^T Z = 20 I_{15}$ is equivalent to
  \begin{equation}\label{eq:qq_gram_master}
    \sum_{x \notin K} u_x u_x^T = W_\Pi, \qquad \sum_{x \notin K} u_x r_x^T = 0, \qquad \sum_{x \notin K} r_x r_x^T = 20 I_{s_{\mathrm{rem}}}.
  \end{equation}
  In particular, for every integer test vector $w \in \mathbb{Z}^m$, $\sum_{x \in C_0} (w^T v_x)^2 \le w^T W w$ holds over $\mathbb{Z}/\mathbb{Q}$ with zero spectral approximation.
  \item Fix canonical anchors in $C_0$ under $O(s_{\mathrm{rem}})$:
  \begin{itemize}
    \item If $s_{\mathrm{rem}} = 1$, $r_x = \varepsilon_x q_x \sqrt{d_x}$ ($\varepsilon_x \in \{0, \pm 1\}$, $\mathrm{res}(x) = q_x^2 d_x$ with $q_x \in \mathbb{Q}_{\ge 0}$ and squarefree $d_x \in \mathbb{Z}_{\ge 1}$).
    \item If $s_{\mathrm{rem}} = 2$ and $\max_{x \in C_0} \mathrm{res}(x) = d_0 > 0$ is attained at $s_0 \in C_0$, rotating $O(2)$ so $r_{s_0} = (\sqrt{d_0}, 0)$ determines $A_x := M_{s_0, x} - g(s_0, x) \in \mathbb{Q}$ ($a_x = A_x / \sqrt{d_0}$) and $\mathrm{rem}_1(x) := \mathrm{res}(x) - A_x^2 / d_0 = q_x^2 d_x \ge 0$, leaving one sign $b_x = \varepsilon_x q_x \sqrt{d_x}$.
    \item If $s_{\mathrm{rem}} = 3$ ($S_1(4,4,4,4)$, $r = 12$) and $\max_{x \in C_0} \mathrm{rem}_1(x) = d_1 > 0$ is attained at $s_1 \in C_0$ (\emph{Subcase~B}), rotating $O(2)$ on the remaining two coordinates so $r_{s_1} = (A_{s_1}/\sqrt{d_0}, \sqrt{d_1}, 0)$ determines $B_x := M_{s_1, x} - g(s_1, x) - \frac{A_{s_1} A_x}{d_0} \in \mathbb{Q}$ and $\mathrm{rem}_2(x) := \mathrm{rem}_1(x) - B_x^2 / d_1 = q_x^2 d_x \ge 0$, leaving one sign $c_x = \varepsilon_x q_x \sqrt{d_x}$.
  \end{itemize}
  \item By classical $\mathbb{Q}$-linear independence of $\{\sqrt{d} : d \ge 1 \text{ squarefree}\}$, two candidates $x, y$ have rational inner product $M_{xy} \in \mathbb{Q}$ if and only if $d_x = 0$, $d_y = 0$, or $d_x = d_y$, and every orthogonality equation $\sum_k x_k u_k (\varepsilon_k q_k \sqrt{d_k}) = 0$, $\sum_k x_k A_k (\varepsilon_k q_k \sqrt{d_k}) = 0$ decomposes into exact $\mathbb{Q}$-linear equations per squarefree radical $d$. Furthermore, in $O(3)$, any $24$-element choice of $C_0$ leaves a rational $14 \times 14$ matrix
  \[
    \tilde{M}_{14} := \diag(W_\Pi, 20 d_0, 20 d_1) - \sum_{x \in C_0} F_x F_x^T = \sum_{y \in B_2 \cup B_3} F_y F_y^T \succeq 0, \qquad F_x := (u_x, A_x, B_x)^T \in \mathbb{Q}^{14},
  \]
  certified in exact $\mathbb{Q}$ via pivoted $LDL^T$ (emitting a rational negative-curvature vector $w^T \tilde{M}_{14} w < 0$ when indefinite) and exact $\mathbb{Q}$ RREF nullspace pruning ($w \in \ker \tilde{M}_{14} \implies w^T F_y = 0$ for all $y \in B_2 \cup B_3$).
\end{enumerate}
\end{theorem}

\begin{proof}
Since $\Pi$ selects $r = \rk(M_{KK})$ linearly independent rows of $M_{KK} = U_r \Lambda U_r^T$, $\tilde{T} = \Pi_{\mathrm{row}} U_r \Lambda^{1/2} \in \mathbb{R}^{r \times r}$ is invertible with $\tilde{T}\tilde{T}^T = \Pi_{\mathrm{row}} M_{KK} \Pi_{\mathrm{row}}^T$. For $v_x = M_{KK}\alpha_x \in \mathrm{col}(M_{KK})$, $Z_K^T \alpha_x = \Lambda^{1/2} U_r^T \alpha_x$ must equal $p_x$ (as $Z_K Z_x^T = v_x$), so $\tilde{T} p_x = \Pi_{\mathrm{row}} U_r \Lambda U_r^T \alpha_x = v_x[\Pi] = u_x$ and $\|p_x\|_2^2 = \alpha_x^T M_{KK} \alpha_x = v_x^T M_{KK}^+ v_x = Q(v_x)$. Applying $\tilde{T} \otimes \tilde{T}$ to $\sum_{x \notin K} p_x p_x^T = 20 I_r - \Lambda$ gives $\sum_{x \notin K} u_x u_x^T = \Pi_{\mathrm{row}}(20 M_{KK} - M_{KK}^2)\Pi_{\mathrm{row}}^T = W_\Pi$, establishing~(1). Parts~(2) and~(3) follow from Gram--Schmidt orthogonalization of $(r_{s_0}, r_{s_1})$ and Besicovitch's linear independence of squarefree radicals over $\mathbb{Q}$.
\end{proof}

Two structural lemmas close subtle corner cases in $S_1$:

\begin{lemma}[$q_0 = 4$ Subcase-A collapse on $S_1(4,4,4,4)$]\label{lem:alg88}
On $S_1(4, 4, 4, 4)$, $M_{KK} = \bigoplus_{c=1}^4 M_4$ with $M_4 = 2(2I_4 - A_{C_4})$, so balanced vectors $v \in \{\pm 1\}^{16}$ ($\mathbf{1}_{C_4}^T v_c = 0$) take values $Q(v) = \sum_{c=1}^4 Q_c(v_c) \in \{2, \frac{5}{2}, 3, \frac{7}{2}, 4\} < \frac{9}{2}$. Whenever $\max_{x \in C_0} \mathrm{res}(x) = \frac{1}{2}$ (i.e.\ $q_0 := Q(s_0) = 4$), every $x \in C_0$ has $Q(x) = 4$, $g(s_0, x) \in \mathbb{Z}$, $A_x \in \{ \pm \frac{1}{2} \}$, and $\mathrm{rem}_1(x) = 0$ (placing $q_0 = 4$ strictly in Subcase~A, $\rho_{C_0} = 0$).
\end{lemma}

\begin{proof}
On a single $4$-cycle, the six balanced vectors in $\{\pm 1\}^4$ consist of two alternating vectors $\pm(1,-1,1,-1)$ ($\lambda = 8$ of $M_4$, $\|v_c\|_2^2 = 4$, $Q_c = v_c^T M_4^+ v_c = \frac{4}{8} = \frac{1}{2}$) and four adjacent-pair vectors $\pm(1,1,-1,-1), \pm(1,-1,-1,1)$ ($\lambda = 4$ of $M_4$, $Q_c = \frac{4}{4} = 1$). On the $\lambda = 4$ eigenspace, inner products $v_c^T M_4^+ w_c = \frac{1}{4} v_c^T w_c \in \{-1, 0, +1\} \subset \mathbb{Z}$. If $\max_{x \in C_0} \mathrm{res}(x) = \frac{1}{2}$, then every $x \in C_0$ satisfies $\frac{9}{2} - Q(x) \le \frac{1}{2} \implies Q(x) \ge 4 \implies Q(x) = 4$, forcing $Q_c(x_c) = 1$ on all four $4$-cycles and $g(s_0, x) = \sum_{c=1}^4 (s_0)_c^T M_4^+ x_c \in \mathbb{Z}$. Since $M_{s_0, x} \in \{\frac{1}{2}, -\frac{3}{2}\} \subset \mathbb{Z} + \frac{1}{2}$, we have $A_x = M_{s_0, x} - g(s_0, x) \in \mathbb{Z} + \frac{1}{2}$, so $A_x^2 / (\frac{1}{2}) = 2A_x^2 \ge \frac{1}{2}$, whence $0 \le \mathrm{rem}_1(x) = \frac{1}{2} - 2A_x^2 \le 0 \implies A_x = \pm \frac{1}{2}$ and $\mathrm{rem}_1(x) = 0$.
\end{proof}

\begin{lemma}[Identical-coordinate exclusion and same-profile $\le 2$-copy cap]\label{lem:alg810}
In every full-core branch of $S_1, S_2, S_3$:
\begin{enumerate}
  \item \emph{(No identical full coordinates)} No two distinct vertices $x \neq y$ in $C_0$ (nor in $B_2 \cup B_3$) satisfy $z_x = z_y$; hence $\{0, 1\}$ indicator variables over distinct full coordinate tuples $(v_x, r_x)$ are exact without multiplicity loss.
  \item \emph{(Exact twin threshold and $\le 2$-copy cap on $C_0$)} For $x \neq y$ in $C_0$ sharing core projection $v_x = v_y = v$ with residual $d := \frac{9}{2} - Q(v)$,
  \[
    \langle r_x, r_y \rangle = M_{xy} - \Big(\frac{9}{2} - d\Big) = \begin{cases} d - 4 & \text{if } S_{xy} = +1\ (M_{xy} = \frac{1}{2}), \\ d - 6 & \text{if } S_{xy} = -1\ (M_{xy} = -\frac{3}{2}). \end{cases}
  \]
  Across all $7$ even partitions of $S_1$ (and all full-core $C_0$ branches of $S_2, S_3$), $\min_v Q(v) = 2$, so $d \le \frac{5}{2} < \frac{8}{3} < 3$. Consequently: (a) $S_{xy} = -1$ is impossible ($|d - 6| > d$ for $d < 3$); (b) $S_{xy} = +1$ requires $|d - 4| \le d \iff d \ge 2$ (so zero same-profile twins exist whenever $d < 2$, in particular at $q_0 \in \{4, \frac{7}{2}, 3\}$ on $S_1(4,4,4,4)$); and (c) if three distinct vertices $x_1, x_2, x_3 \in C_0$ shared $v$, then
  \[
    0 \le \|r_{x_1} + r_{x_2} + r_{x_3}\|_2^2 = 3d + 6(d - 4) = 9d - 24 \le 9\Big(\frac{5}{2}\Big) - 24 = -\frac{3}{2} < 0,
  \]
  contradiction. Thus at most $2$ vertices in $C_0$ can share any core vector $v$ (and similarly on $B_2 \cup B_3$ where $M_{yy} = 4$ and $\max \mathrm{res}_B \le 2$ caps multiplicity at $\lfloor 4/(4 - \mathrm{res}_B) \rfloor \le 2$). Moreover, for any twin $y \neq s_0$ of $s_0$, $A_y = d_0 - 4 \neq d_0 = A_{s_0}$ ($\mathrm{rem}_1(y) = 8 - 16/d_0 \ge 0 \iff d_0 \ge 2$), and for any twin $y \neq s_1$ of $s_1$ in $O(3)$, $B_y = d_1 - 4 \neq d_1 = B_{s_1}$ ($\mathrm{rem}_2(y) = 8 - 16/d_1 \ge 0 \iff d_1 \ge 2$), so retaining candidate profiles equal to the anchor projection of $s_0$ (and of $s_1$ in $O(3)$) covers every same-profile twin without omission.
\end{enumerate}
\end{lemma}

\begin{proof}
If $z_x = z_y$ ($x \neq y$), then $M_{xx} = \|z_x\|_2^2 = \langle z_x, z_y \rangle = M_{xy}$, contradicting $\frac{9}{2} \notin \{\frac{1}{2}, -\frac{3}{2}\}$ on $C_0$ and $4 \notin \{-2, 0, \pm 1\}$ on $B_2 \cup B_3$. All assertions in~(2) follow directly from Cauchy--Schwarz and $\sum_{i,j=1}^3 \langle r_{x_i}, r_{x_j} \rangle \ge 0$.
\end{proof}

Finally, cumulative $\operatorname{Stab}$-orbit ordering (by descending $\mathrm{res}(s_0)$, then descending $\mathrm{rem}_1(s_1)$, and cumulative $\operatorname{Stab}(\text{base})$ orbits on the mandatory $N_{-3/2}(s_0)$ neighborhood where every $x \in C_0$ has $8$ negative neighbors) together with global reflection $b \mapsto -b$ ($O(2)$) or $c \mapsto -c$ ($O(3)$) preserves every realization without loss since $\operatorname{Aut}(M_{KK})$ commutes with $M_{KK}^+$.

\subsection{Complete exact rational elimination across \texorpdfstring{$S_3, S_2, S_1$}{S3, S2, S1}}
Applying Theorem~\ref{thm:alg86_87} and Lemmas~\ref{lem:alg88}--\ref{lem:alg810} in exact rational arithmetic (see~\S\ref{sec:compute} for implementation and verification details) rules out all three templates $S_3, S_2, S_1$:
\begin{enumerate}
  \item \emph{Template $S_3 = (26, 42, 6, 0)$ ($734.26\text{s}$ combined with $S_2$):} On $C_2$ ($|C_2| = 6$, $(E_{11})_{uu} = \frac{1}{12}$, with $2$-vertex Jacobi on $(2,2)$ restricting off-diagonals), $C_2$ splits up to $S_3$-symmetry over $(B_1, B_2, B_3)$ into $10$ symmetry orbits with external negative-neighbor count in $\{0, 2, 4\}$. Extending to the $20 \times 20$ core $K = B_1 \cup C_2$ ($r \in \{13, 14, 15\}$), every core configuration across all $10$ orbits is eliminated via exact $\mathbb{Q}$ $LDL^T$ indefiniteness witnesses, iterative degree peeling ($|C_0| < 26$), $O(1)/O(2)$ anchor LP Farkas duals, and on the tenth orbit (pair~$1$, $r = 14$) an exact $\mathrm{lcm}$-cleared integer CP-SAT refutation proving $\min(A_1, A_2) \ge \frac{7}{3}$.
  \item \emph{Template $S_2 = (25, 45, 3, 1)$:} The single $C_5$-vertex decouples rationally. For sign $s = +1$ on $C_2$ ($|C_2| = 3$), vertices in $B_1$ have degree in $\{0, 2\}$ with zero cycles (via $20E_{15}[v, y_1] = -\frac{1}{3}$), and the all-degree-$0$ core has Schur complement $-\frac{215}{12} < 0$ ($16 \times 16$ rational $LDL^T$ witness). For $s = -1$, only $z = (1, 1, 1)$ is admissible on $C_2$, yielding $15$ path-cycle core configurations (from $(4,4,6)$ down to the empty cycle set) on $K = B_1 \cup C_2$ ($|K| = 18$, $r \in \{13, 14, 15\}$), all $15$ of which are eliminated ($0$ remaining configurations).
  \item \emph{Template $S_1 = (24, 48, 0, 2)$, $t_1 \le 2$ cycles ($168.49\text{s}$):}
  \begin{itemize}
    \item $(16)$ ($r = 15$, $s_{\mathrm{rem}} = 0$): the sum of the $24$ smallest residuals on $C_0$ is $12 > 0$, impossible since $s_{\mathrm{rem}} = 0$ requires $\mathrm{res}(x) = 0$;
    \item $(10, 6)$ ($r = 14$): iterative degree-$8$ core pruning empties the candidate pool;
    \item $(8, 8)$ ($r = 14$): all $8$ anchor cases ($7$ high-residual orbits together with the $\mathrm{res} \le \frac{1}{2}$ tail) are infeasible;
    \item $(12, 4)$ ($r = 14$): the zero-residual subcase ($8.39\text{s}$) and all $22$ positive-residual anchor orbits are infeasible.
  \end{itemize}
  \item \emph{Template $S_1$, $(6, 6, 4)$ ($57.91\text{s}$):} $r = 13$ ($O(2)$). The zero-residual subcase ($96$ vectors, $0.0085\text{s}$) and all $7$ positive-residual anchor orbits are ruled out via $53$ degree-peeling eliminations, $301$ certified rational Farkas duals ($\min\text{gap} = \frac{32852784940321648105016}{24709849076520975126375} > 0$), and $8$ zero-slack $\mathrm{lcm}$-cleared integer CP-SAT refutations ($0$ remaining configurations).
  \item \emph{Template $S_1$, $(8, 4, 4)$ ($80.42\text{s}$):} $r = 13$ ($O(2)$). The zero-residual subcase ($256$ vectors, $0.1107\text{s}$) and all $14$ positive-residual anchor orbits are ruled out via $672$ degree-peeling eliminations, $671$ certified rational Farkas duals ($\min\text{gap} = \frac{19}{6} > 0$), and $4$ zero-slack integer CP-SAT refutations ($0$ remaining configurations).
  \item \emph{Template $S_1$, $(4, 4, 4, 4)$:} $r = 12$ ($O(3)$). Lemma~\ref{lem:alg88} places $q_0 = 4$ inside Subcase~A, and Subcase~A across $q_0 \in \{4, \frac{7}{2}, 3, \frac{5}{2}, 2\}$ is ruled out in $0.13\text{s}$ ($0.12\text{s}$ at $q_0 = 4$, $0.01\text{s}$ at $q_0 = \frac{7}{2}$, and $<0.005\text{s}$ elsewhere). Subcase~B spans $135$ $(s_0, s_1)$ $O(3)$-orbits:
  \begin{itemize}
    \item $q_0 = \frac{7}{2}$ ($d_0 = 1$, $19$ orbits, $28.53\text{s}$): $285$ degree-peeling + $412$ Farkas + $2$ complete $24$-vertex $C_0$-configurations, all refuted ($0$ remaining);
    \item $q_0 = 5/2$ ($d_0 = 2$, twins active, $48$ orbits, $97.99\text{s}$): $1{,}559$ degree-peeling + $708$ Farkas ($0$ remaining);
    \item $q_0 = 2$ ($d_0 = \frac{5}{2}$, twins active including twin $s_1$-orbit, $23$ orbits, $119.84\text{s}$): $430$ degree-peeling + $276$ Farkas ($0$ remaining);
    \item $q_0 = 3$ ($d_0 = \frac{3}{2} < 2$, $0$ twins by Lemma~\ref{lem:alg810}(2), $45$ orbits, $3105.67\text{s}$): $21{,}543$ degree-peeling + $18{,}723$ $C_0$-branch Farkas + $1{,}985$ complete $24$-vertex $C_0$-configurations ($34{,}376$ leaves reduced via exact $14 \times 14$ $\mathbb{Q}$ $LDL^T$/RREF to $33{,}336$ $(B_2 \cup B_3)$ Farkas duals and $1{,}040$ zero-slack integer CP-SAT refutations; $0$ remaining).
  \end{itemize}
  Across $(4, 4, 4, 4)$, all $53{,}455$ Farkas duals satisfy $\text{gap} > 0$ in exact $\mathbb{Q}$, giving $54{,}427$ certified Farkas duals across the three $O(2)/O(3)$ partitions $(6,6,4), (8,4,4), (4,4,4,4)$ and $1{,}052$ zero-slack CP-SAT refutations ($8 + 4 + 1{,}040$) across $(6,6,4)$, $(8,4,4)$, and complete $24$-vertex $C_0$-configurations at $q_0 = 3$ (alongside separate exact CP-SAT refutations in $(8,8)$, $(12,4)$, Subcase~A, $S_2$, and $S_3$).
\end{enumerate}

Combining Theorem~\ref{thm:srg136}, Propositions~\ref{prop:kill47} and~\ref{prop:other_19_74}, Lemma~\ref{lem:idx13_74to3}, and the exact rational elimination of $S_3, S_2, S_1$ proves Theorem~\ref{thm:main19}; combining Propositions~\ref{prop:kill861}--\ref{prop:kill864} with the exact certificates for the remaining $(18, 59)$ spectra proves Theorem~\ref{thm:main18}.

% ============================================================================
\section{Computational verification, constructive searches, and independent implementation}\label{sec:compute}
% ============================================================================

\subsection{Characteristic polynomial enumeration and exhaustive \texorpdfstring{$\mathcal{P}_{n,7}$}{P_(n,7)} saturation}
Our parallel C++20 polynomial enumerator \nolinkurl{compute/poly/enum_polys_v4.cpp} generates every monic cofactor $\phi(x) = Q(x)/(x - \kappa)^{m_\kappa}$ ($\deg\phi = \eta \le d - m_\kappa$) compatible with Theorem~\ref{thm:twoadic_val}, Theorem~\ref{thm:pfaffian}, real-rootedness on $[-5, \infty)$, and shifted power-sum moments, using $512$-bit stack-allocated signed \texttt{BigInt} arithmetic (\texttt{int64\_t}/\texttt{\_\_int128} fast path) and zero IEEE-754 operations:
\begin{enumerate}
  \item \emph{Top-down prefix tree with exact Cauchy--Schwarz and AM--GM bounds:} Working in the shifted variable $y = z + 2$ ($\lambda = 2y - 5 \ge -5$, so $y_i \ge 0$), top-down coefficient steps $b_1, \dots, b_\eta$ are bounded at depth $k$ by the $(\eta - k)$-th Rolle derivative $g^{(k)}(y)$:
  \begin{itemize}
    \item Cauchy--Schwarz on remaining roots bounds $b_k$ inside the closed integer interval $[\lceil A_{\mathrm{num}}/8 \rceil, \lfloor B_{\mathrm{num}}/8 \rfloor]$ (cleared of divisions).
    \item Symmetric Bareiss $LDL^T$ enforces exact integer positive semidefiniteness of the four Hamburger and Stieltjes Hankel moment matrices $H(s_0, \dots, s_{2m})$ and $H(s_1, \dots, s_{2m+1})$ available at depth $k$.
    \item At full depth $k = \eta$, exact integer Sturm chains and Rolle derivative sign tables certify $\eta$ real roots in $[0, \infty)$ ($\lambda_i \ge -5$), explicitly retaining $y = 0$ ($\lambda = -5$, i.e.\ $\rk(5I_n + S) < d$). Notice that for odd $n$ at $k = \eta$, $b_\eta$ has $2$-adic step $2^{\eta - 1}$ (Weakly Type~2 on $\phi$), while full-degree parity $2^{2r} \mid e_{2r}$ on $Q$ (Theorem~\ref{thm:twoadic_val}(3)) is enforced exactly.
  \end{itemize}
  \item \emph{Corollary~\ref{cor:cor43_full} mod-$4$ prefix filter and $\mathcal{P}_{n,7}$ mod-$128$ lookup:} For odd $n$, a precomputed bitmask over the $(d+1)/2$ choices of $Q_0(z) \bmod 2$ prunes invalid prefixes modulo $2^{n+1}$ at every depth $k \ge 3$. At the leaves, each polynomial is checked against the Burnside walk-orbit congruence family $\mathcal{P}_{n,7} \subseteq (\mathbb{Z}/128\mathbb{Z})[z]$:
  \begin{lemma}[Exact $\mathcal{P}_{n,7}$ saturation]\label{lem:p7_saturation}
  By~\cite[Cor.~3.3 and Cor.~3.13]{GY2019}, $|\mathcal{P}_{n,e}| \le 2^{\binom{e-2}{2} + (n \bmod 2)}$ for all $e \ge 2$; at $e = 7$, $\binom{5}{2} = 10$, giving theoretical upper bounds $|\mathcal{P}_{n,7}| \le 1{,}024$ for even $n$ and $|\mathcal{P}_{n,7}| \le 2{,}048$ for odd $n$. Sampling random $n \times n$ Seidel matrices modulo $128$ collects $1{,}024$ distinct residues for $n \in \{48, 58, 60, 72, 74\}$ and $2{,}048$ distinct residues for $n \in \{49, 59, 73, 75\}$. Because a subset of a finite set of cardinality $\le M$ having $M$ observed elements equals the entire set, the sampled lookup tables equal $\mathcal{P}_{n,7}$ with $100\%$ mathematical completeness. Moreover, \nolinkurl{stage2_deck_par.cpp} (\texttt{!= 2048}) and \nolinkurl{stage2_deck_odd_v3.cpp} (\texttt{!= 1024}) abort immediately if $\mathcal{P}_{n-1,7}$ fails to saturate, while \nolinkurl{enum_polys_v4.cpp} sets \texttt{enabled = false} (conservatively retaining all candidates) if unsaturated.
  \end{lemma}
\end{enumerate}

Table~\ref{tab:stage1_runtimes} records exact enumeration counts and runtimes across all five resolved cases, together with prefix-depth benchmarks illustrating the combinatorial growth on the open single-step boundary pairs $(18, 58)$ and $(19, 73)$ (where $\mathcal{T} = 84$ and $\mathcal{T} = 81$ give $\theta = 22 > d$, yielding only the trivial bound $m_\kappa \ge 0$ with no guaranteed root at $\kappa$).

\begin{table}[htbp]
\centering
\caption{Characteristic polynomial enumeration counts, filtration stages, and prefix-depth benchmarks (\nolinkurl{enum_polys_v4.cpp}, 8 threads).}\label{tab:stage1_runtimes}
\scriptsize
\setlength{\tabcolsep}{4pt}
\begin{tabular}{@{}lccccrrrl@{}}
\toprule
Case $(d,n)$ & $\deg\phi$ & $\mathcal{T}$ & $|\mathcal{P}_{n,7}|$ & Raw Type-2 & Post-$\mathcal{P}_{n,7}$ & Post-Thm.~\ref{thm:pfaffian} & Runtime & Status \\
\midrule
$(19, 75)$ & $7$ & $25$ & $2{,}048$ & $14$ & $8$ & $8$ & $0.007\text{s}^{\text{b}}$ & Matches~\cite{GSY2021} \\
$(18, 60)$ & $12$ & $48$ & $1{,}024$ & $44$ & $44$ & $44$ & $0.15\text{s}^{\text{b}}$ & Matches~\cite{GS2024} \\
$(17, 49)$ & $13$ & $49$ & $2{,}048$ & $704^{\text{a}}$ & $188^{\text{a}}$ & $109$ & $2.27\text{s}^{\text{b}}$ & Matches~\cite{GSY2023} \\
$\mathbf{(19, 74)}$ & $\mathbf{13}$ & $\mathbf{52}$ & $\mathbf{1{,}024}$ & $\mathbf{169}$ & $\mathbf{169}$ & $\mathbf{169}$ & $\mathbf{1.35\text{s}}^{\text{b}}$ & \textbf{Complete enumeration} \\
$\mathbf{(18, 59)}$ & $\mathbf{17}$ & $\mathbf{65}$ & $\mathbf{2{,}048}$ & --- & --- & $\mathbf{1{,}214}$ & $\mathbf{393.11\text{s}}$ & \textbf{Complete enumeration} \\
\midrule
$(18, 58)$ & $18$ & $84$ & $1{,}024$ & \multicolumn{3}{c}{$17{,}481{,}969$ prefixes at depth $k = 7 / 18$} & $100.41\text{s}$ & Open boundary \\
$(19, 73)$ & $19$ & $81$ & $2{,}048$ & \multicolumn{3}{c}{$10{,}682{,}723$ prefixes at depth $k = 8 / 19$} & $486.41\text{s}$ & Open boundary \\
\bottomrule
\end{tabular}
\par\vspace{3pt}
\raggedright\scriptsize
$^{\text{a}}$Weakly Type-2 tree search without $2^{2r}\mid e_{2r}$ yields $2{,}291$ unfiltered and $194$ post-$\mathcal{P}_{49,7}$ (\cite[Lem.~3.6]{GSY2023}); enforcing exact full-degree parity gives $704 \to 188 \to 109$.\\
$^{\text{b}}$Archived wall-clock \texttt{elapsed\_s} in \nolinkurl{certificates/stage1/} (\nolinkurl{raw_19_75.json}: $0.0068\text{s}$; \nolinkurl{raw_18_60.json}: $0.1471\text{s}$ including $\mathcal{P}_{60,7}$ table setup, pure tree $0.0172\text{s}$; \nolinkurl{cor43_17_49.json}: $2.2661\text{s}$; \nolinkurl{raw_19_74.json}: $1.3537\text{s}$ including $\mathcal{P}_{74,7}$ table setup, pure tree $0.1419\text{s}$).
\end{table}

\subsection{Interlacing deck generation and exact rational certification}
For each admissible polynomial $\chi_S = \mathrm{Gcd}_x \cdot H_x$ ($\deg H_x = s$), \nolinkurl{stage2_deck_par.cpp} (even $n \to$ odd $n-1$) and \nolinkurl{stage2_deck_odd_v3.cpp} (odd $n \to$ even $n-1$) enumerate the exact vertex-deletion interlacing deck $\Deck(\chi_S)$ using 36-bit integer-bracketed roots of $H_x$ and exact divided-difference interval bounds $[L_k, U_k]$:
\begin{itemize}
  \item On $(19, 74)$ ($169$ candidates, $s \in [4, 9]$), \nolinkurl{stage2_deck_par.cpp} generates $5{,}069$ deck polynomials in $98.86\text{s}$ ($2{,}744$ of which fail Theorem~\ref{thm:pfaffian} at odd order $73$, ruling out $8$ candidates immediately via \texttt{t42\_farkas}).
  \item On $(18, 59)$ ($1{,}214$ candidates, $s \in [4, 11]$), \nolinkurl{stage2_deck_odd_v3.cpp} generates $781{,}998$ Strict Type-2 ($2^{58}$-divisible) deck polynomials in $21.59\text{s}$.
\end{itemize}

All one-vertex, Eulerian/Galois, and higher-order certification programs (\nolinkurl{phase2_even.py}, \nolinkurl{phase2_odd_v2.py}, \nolinkurl{run_stage3_18_59.py}, \nolinkurl{run_stage3_19_74.py}, \nolinkurl{run_705_cone0_fast.py}, \nolinkurl{enum705.cpp}, and \nolinkurl{compute/poly/idx13_exact/}) enforce four strict exactness safeguards:
\begin{enumerate}
  \item \emph{Certified integer Farkas duals with exact Bland Simplex fallback:} Floating-point LP (\texttt{HiGHS}) is used solely as an untrusted dual hint; if rounding fails $c^T A \ge 0, c^T b < 0$ in exact \texttt{Fraction}, an exact rational Phase-I Simplex solver with Bland's anti-cycling rule computes a certified rational Farkas vector cleared to exact integers. Whenever an LP variable-bound query $x_j \le \lfloor u_j \rfloor$ lacks an exact dual certificate, its bound falls back conservatively to $[0, n]$.
  \item \emph{Exact $K^\times/(K^\times)^2$ square-class invariants and algebraic root comparison:} Rather than hashing via heuristic prime evaluations (which can fail to separate distinct square classes over degree-$e \ge 3$ fields, e.g.\ on $q_{666}$ and $q_{778}$ of $(18, 59)$), square-class equivalence $a \cdot b \in (K^\times)^2$ in $K = \mathbb{Q}[x]/(q(x))$ is decided algebraically via square-root factoring $\mathrm{factor}(t^2 - a(\bar{\lambda})b(\bar{\lambda}))$ over $K$ pre-bucketed by Dedekind split-prime Legendre signatures, while $2$-vertex Cauchy interlacing is verified via exact \texttt{CRootOf} algebraic root isolation.
  \item \emph{Strict odd-$n$ domain guard on Theorem~\ref{thm:eulerian}:} On odd-to-even vertex deletion ($59 \to 58$), even-order $M_{uu} \bmod 4$ Eulerian congruences are strictly disabled (Remark~\ref{rem:odd_guard}), relying solely on Galois orthogonality, Proposition~\ref{prop:prop61}, $2$-vertex Jacobi interlacing, and coupled one- and two-vertex linear Farkas duals ($\sum_u c_u (x_u - \delta_{uv} - \sum_q m_{v,u,q}) \le 0$). On $q_{705}$ ($s = 7$, $6$ cones), custom C++ exact RREF enumerator \nolinkurl{enum705.cpp} coupled with \nolinkurl{run_705_cone0_fast.py} ($2{,}234.24\text{s}$ combined) exhaustively enumerates all $19{,}975$ nonnegative integer solutions of cone~$0$, refuting $19{,}456$ via two-vertex Farkas duality and $519$ via $S_v$-Galois square classes, while cones $1$--$5$ are eliminated by $25$ coupled duals and Proposition~\ref{prop:prop61}.
  \item \emph{Exact rational verification modules for $p_{13}$ (\nolinkurl{compute/poly/idx13_exact/}):} Static AST verification bans \texttt{eigh}, \texttt{eigvalsh}, \texttt{pinv}, \texttt{rint}, \texttt{linprog}, and float tolerances across the five scripts generating the archived JSON certificates for Section~\ref{sec:idx13}: \nolinkurl{certify_s3_s2.py} $\to$ \nolinkurl{cert_s3_s2.json} ($S_3$ and $S_2$), \nolinkurl{certify_s1_t2.py} $\to$ \nolinkurl{cert_s1_t2.json} ($S_1$ partitions $(16)$, $(10,6)$, $(8,8)$, $(12,4)$), \nolinkurl{certify_s1_6_6_4.py} $\to$ \nolinkurl{s1_6_6_4_certs.json}, \nolinkurl{certify_s1_8_4_4.py} $\to$ \nolinkurl{s1_8_4_4_certs.json}, and \nolinkurl{certify_s1_4_4_4_4.py} $\to$ \nolinkurl{s1_4_4_4_4_certs.json}. (In the public repository release, large JSON certificate files are stored compressed as \texttt{.json.gz}, while \nolinkurl{certificates/verify_all.py} verifies their decompressed byte streams bit-for-bit against \nolinkurl{certificates/SHA256SUMS.txt}.)
\end{enumerate}

\subsection{Calibration on \texorpdfstring{$12$}{12} extremal Seidel matrices and constructive searches}
Every enumeration, interlacing, Eulerian, and Galois routine was calibrated against $12$ reconstructed extremal Seidel matrices archived in \nolinkurl{certificates/calibration/} (Table~\ref{tab:calibrated_12}): each matrix passed exact rational pivoted $LDL^T$ ($\rk(5I_n + S) = d$, $5I_n + S \succeq 0$), satisfied Theorem~\ref{thm:twoadic_val}, Theorem~\ref{thm:pfaffian}, Theorem~\ref{thm:weapon_b}, Theorem~\ref{thm:eulerian}, and Theorem~\ref{thm:galois}, and produced a non-eliminated interlacing deck containing all $n$ true vertex-deleted submatrix polynomials with zero false exclusions.

\begin{table}[htbp]
\centering
\caption{The $12$ reconstructed extremal Seidel matrices in \nolinkurl{certificates/calibration/} (all pass $\rk(5I_n+S)=d$ with zero false exclusions).}\label{tab:calibrated_12}
\resizebox{\textwidth}{!}{%
\begin{tabular}{@{}llcl@{}}
\toprule
Archive File & Source & $(d,n)$ & Exact Characteristic Polynomial $\chi_S(x)$ \\
\midrule
\nolinkurl{cal_17_48_netto.json} & Netto~\cite{GKMS2016} & $(17,48)$ & $(x+5)^{31}(x-11)^9(x-7)^8$ \\
\nolinkurl{cal_17_48_a5_3_a1_4.json} & MSY26~\cite{MSY2026} & $(17,48)$ & $(x+5)^{31}(x-11)^9(x-7)^8$ \\
\nolinkurl{cal_18_57_gsy21_f1.json} & GSY23 $F_1$~\cite{GSY2023} & $(18,57)$ & $(x+5)^{39}(x-15)(x-13)^4(x-11)^9(x-9)^2(x-7)(x-4)$ \\
\nolinkurl{cal_18_57_gsy21_f2.json} & GSY23 $F_2$~\cite{GSY2023} & $(18,57)$ & $(x+5)^{39}(x-13)^6(x-11)^5(x-9)(x-7)^4(x^2-25x+152)$ \\
\nolinkurl{cal_18_57_gsy21_f3.json} & GSY23 $F_3$~\cite{GSY2023} & $(18,57)$ & $(x+5)^{39}(x-15)(x-11)^{11}(x^2-20x+87)^2(x^2-19x+76)$ \\
\nolinkurl{cal_18_57_gsy21_f4.json} & GSY23 $F_4$~\cite{GSY2023}$^{\text{a}}$ & $(18,57)$ & $(x+5)^{39}(x-11)^{13}(x^2-20x+87)(x^3-32x^2+303x-716)$ \\
\nolinkurl{cal_18_57_a9_2_a1.json} & LMTY25~\cite{LMTY2025}$^{\text{b}}$ & $(18,57)$ & $(x+5)^{39}(x-15)(x-13)^2(x-11)^7(x-9)(x-7)(x^2-22x+113)r_4(x)$ \\
\nolinkurl{cal_18_57_a9_2_a1_6ev_first.json} & LMTY25~\cite{LMTY2025}$^{\text{b}}$ & $(18,57)$ & $(x+5)^{39}(x-13)^3(x-11)^{11}(x-9)^2(x^2-17x+36)$ \\
\nolinkurl{cal_18_57_a9_2_a1_6ev_last.json} & LMTY25~\cite{LMTY2025}$^{\text{b}}$ & $(18,57)$ & $(x+5)^{39}(x-13)^2(x-11)^{13}(x-9)(x^2-17x+32)$ \\
\nolinkurl{cal_18_57_a5_3_a1_4_spec1_111.json} & MSY26~\cite{MSY2026} & $(18,57)$ & $(x+5)^{39}(x-13)^8(x-9)^8(x^2-19x+72)$ \\
\nolinkurl{cal_18_57_a5_3_a1_4_spec2_11m1.json} & MSY26~\cite{MSY2026} & $(18,57)$ & $(x+5)^{39}(x-13)^6(x-12)(x-11)(x-9)^6(x^2-20x+87)^2$ \\
\nolinkurl{cal_19_72_asche.json} & Asche~\cite{Taylor1977,LinYu2020} & $(19,72)$ & $(x+5)^{53}(x-19)^3(x-13)^{16}$ \\
\bottomrule
\end{tabular}%
}
\par\vspace{3pt}
\raggedright\scriptsize
$^{\text{a}}$Corrects printed typo in~\cite[p.~1876]{GSY2023} (\S\ref{subsec:lit_review}). \quad $^{\text{b}}$Converted via $S_{\mathrm{std}} = -S_{\mathrm{LMTY}}$; here $r_4(x) = x^4 - 39x^3 + 551x^2 - 3329x + 7264$.
\end{table}

\begin{remark}[Constructive and local-neighborhood searches on $(18, 58)$ and $(19, 73)$ and explicit non-proof disclaimer]\label{rem:trackb_disclaimer}
Alongside our spectral upper-bound proofs, we carried out constructive and local-neighborhood searches for $(18, 58)$ and $(19, 73)$ (\nolinkurl{compute/search/}):
\begin{enumerate}
  \item \emph{Exact $1$-line and $k$-opt extensions:} Exact \texttt{\_\_int128} column-space extension (\nolinkurl{extender.cpp}) tested $10$ calibrated systems ($2$ in $(17,48)$, $7$ in $(18,57)$, and Asche's $(19,72)$ set among the $12$ archived \nolinkurl{cal_*.json} matrices; $0$ same-dimension extensions), all $768$ non-strongly-maximal MSY26 matrices across all $12$ Latin-square classes ($0$ same-dimension extensions; exact $\mathbb{R}^{19}$ max-clique lifts reach at most $58 + 11 = 69 < 72$ lines), and $120$ sampled LMTY25 $A_9^2 \oplus A_1$ constructions ($0$ extensions). Deleting $k \in \{1, 2, 3, 4, 5\}$ lines and running exact \texttt{\_\_int128} Bareiss + Bron--Kerbosch maximum-clique re-extension (\nolinkurl{kopt_search.cpp}) across $65{,}505$ jobs in $d \in \{18, 19\}$ ($513$ exhaustive $k=1$, $14{,}364$ exhaustive $k=2$, and $42{,}000$ random $k \in \{3,4,5\}$ jobs across the $9$ calibrated $(18,57)$ matrices; $72 + 2{,}556 + 6{,}000 = 8{,}628$ jobs on Asche's $(19,72)$ set) yielded $0$ improvements beyond $57$ and $72$.
  \item \emph{Lattice, $\mathbb{Z}_p$-Fourier, spectral/Gram, and CP-SAT searches:} Integral root-overlattice searches (\nolinkurl{overlattice_search.py}) tested calibrated $A_9^2 \oplus A_1$, $A_5^3 \oplus A_1^4$ (max clique $57$) and rank-$19$ glue lattices $A_7^2 \oplus A_3 \oplus A_1^2$, $D_6^3 \oplus A_1$, $A_3^6 \oplus A_1$ in $d = 18$ (max clique $54$), and rank-$20$ glue lattices $A_5^3 \oplus A_1^5$, $A_9^2 \oplus A_1^2$, $A_7^2 \oplus A_3 \oplus A_1^3$, $D_6^3 \oplus A_1^2$, $A_5^3 \oplus A_3 \oplus A_1^2$ in $d = 19$ (max clique $69$); direct $\mathbb{Z}^d$ norm-$10$ clique chains (\nolinkurl{z_lattice_search.cpp}, reaching $55$ in $\mathbb{R}^{18}$ and $64$ in $\mathbb{R}^{19}$); exact $\mathbb{Z}_p$-equivariant rank-$1$ Hermitian Fourier block searches (\nolinkurl{symmetric_seidel_search.py}) at primes $p \in \{5, 7, 11, 13, 19\}$ for $(18,58)$ and $p \in \{5, 7, 11, 13\}$ for $(19,73)$ ($0$ valid completions); spectral-energy tabu (\nolinkurl{spectral_tabu.cpp}) and continuous Gram L-BFGS + rounding (\nolinkurl{run_gram_frontier.py}, best rounded spectral excess $E = 0.588171$ on $(18,58)$ and $E = 0.826110$ on $(19,73)$); and OR-Tools CP-SAT large-neighborhood destroys ($k \in \{12, 17, 22, 27\}$ across $F_1$--$F_4$, \nolinkurl{sat_search.py}, max recovered $57$) all terminated at $\le 57$ and $\le 72$.
\end{enumerate}
We emphasize explicitly that these computations are \textbf{bounded constructive and local-neighborhood searches, not exhaustive nonexistence proofs}: whether $N(18)$ equals $57$ or $58$, and whether $N(19)$ equals $72$ or $73$, remain open problems.
\end{remark}

\subsection{Independent verification suite (\texttt{audit/})}
To guard against implementation bugs, shared helper assumptions, or serialization errors, a separate verification suite (\nolinkurl{audit/}) was implemented independently in C++20 and Python, sharing zero source files or imports with \nolinkurl{compute/}. Table~\ref{tab:clean_room_audit} summarizes every independent verification check.

\begin{table}[htbp]
\centering
\caption{Independent verification suite (\nolinkurl{audit/}, zero shared code with \nolinkurl{compute/}).}\label{tab:clean_room_audit}
\resizebox{\textwidth}{!}{%
\begin{tabular}{@{}llrl@{}}
\toprule
Verification Script & Verified Scope & Runtime & Outcome \\
\midrule
\nolinkurl{small_n_sanity.py} & $1{,}234$ Seidel matrices ($n \le 16$, exhaustive $n \le 6$) & $32.6\text{s}$ & $100\%$ Thm.~\ref{thm:twoadic_val}, Thm.~\ref{thm:pfaffian}, mult.\ bounds: Verified \\
\nolinkurl{clean_enum.cpp} & Full polynomial tree: $(19,75),(18,60),(17,49),(19,74)$ ($(18,59)$ not re-run) & --- & Exact candidate set \& count match: Verified \\
\nolinkurl{verify_polys.py} & All $10$ candidate tables incl.\ all $1{,}214$ admissible $(18,59)$ polys & --- & Exact Sturm, moments, $2$-adic, Thm.~\ref{thm:pfaffian}: Verified \\
\nolinkurl{verify_gram.py} & All $12$ calibrated matrices (\nolinkurl{certificates/calibration/}) & $1.6\text{s}$ & Exact $\mathbb{Q}$ rank $= d$, PSD, $\chi_S$: Verified \\
\nolinkurl{verify_farkas.py} & Interlacing decks + $1$-vertex Farkas ($(18,60)$ $39$; $(19,74)$ $147$) & $121.5\text{s}$ & Independent Sturm deck + exact duals: Verified \\
\nolinkurl{verify_phase2_certs.py} & All $4$ \nolinkurl{phase2_*.json} bundles ($1{,}432$ certs: $42{+}162{+}100{+}1{,}128$) & $1{,}524.3\text{s}$ & Exact Farkas, Thm.~\ref{thm:pfaffian}, $K^\times$ RREF, Prop.~\ref{prop:prop61}: Verified \\
\nolinkurl{test_weapon_b.py} & Hensel Thm.~\ref{thm:weapon_b}(1)--(2) ($126$ odd + $9$ $(18,57)$ + $24$ rem.\ cands) & $302.4\text{s}$ & $100\%$ mod-$8$ match; $0$ extra exclusions: Verified \\
\nolinkurl{verify_algebraist_sec7.py} & Thms.~\ref{thm:eulerian}--\ref{thm:galois} ($116$ odd + real $(18,57)$ + $46$ Gal.) & $203.7\text{s}$ & Exact gauge uniqueness \& Gal.\ Farkas: Verified \\
\nolinkurl{verify_sec7_certs.py} & All $77/77$ \nolinkurl{sec7_*.json} certificates & $1{,}420.5\text{s}$ & $77/77$ verified ($0$ missed Galois sols): Verified \\
\nolinkurl{verify_phase3_47_861.py} & Independent Steps 1--2 for $p_{47}$ and $q_{861}$ & $248.5\text{s}$ & $928 > 896$ and $840 > 648$: Verified \\
\nolinkurl{verify_phase3_18_59_endgame.py} & All $20/20$ higher-order $(18, 59)$ eliminations & $2{,}041.8\text{s}$ & $20/20$ verified ($86/86$ $(18,59)$ excluded): Verified \\
\nolinkurl{verify_phase3_19_74_idx13.py} & Independent implementation for $(19, 74)$ $p_{13}$ ($S_3,S_2,S_1$) & $830.8\text{s}$ & $74\to 21\to 3\to 0$ ($54{,}427$ duals + twins): Verified \\
\bottomrule
\end{tabular}%
}
\end{table}

% ============================================================================
\section{Lean 4 formalization and explicit hypotheses}\label{sec:lean}
% ============================================================================

Under \nolinkurl{leanprover/lean4:v4.36.0}~\cite{MouraUllrich2021} and Mathlib \texttt{v4.36.0}, we developed a $20{,}765$-line Lean~4 library (\nolinkurl{lean/}, zero \texttt{sorry}, zero \nolinkurl{native_decide}, zero custom axioms) covering the algebraic and spectral theory of Sections~\ref{sec:phase1}--\ref{sec:idx13} together with pure-kernel one-vertex Farkas verification for $(19, 74)$.

\subsection{Module inventory and pure-kernel certificate reflection}
Table~\ref{tab:lean_modules} lists all $20$ Lean source files ($17$ library modules imported in \nolinkurl{EquiangularLines.lean}, \nolinkurl{Audit.lean}, and \nolinkurl{Smoke.lean}). In \nolinkurl{CertKernel.lean}, the checker \nolinkurl{checkCertZ} operates purely on primitive \texttt{Int} lists without importing \texttt{Mathlib.LinearAlgebra}, allowing Lean's logical kernel (\texttt{by decide}) to verify all $147/147$ one-vertex \texttt{1-vertex\_farkas} certificates for $(19, 74)$ across \nolinkurl{Certs19_74A/B/C.lean} using only \texttt{[propext]} (and \texttt{[propext, Quot.sound]} when bridged to real Seidel matrices via \nolinkurl{charpoly_eliminated_of_deck_farkas} in \nolinkurl{Certificates.lean}).

\begin{table}[htbp]
\centering
\caption{Complete module inventory of the \nolinkurl{EquiangularLines} Lean~4 formalization ($20{,}765$ lines, $127$ audited declarations).}\label{tab:lean_modules}
\resizebox{\textwidth}{!}{%
\begin{tabular}{@{}lrrl@{}}
\toprule
Module (\texttt{lean/EquiangularLines/}) & Lines & Audited Decls & Mathematical Contents Formalized \\
\midrule
\nolinkurl{Basic.lean} & $187$ & $7$ & \nolinkurl{IsEquiangularSystem} ($\alpha < 1$), distinctness, $n > d \implies \alpha > 0$ \\
\nolinkurl{Seidel.lean} & $440$ & $17$ & \nolinkurl{IsSeidelMatrix}, $\tr(S)=0$, $\tr(S^2)=n(n-1)$, Thm.~\ref{thm:seidel_gram} ($\rk \le d$) \\
\nolinkurl{Bounds.lean} & $553$ & $5$ & Gerzon bound, relative bound, $(18,58..)$/$(19,73..)$ exclude $1/k$ ($k \ge 7$) \\
\nolinkurl{LemmensSeidel.lean} & $1{,}538$ & $1$ & Pillars + $(K_3,2K_2)$-free Gram proof of $N_{1/3}(d) \le \max(28, 2(d-1))$ \\
\nolinkurl{Spectral.lean} & $2{,}010$ & $25$ & Neumann's theorem, Cor.~\ref{cor:one_fifth}, Prop.~\ref{prop:spectral_fac}, Cauchy interlacing \\
\nolinkurl{TwoAdic.lean} & $2{,}098$ & $3$ & Strict/Weak Type-2 $2$-adic valuations, Pfaffian Thm.~\ref{thm:pfaffian} ($n \ge 1$) \\
\nolinkurl{Multiplicity.lean} & $778$ & $11$ & Shifted Prop.~\ref{prop:gsy_mult}, $m_{15} \ge 6$ on $(19,74)$, $m_{11} \ge 1$ on $(18,59)$ \\
\nolinkurl{Obstructions.lean} & $1{,}075$ & $14$ & $\sum \chi_{S_u} = \chi_S'$, $K^\times/(K^\times)^2$ class, Farkas, Prop.~\ref{prop:prop61} \\
\nolinkurl{GaloisObstructions.lean} & $1{,}113$ & $11$ & Hensel Thm.~\ref{thm:weapon_b}, Thm.~\ref{thm:eulerian} (odd-$n$), Generalized Galois Thm.~\ref{thm:galois} \\
\nolinkurl{SRG136.lean} & $1{,}866$ & $5$ & Bidirectional $p_{136} \Leftrightarrow \SRG(75, 32, 10, 16)$ (Thm.~\ref{thm:srg136}) \\
\nolinkurl{Phase3Algebra.lean} & $2{,}733$ & $13$ & Lem.~\ref{lem:block_comm}, Step~2 of $p_{47}, q_{861}, q_{488}$, Lem.~\ref{lem:resolvent_pole}(1)--(2) \\
\nolinkurl{CertKernel.lean} + \nolinkurl{Certs19_74A/B/C.lean} & $5{,}695$ & $3$ & Pure-\texttt{Int} Farkas kernel + $147/147$ $(19,74)$ \texttt{by decide} shards \\
\nolinkurl{Certificates.lean} & $282$ & $6$ & Polynomial-list bridge, \nolinkurl{charpoly_eliminated_of_deck_farkas} \\
\nolinkurl{Main.lean} + \nolinkurl{EquiangularLines.lean} + \nolinkurl{Audit.lean} + \nolinkurl{Smoke.lean} & $397$ & $6$ & Master theorems + $127$-declaration \texttt{\#print axioms} harness \\
\midrule
\textbf{Total} & $\mathbf{20{,}765}$ & $\mathbf{127}$ & \textbf{Zero \texttt{sorry}, zero \texttt{native\_decide}, zero custom axioms} \\
\bottomrule
\end{tabular}%
}
\end{table}

\subsection{Formalization scope and explicit hypotheses in \texttt{Main.lean}}
To document the formalization boundary unambiguously, we state what is formally proved inside Lean~4 versus what enters the two top-level Lean~4 theorems \nolinkurl{frontier_19_74_no_lines_of_srg_and_cert_exhaustion} and \nolinkurl{frontier_18_59_no_lines_of_cert_exhaustion} in \nolinkurl{EquiangularLines/Main.lean} (whose full $127$-declaration axiom audit appears in Appendix~\ref{app:lean_audit}):
\begin{enumerate}
  \item \emph{Formally proved inside Lean~4:} All preliminary spectral reductions ($\alpha = 1/5$, $\rk(5I_n+S) \le d$, $(x+5)^{n-d}\mid\chi_S$); Theorem~\ref{thm:twoadic_val}, Theorem~\ref{thm:pfaffian}, and multiplicity bounds $m_{15} \ge 6$ and $m_{11} \ge 1$; Jacobi's derivative identity, $K^\times/(K^\times)^2$ square-class invariance, Proposition~\ref{prop:prop61}, Theorem~\ref{thm:weapon_b}, Theorem~\ref{thm:eulerian} (odd-$n$), and Theorem~\ref{thm:galois}; the complete bidirectional equivalence $p_{136} \Leftrightarrow \SRG(75, 32, 10, 16)$ (Theorem~\ref{thm:srg136}); the spectral refutations $\neg(\chi_S = \texttt{idx47Charpoly} \wedge \texttt{RealizesIdx47Block}\,S)$, $\neg\,$\nolinkurl{RealizesIdx861Block}$\,S$, and $\neg\,$\nolinkurl{RealizesIdx488Proj}$\,S$ (proved in \nolinkurl{idx47_eliminated}, \nolinkurl{idx861_eliminated}, and \nolinkurl{idx488_rank_obstruction}), together with resolvent pole cancellation \nolinkurl{idx13_pole_cancellation_at_15}; soundness of integer Farkas deck elimination (\nolinkurl{charpoly_eliminated_of_deck_farkas}); and pure-kernel \texttt{by decide} verification of all $147/147$ one-vertex LP Farkas certificates for $(19, 74)$ (\nolinkurl{certs19_74_all_147_verified}).
  \item \emph{Explicit hypotheses in \nolinkurl{Main.lean} (proved and certified in Sections~\ref{sec:jacobi_galois}--\ref{sec:compute}, not inside Lean):}
  \begin{itemize}
    \item \texttt{hno\_srg : $\neg$\,ExistsSRG 75 32 10 16}, corresponding to Azarija and Marc~\cite{AzarijaMarc2018} (appearing \emph{only} in the $(19, 74)$ theorem, never in $(18, 59)$);
    \item \texttt{hexh}, the finite-enumeration exhaustion disjunction routing any hypothetical realization $S$ into the Lean-refuted cases above (encapsulating C++ polynomial enumeration completeness, Sturm interlacing deck completeness, index linkage between candidate polynomials and certificate records, the remaining $22 = 169 - 147$ $(19, 74)$ candidates not ruled out by one-vertex LP Farkas duality, all $1{,}214$ $(18, 59)$ external JSON certificates, and our exact rational elimination of $p_{13}$);
    \item Within \texttt{hexh}, the interlacing-to-block derivations establishing \nolinkurl{RealizesIdx47Block}\,$S$, \nolinkurl{RealizesIdx861Block}\,$S$, and \nolinkurl{RealizesIdx488Proj}\,$S$ from the corresponding candidate spectra (proved from $1$- and $2$-vertex Jacobi systems in Step~1 of Propositions~\ref{prop:kill47}, \ref{prop:kill861}, and~\ref{prop:kill488}, and then refuted inside Lean~4 by $\neg(\chi_S = \texttt{idx47Charpoly} \wedge \texttt{RealizesIdx47Block}\,S)$, $\neg\,$\nolinkurl{RealizesIdx861Block}$\,S$, and $\neg\,$\nolinkurl{RealizesIdx488Proj}$\,S$ in \nolinkurl{frontier_19_74_core_cases_eliminated} and \nolinkurl{frontier_18_59_core_cases_eliminated}).
  \end{itemize}
\end{enumerate}

\appendix

% ============================================================================
\section{Complete enumeration and disposition of candidate spectra for \texorpdfstring{$(19, 74)$}{(19, 74)} (\texorpdfstring{$169$}{169} candidates)}\label{app:census_19_74}
% ============================================================================

Table~\ref{tab:all_19_74} lists all $169$ admissible characteristic polynomials $p_0, \dots, p_{168}$ ($\chi_S(x) = (x + 5)^{55} Q(x)$) for $(19, 74)$ from \nolinkurl{certificates/stage1/table_19_74.json}, together with the number $s = \deg H_x$ of distinct eigenvalues, the cardinality $|\Deck|$ of the vertex-deletion interlacing deck (\nolinkurl{deck_19_74.json}), and the exact elimination or reduction mechanism.

{\scriptsize
\setlength{\tabcolsep}{3.5pt}
\begin{longtable}{@{}rrr>{\raggedright\arraybackslash}p{2.50in}>{\raggedright\arraybackslash}p{2.30in}@{}}
\caption{Complete enumeration and exact disposition of all $169$ admissible characteristic polynomials for $(19, 74)$ ($\chi_S(x) = (x+5)^{55} Q(x)$).}\label{tab:all_19_74}\\
\toprule
Index & $s$ & $|\Deck|$ & Cofactor $Q(x)$ & Exact Elimination or Reduction \\
\midrule
\endfirsthead
\multicolumn{5}{@{}l}{\scriptsize\emph{Table~\ref{tab:all_19_74} (continued)}}\\
\toprule
Index & $s$ & $|\Deck|$ & Cofactor $Q(x)$ & Exact Elimination or Reduction \\
\midrule
\endhead
\midrule
\multicolumn{5}{r@{}}{\scriptsize\emph{Continued on next page}}\\
\endfoot
\bottomrule
\endlastfoot
0 & 6 & 15 & $(x-13)^6$\allowbreak $(x-15)^{10}$\allowbreak $(x^3-47x^2+723x-3637)$ & One-vertex Farkas infeasibility (Lean \texttt{decide}) \\
1 & 7 & 36 & $(x-13)^5$\allowbreak $(x-15)^{10}$\allowbreak $(x^4-60x^3+1334x^2-13028x+47161)$ & One-vertex Farkas infeasibility (Lean \texttt{decide}) \\
2 & 6 & 3 & $(x-13)^8$\allowbreak $(x-15)^7$\allowbreak $(x-17)^2$\allowbreak $(x^2-32x+251)$ & One-vertex Farkas infeasibility (Lean \texttt{decide}) \\
3 & 8 & 13 & $(x-13)^5$\allowbreak $(x-15)^9$\allowbreak $(x^2-28x+191)$\allowbreak $(x^3-47x^2+727x-3697)$ & One-vertex Farkas infeasibility (Lean \texttt{decide}) \\
4 & 8 & 24 & $(x-13)^6$\allowbreak $(x-15)^8$\allowbreak $(x-17)$\allowbreak $(x^2-28x+191)$\allowbreak $(x^2-32x+251)$ & One-vertex Farkas infeasibility (Lean \texttt{decide}) \\
5 & 7 & 14 & $(x-13)^6$\allowbreak $(x-15)^8$\allowbreak $(x-17)^2$\allowbreak $(x^3-43x^2+607x-2813)$ & One-vertex Farkas infeasibility (Lean \texttt{decide}) \\
6 & 5 & 9 & $(x-11)$\allowbreak $(x-13)^6$\allowbreak $(x-15)^9$\allowbreak $(x-17)^3$ & $6$ sols $\to 0$ ($2$-vertex Jacobi + odd-$73$ Farkas; Prop.~\ref{prop:other_19_74}) \\
7 & 8 & 123 & $(x-13)^3$\allowbreak $(x-15)^{11}$\allowbreak $(x^2-28x+191)$\allowbreak $(x^3-43x^2+603x-2753)$ & One-vertex Farkas infeasibility (Lean \texttt{decide}) \\
8 & 7 & 95 & $(x-13)^3$\allowbreak $(x-15)^{12}$\allowbreak $(x^2-26x+161)$\allowbreak $(x^2-30x+217)$ & $K^\times$ ($0$ feasible cones) \\
9 & 9 & 20 & $(x-13)^3$\allowbreak $(x-15)^{10}$\allowbreak $(x-17)$\allowbreak $(x^2-28x+191)$\allowbreak $(x^3-41x^2+551x-2423)$ & One-vertex Farkas infeasibility (Lean \texttt{decide}) \\
10 & 6 & 17 & $(x-13)^4$\allowbreak $(x-15)^{11}$\allowbreak $(x-17)^2$\allowbreak $(x^2-24x+139)$ & One-vertex Farkas infeasibility (Lean \texttt{decide}) \\
11 & 9 & 21 & $(x-13)$\allowbreak $(x-15)^{12}$\allowbreak $(x^2-28x+191)$\allowbreak $(x^4-54x^3+1080x^2-9482x+30839)$ & One-vertex Farkas infeasibility (Lean \texttt{decide}) \\
12 & 6 & 2 & $(x-13)$\allowbreak $(x-15)^{12}$\allowbreak $(x^3-41x^2+551x-2423)^2$ & One-vertex Farkas infeasibility (Lean \texttt{decide}) \\
13 & 4 & 6 & $(x-11)^3$\allowbreak $(x-15)^{15}$\allowbreak $(x-17)$ & $74\to 21\to \{S_1,S_2,S_3\}\to 0$ (exact $\mathbb{Q}$ elimination; \S\ref{sec:idx13}) \\
14 & 6 & 13 & $(x-11)$\allowbreak $(x-15)^{15}$\allowbreak $(x^3-39x^2+495x-2049)$ & One-vertex Farkas infeasibility (Lean \texttt{decide}) \\
15 & 7 & 26 & $(x-13)$\allowbreak $(x-15)^{14}$\allowbreak $(x^4-52x^3+998x^2-8372x+25873)$ & One-vertex Farkas infeasibility (Lean \texttt{decide}) \\
16 & 7 & 25 & $(x-13)$\allowbreak $(x-15)^{14}$\allowbreak $(x^4-52x^3+998x^2-8372x+25889)$ & One-vertex Farkas infeasibility (Lean \texttt{decide}) \\
17 & 6 & 1 & $(x-15)^{15}$\allowbreak $(x^2-24x+139)$\allowbreak $(x^2-26x+161)$ & One-vertex Farkas infeasibility (Lean \texttt{decide}) \\
18 & 5 & 5 & $(x-13)^5$\allowbreak $(x-15)^{10}$\allowbreak $(x^2-30x+217)^2$ & Thm.~\ref{thm:pfaffian} ($n_v=73$) + LP (Farkas) \\
19 & 6 & 8 & $(x-13)^7$\allowbreak $(x-15)^8$\allowbreak $(x-17)^2$\allowbreak $(x^2-30x+217)$ & One-vertex Farkas infeasibility (Lean \texttt{decide}) \\
20 & 8 & 70 & $(x-13)^4$\allowbreak $(x-15)^{10}$\allowbreak $(x^5-73x^4+2114x^3-30354x^2+216093x-610213)$ & One-vertex Farkas infeasibility (Lean \texttt{decide}) \\
21 & 7 & 36 & $(x-13)^6$\allowbreak $(x-15)^9$\allowbreak $(x-17)$\allowbreak $(x^3-45x^2+663x-3187)$ & One-vertex Farkas infeasibility (Lean \texttt{decide}) \\
22 & 8 & 103 & $(x-13)^3$\allowbreak $(x-15)^{11}$\allowbreak $(x^5-71x^4+1998x^3-27850x^2+192241x-525583)$ & One-vertex Farkas infeasibility (Lean \texttt{decide}) \\
23 & 9 & 38 & $(x-13)^3$\allowbreak $(x-15)^{10}$\allowbreak $(x^2-28x+191)$\allowbreak $(x^4-58x^3+1248x^2-11798x+41311)$ & One-vertex Farkas infeasibility (Lean \texttt{decide}) \\
24 & 6 & 7 & $(x-13)^3$\allowbreak $(x-15)^9$\allowbreak $(x-17)$\allowbreak $(x^2-28x+191)^3$ & $4$ sols $\to 0$ ($2$-vertex Jacobi + odd-$73$ Farkas + $E_{17}$; Prop.~\ref{prop:other_19_74}) \\
25 & 7 & 26 & $(x-11)$\allowbreak $(x-13)^2$\allowbreak $(x-15)^{13}$\allowbreak $(x^3-43x^2+603x-2753)$ & Thm.~\ref{thm:pfaffian} ($n_v=73$) + LP (Farkas) \\
26 & 5 & 6 & $(x-11)^2$\allowbreak $(x-13)^3$\allowbreak $(x-15)^{12}$\allowbreak $(x-17)^2$ & One-vertex Farkas infeasibility (Lean \texttt{decide}) \\
27 & 7 & 37 & $(x-13)^2$\allowbreak $(x-15)^{13}$\allowbreak $(x^4-54x^3+1076x^2-9378x+30179)$ & One-vertex Farkas infeasibility (Lean \texttt{decide}) \\
28 & 6 & 19 & $(x-11)$\allowbreak $(x-13)^2$\allowbreak $(x-15)^{14}$\allowbreak $(x^2-28x+183)$ & One-vertex Farkas infeasibility (Lean \texttt{decide}) \\
29 & 6 & 35 & $(x-13)^3$\allowbreak $(x-15)^{13}$\allowbreak $(x^3-41x^2+543x-2311)$ & One-vertex Farkas infeasibility (Lean \texttt{decide}) \\
30 & 7 & 45 & $(x-13)^2$\allowbreak $(x-15)^{13}$\allowbreak $(x^4-54x^3+1076x^2-9370x+30091)$ & One-vertex Farkas infeasibility (Lean \texttt{decide}) \\
31 & 7 & 42 & $(x-13)^2$\allowbreak $(x-15)^{13}$\allowbreak $(x-17)$\allowbreak $(x^3-37x^2+447x-1763)$ & One-vertex Farkas infeasibility (Lean \texttt{decide}) \\
32 & 7 & 14 & $(x-13)$\allowbreak $(x-15)^{14}$\allowbreak $(x^4-52x^3+998x^2-8364x+25769)$ & One-vertex Farkas infeasibility (Lean \texttt{decide}) \\
33 & 6 & 10 & $(x-13)$\allowbreak $(x-15)^{15}$\allowbreak $(x^3-37x^2+443x-1711)$ & One-vertex Farkas infeasibility (Lean \texttt{decide}) \\
34 & 5 & 7 & $(x-13)^6$\allowbreak $(x-15)^{11}$\allowbreak $(x^2-32x+243)$ & One-vertex Farkas infeasibility (Lean \texttt{decide}) \\
35 & 4 & 2 & $(x-13)^9$\allowbreak $(x-15)^6$\allowbreak $(x-17)^4$ & Prop.~\ref{prop:prop61} ($k=44$, $\tr T^2 \ge 1625 > 870$) \\
36 & 6 & 24 & $(x-13)^4$\allowbreak $(x-15)^{12}$\allowbreak $(x^3-43x^2+599x-2709)$ & One-vertex Farkas infeasibility (Lean \texttt{decide}) \\
37 & 7 & 71 & $(x-13)^5$\allowbreak $(x-15)^{10}$\allowbreak $(x-17)$\allowbreak $(x^3-43x^2+603x-2753)$ & One-vertex Farkas infeasibility (Lean \texttt{decide}) \\
38 & 6 & 4 & $(x-13)^5$\allowbreak $(x-15)^8$\allowbreak $(x-17)^2$\allowbreak $(x^2-28x+191)^2$ & Thm.~\ref{thm:pfaffian} ($n_v=73$) + LP (Farkas) \\
39 & 7 & 69 & $(x-13)^4$\allowbreak $(x-15)^{11}$\allowbreak $(x^4-58x^3+1244x^2-11678x+40427)$ & One-vertex Farkas infeasibility (Lean \texttt{decide}) \\
40 & 6 & 40 & $(x-13)^5$\allowbreak $(x-15)^{11}$\allowbreak $(x-17)$\allowbreak $(x^2-28x+183)$ & $K^\times$ RREF ($2$ cones, $0$ $\mathbb{Z}_{\ge 0}$ sols) \\
41 & 8 & 45 & $(x-11)$\allowbreak $(x-13)^3$\allowbreak $(x-15)^{11}$\allowbreak $(x-17)$\allowbreak $(x^3-43x^2+607x-2813)$ & One-vertex Farkas infeasibility (Lean \texttt{decide}) \\
42 & 8 & 7 & $(x-13)^2$\allowbreak $(x-15)^{10}$\allowbreak $(x^2-28x+191)^2$\allowbreak $(x^3-43x^2+607x-2813)$ & One-vertex Farkas infeasibility (Lean \texttt{decide}) \\
43 & 7 & 11 & $(x-13)^3$\allowbreak $(x-15)^{12}$\allowbreak $(x^2-24x+139)$\allowbreak $(x^2-32x+251)$ & One-vertex Farkas infeasibility (Lean \texttt{decide}) \\
44 & 7 & 53 & $(x-13)^3$\allowbreak $(x-15)^{12}$\allowbreak $(x^4-56x^3+1158x^2-10464x+34817)$ & One-vertex Farkas infeasibility (Lean \texttt{decide}) \\
45 & 8 & 98 & $(x-13)^2$\allowbreak $(x-15)^{12}$\allowbreak $(x^2-28x+191)$\allowbreak $(x^3-41x^2+547x-2371)$ & One-vertex Farkas infeasibility (Lean \texttt{decide}) \\
46 & 7 & 15 & $(x-11)^2$\allowbreak $(x-13)$\allowbreak $(x-15)^{13}$\allowbreak $(x-17)$\allowbreak $(x^2-28x+191)$ & One-vertex Farkas infeasibility (Lean \texttt{decide}) \\
47 & 5 & 3 & $(x-11)$\allowbreak $(x-15)^{12}$\allowbreak $(x^2-28x+191)^3$ & $[52,16,6]$, $WB+BX=10B \Rightarrow 928\le 896$ (Prop.~\ref{prop:kill47}) \\
48 & 6 & 23 & $(x-11)$\allowbreak $(x-13)$\allowbreak $(x-15)^{15}$\allowbreak $(x^2-26x+157)$ & One-vertex Farkas infeasibility (Lean \texttt{decide}) \\
49 & 6 & 13 & $(x-13)^2$\allowbreak $(x-15)^{14}$\allowbreak $(x^3-39x^2+491x-1981)$ & One-vertex Farkas infeasibility (Lean \texttt{decide}) \\
50 & 6 & 18 & $(x-13)^5$\allowbreak $(x-15)^{11}$\allowbreak $(x^3-45x^2+659x-3151)$ & One-vertex Farkas infeasibility (Lean \texttt{decide}) \\
51 & 5 & 20 & $(x-11)$\allowbreak $(x-13)^5$\allowbreak $(x-15)^{12}$\allowbreak $(x-19)$ & $K^\times$ RREF (non-integer sols only) \\
52 & 6 & 31 & $(x-13)^6$\allowbreak $(x-15)^{10}$\allowbreak $(x-17)$\allowbreak $(x^2-30x+213)$ & One-vertex Farkas infeasibility (Lean \texttt{decide}) \\
53 & 6 & 41 & $(x-13)^5$\allowbreak $(x-15)^{11}$\allowbreak $(x^3-45x^2+659x-3127)$ & One-vertex Farkas infeasibility (Lean \texttt{decide}) \\
54 & 7 & 8 & $(x-13)^3$\allowbreak $(x-15)^{10}$\allowbreak $(x^2-28x+191)^2$\allowbreak $(x^2-30x+217)$ & One-vertex Farkas infeasibility (Lean \texttt{decide}) \\
55 & 6 & 25 & $(x-13)^4$\allowbreak $(x-15)^{12}$\allowbreak $(x^3-43x^2+599x-2693)$ & One-vertex Farkas infeasibility (Lean \texttt{decide}) \\
56 & 7 & 52 & $(x-13)^3$\allowbreak $(x-15)^{12}$\allowbreak $(x^4-56x^3+1158x^2-10480x+35041)$ & One-vertex Farkas infeasibility (Lean \texttt{decide}) \\
57 & 8 & 120 & $(x-13)^3$\allowbreak $(x-15)^{11}$\allowbreak $(x^2-30x+217)$\allowbreak $(x^3-41x^2+551x-2423)$ & $K^\times$ ($0$ feasible cones) \\
58 & 7 & 52 & $(x-13)^2$\allowbreak $(x-15)^{13}$\allowbreak $(x^4-54x^3+1076x^2-9386x+30267)$ & One-vertex Farkas infeasibility (Lean \texttt{decide}) \\
59 & 6 & 15 & $(x-13)^5$\allowbreak $(x-15)^{10}$\allowbreak $(x-17)^2$\allowbreak $(x^2-26x+161)$ & One-vertex Farkas infeasibility (Lean \texttt{decide}) \\
60 & 9 & 13 & $(x-11)$\allowbreak $(x-13)$\allowbreak $(x-15)^{12}$\allowbreak $(x^2-28x+191)$\allowbreak $(x^3-43x^2+607x-2813)$ & One-vertex Farkas infeasibility (Lean \texttt{decide}) \\
61 & 7 & 25 & $(x-13)^3$\allowbreak $(x-15)^{12}$\allowbreak $(x-17)$\allowbreak $(x^3-39x^2+495x-2049)$ & One-vertex Farkas infeasibility (Lean \texttt{decide}) \\
62 & 6 & 20 & $(x-13)^4$\allowbreak $(x-15)^{12}$\allowbreak $(x-17)$\allowbreak $(x^2-26x+157)$ & One-vertex Farkas infeasibility (Lean \texttt{decide}) \\
63 & 7 & 15 & $(x-11)$\allowbreak $(x-15)^{14}$\allowbreak $(x^4-54x^3+1080x^2-9482x+30839)$ & One-vertex Farkas infeasibility (Lean \texttt{decide}) \\
64 & 7 & 48 & $(x-11)$\allowbreak $(x-13)^2$\allowbreak $(x-15)^{13}$\allowbreak $(x-17)$\allowbreak $(x^2-26x+161)$ & One-vertex Farkas infeasibility (Lean \texttt{decide}) \\
65 & 7 & 42 & $(x-11)$\allowbreak $(x-13)$\allowbreak $(x-15)^{14}$\allowbreak $(x-17)$\allowbreak $(x^2-24x+139)$ & Thm.~\ref{thm:pfaffian} ($n_v=73$) + LP (Farkas) \\
66 & 7 & 16 & $(x-13)^2$\allowbreak $(x-15)^{13}$\allowbreak $(x^4-54x^3+1076x^2-9354x+29851)$ & One-vertex Farkas infeasibility (Lean \texttt{decide}) \\
67 & 5 & 4 & $(x-13)^8$\allowbreak $(x-15)^9$\allowbreak $(x-17)$\allowbreak $(x-19)$ & $K^\times$ ($0$ feasible cones) \\
68 & 7 & 61 & $(x-13)^5$\allowbreak $(x-15)^{10}$\allowbreak $(x^4-60x^3+1334x^2-13020x+47041)$ & One-vertex Farkas infeasibility (Lean \texttt{decide}) \\
69 & 6 & 30 & $(x-13)^4$\allowbreak $(x-15)^{12}$\allowbreak $(x^3-43x^2+599x-2701)$ & One-vertex Farkas infeasibility (Lean \texttt{decide}) \\
70 & 6 & 23 & $(x-11)$\allowbreak $(x-13)^3$\allowbreak $(x-15)^{13}$\allowbreak $(x^2-30x+213)$ & One-vertex Farkas infeasibility (Lean \texttt{decide}) \\
71 & 8 & 25 & $(x-13)^4$\allowbreak $(x-15)^{10}$\allowbreak $(x^2-28x+191)$\allowbreak $(x^3-45x^2+663x-3187)$ & One-vertex Farkas infeasibility (Lean \texttt{decide}) \\
72 & 7 & 26 & $(x-11)$\allowbreak $(x-13)^3$\allowbreak $(x-15)^{12}$\allowbreak $(x^3-45x^2+663x-3187)$ & One-vertex Farkas infeasibility (Lean \texttt{decide}) \\
73 & 7 & 8 & $(x-13)^5$\allowbreak $(x-15)^9$\allowbreak $(x-17)^2$\allowbreak $(x^3-41x^2+551x-2423)$ & One-vertex Farkas infeasibility (Lean \texttt{decide}) \\
74 & 5 & 27 & $(x-13)^4$\allowbreak $(x-15)^{13}$\allowbreak $(x^2-28x+179)$ & Odd-$73$ submatrix Galois RREF: $2(x_1+x_4)=29$ (Prop.~\ref{prop:other_19_74}) \\
75 & 8 & 76 & $(x-13)^2$\allowbreak $(x-15)^{12}$\allowbreak $(x^5-69x^4+1886x^3-25526x^2+171073x-454213)$ & One-vertex Farkas infeasibility (Lean \texttt{decide}) \\
76 & 6 & 13 & $(x-11)^2$\allowbreak $(x-13)$\allowbreak $(x-15)^{14}$\allowbreak $(x^2-30x+217)$ & One-vertex Farkas infeasibility (Lean \texttt{decide}) \\
77 & 7 & 47 & $(x-13)^2$\allowbreak $(x-15)^{13}$\allowbreak $(x^2-24x+139)$\allowbreak $(x^2-30x+217)$ & One-vertex Farkas infeasibility (Lean \texttt{decide}) \\
78 & 8 & 21 & $(x-11)$\allowbreak $(x-13)$\allowbreak $(x-15)^{13}$\allowbreak $(x^4-56x^3+1162x^2-10584x+35701)$ & One-vertex Farkas infeasibility (Lean \texttt{decide}) \\
79 & 7 & 36 & $(x-13)^3$\allowbreak $(x-15)^{12}$\allowbreak $(x^4-56x^3+1158x^2-10456x+34681)$ & One-vertex Farkas infeasibility (Lean \texttt{decide}) \\
80 & 8 & 83 & $(x-13)^2$\allowbreak $(x-15)^{12}$\allowbreak $(x^5-69x^4+1886x^3-25510x^2+170625x-451093)$ & One-vertex Farkas infeasibility (Lean \texttt{decide}) \\
81 & 8 & 22 & $(x-13)$\allowbreak $(x-15)^{13}$\allowbreak $(x^2-28x+191)$\allowbreak $(x^3-39x^2+495x-2049)$ & One-vertex Farkas infeasibility (Lean \texttt{decide}) \\
82 & 6 & 30 & $(x-13)^3$\allowbreak $(x-15)^{13}$\allowbreak $(x^3-41x^2+543x-2303)$ & One-vertex Farkas infeasibility (Lean \texttt{decide}) \\
83 & 7 & 6 & $(x-11)$\allowbreak $(x-15)^{14}$\allowbreak $(x^2-26x+161)$\allowbreak $(x^2-28x+191)$ & One-vertex Farkas infeasibility (Lean \texttt{decide}) \\
84 & 5 & 0 & $(x-15)^{16}$\allowbreak $(x^3-35x^2+399x-1477)$ & One-vertex Farkas infeasibility (Lean \texttt{decide}) \\
85 & 7 & 82 & $(x-13)^5$\allowbreak $(x-15)^{10}$\allowbreak $(x^4-60x^3+1334x^2-13012x+46921)$ & One-vertex Farkas infeasibility (Lean \texttt{decide}) \\
86 & 7 & 12 & $(x-13)^4$\allowbreak $(x-15)^9$\allowbreak $(x^2-28x+191)^2$\allowbreak $(x^2-32x+251)$ & One-vertex Farkas infeasibility (Lean \texttt{decide}) \\
87 & 8 & 17 & $(x-13)^5$\allowbreak $(x-15)^9$\allowbreak $(x-17)$\allowbreak $(x^2-28x+191)$\allowbreak $(x^2-30x+217)$ & One-vertex Farkas infeasibility (Lean \texttt{decide}) \\
88 & 8 & 16 & $(x-13)^4$\allowbreak $(x-15)^{10}$\allowbreak $(x^2-32x+251)$\allowbreak $(x^3-41x^2+551x-2423)$ & One-vertex Farkas infeasibility (Lean \texttt{decide}) \\
89 & 7 & 79 & $(x-13)^4$\allowbreak $(x-15)^{11}$\allowbreak $(x^4-58x^3+1244x^2-11686x+40563)$ & One-vertex Farkas infeasibility (Lean \texttt{decide}) \\
90 & 7 & 64 & $(x-13)^4$\allowbreak $(x-15)^{11}$\allowbreak $(x^2-26x+161)$\allowbreak $(x^2-32x+251)$ & One-vertex Farkas infeasibility (Lean \texttt{decide}) \\
91 & 8 & 32 & $(x-13)^4$\allowbreak $(x-15)^{10}$\allowbreak $(x-17)$\allowbreak $(x^4-56x^3+1162x^2-10584x+35701)$ & One-vertex Farkas infeasibility (Lean \texttt{decide}) \\
92 & 7 & 63 & $(x-13)^4$\allowbreak $(x-15)^{11}$\allowbreak $(x^4-58x^3+1244x^2-11670x+40291)$ & One-vertex Farkas infeasibility (Lean \texttt{decide}) \\
93 & 9 & 23 & $(x-13)^2$\allowbreak $(x-15)^{11}$\allowbreak $(x^3-41x^2+551x-2423)$\allowbreak $(x^3-43x^2+607x-2813)$ & One-vertex Farkas infeasibility (Lean \texttt{decide}) \\
94 & 7 & 14 & $(x-11)$\allowbreak $(x-13)^2$\allowbreak $(x-15)^{11}$\allowbreak $(x-17)$\allowbreak $(x^2-28x+191)^2$ & One-vertex Farkas infeasibility (Lean \texttt{decide}) \\
95 & 5 & 2 & $(x-13)$\allowbreak $(x-15)^{10}$\allowbreak $(x^2-28x+191)^4$ & One-vertex Farkas infeasibility (Lean \texttt{decide}) \\
96 & 8 & 12 & $(x-13)$\allowbreak $(x-15)^{13}$\allowbreak $(x^2-26x+161)$\allowbreak $(x^3-41x^2+551x-2423)$ & One-vertex Farkas infeasibility (Lean \texttt{decide}) \\
97 & 8 & 10 & $(x-13)$\allowbreak $(x-15)^{13}$\allowbreak $(x^2-28x+191)$\allowbreak $(x^3-39x^2+495x-2033)$ & One-vertex Farkas infeasibility (Lean \texttt{decide}) \\
98 & 6 & 19 & $(x-13)$\allowbreak $(x-15)^{15}$\allowbreak $(x-9)$\allowbreak $(x^2-28x+191)$ & $K^\times$ ($0$ feasible cones) \\
99 & 6 & 4 & $(x-15)^{15}$\allowbreak $(x^4-50x^3+924x^2-7470x+22259)$ & One-vertex Farkas infeasibility (Lean \texttt{decide}) \\
100 & 5 & 2 & $(x-13)^7$\allowbreak $(x-15)^8$\allowbreak $(x^2-32x+251)^2$ & One-vertex Farkas infeasibility (Lean \texttt{decide}) \\
101 & 7 & 10 & $(x-13)^7$\allowbreak $(x-15)^8$\allowbreak $(x-17)$\allowbreak $(x^3-47x^2+727x-3697)$ & One-vertex Farkas infeasibility (Lean \texttt{decide}) \\
102 & 8 & 19 & $(x-13)^4$\allowbreak $(x-15)^{10}$\allowbreak $(x^2-28x+191)$\allowbreak $(x^3-45x^2+663x-3203)$ & One-vertex Farkas infeasibility (Lean \texttt{decide}) \\
103 & 6 & 70 & $(x-13)^5$\allowbreak $(x-15)^{11}$\allowbreak $(x^3-45x^2+659x-3119)$ & One-vertex Farkas infeasibility (Lean \texttt{decide}) \\
104 & 8 & 104 & $(x-11)$\allowbreak $(x-13)^3$\allowbreak $(x-15)^{11}$\allowbreak $(x^2-28x+191)$\allowbreak $(x^2-32x+251)$ & One-vertex Farkas infeasibility (Lean \texttt{decide}) \\
105 & 7 & 64 & $(x-11)$\allowbreak $(x-13)^4$\allowbreak $(x-15)^{11}$\allowbreak $(x-17)$\allowbreak $(x^2-30x+217)$ & One-vertex Farkas infeasibility (Lean \texttt{decide}) \\
106 & 7 & 33 & $(x-13)^4$\allowbreak $(x-15)^{11}$\allowbreak $(x-17)$\allowbreak $(x^3-41x^2+547x-2371)$ & Thm.~\ref{thm:pfaffian} ($n_v=73$) + LP (Farkas) \\
107 & 7 & 96 & $(x-13)^3$\allowbreak $(x-15)^{12}$\allowbreak $(x^2-28x+183)$\allowbreak $(x^2-28x+191)$ & Thm.~\ref{thm:pfaffian} ($n_v=73$) + LP (Farkas) \\
108 & 9 & 27 & $(x-13)^2$\allowbreak $(x-15)^{11}$\allowbreak $(x^2-28x+191)$\allowbreak $(x^4-56x^3+1162x^2-10584x+35701)$ & One-vertex Farkas infeasibility (Lean \texttt{decide}) \\
109 & 6 & 29 & $(x-13)^3$\allowbreak $(x-15)^{13}$\allowbreak $(x^3-41x^2+543x-2319)$ & One-vertex Farkas infeasibility (Lean \texttt{decide}) \\
110 & 8 & 29 & $(x-11)$\allowbreak $(x-13)^2$\allowbreak $(x-15)^{12}$\allowbreak $(x-17)$\allowbreak $(x^3-41x^2+551x-2423)$ & One-vertex Farkas infeasibility (Lean \texttt{decide}) \\
111 & 8 & 10 & $(x-13)$\allowbreak $(x-15)^{13}$\allowbreak $(x^2-24x+139)$\allowbreak $(x^3-43x^2+607x-2813)$ & One-vertex Farkas infeasibility (Lean \texttt{decide}) \\
112 & 8 & 38 & $(x-13)$\allowbreak $(x-15)^{13}$\allowbreak $(x^5-67x^4+1778x^3-23350x^2+151693x-389863)$ & One-vertex Farkas infeasibility (Lean \texttt{decide}) \\
113 & 6 & 12 & $(x-11)$\allowbreak $(x-15)^{15}$\allowbreak $(x^3-39x^2+495x-2033)$ & One-vertex Farkas infeasibility (Lean \texttt{decide}) \\
114 & 5 & 1 & $(x-13)$\allowbreak $(x-15)^{16}$\allowbreak $(x^2-22x+113)$ & One-vertex Farkas infeasibility (Lean \texttt{decide}) \\
115 & 6 & 4 & $(x-13)^7$\allowbreak $(x-15)^9$\allowbreak $(x^3-49x^2+791x-4207)$ & One-vertex Farkas infeasibility (Lean \texttt{decide}) \\
116 & 6 & 26 & $(x-13)^6$\allowbreak $(x-15)^{10}$\allowbreak $(x-19)$\allowbreak $(x^2-28x+191)$ & One-vertex Farkas infeasibility (Lean \texttt{decide}) \\
117 & 7 & 8 & $(x-13)^6$\allowbreak $(x-15)^9$\allowbreak $(x^2-30x+217)$\allowbreak $(x^2-32x+251)$ & One-vertex Farkas infeasibility (Lean \texttt{decide}) \\
118 & 7 & 21 & $(x-13)^4$\allowbreak $(x-15)^{11}$\allowbreak $(x^4-58x^3+1244x^2-11702x+40771)$ & One-vertex Farkas infeasibility (Lean \texttt{decide}) \\
119 & 7 & 51 & $(x-13)^5$\allowbreak $(x-15)^{10}$\allowbreak $(x^4-60x^3+1334x^2-13020x+47057)$ & Thm.~\ref{thm:pfaffian} ($n_v=73$) + LP (Farkas) \\
120 & 7 & 27 & $(x-11)$\allowbreak $(x-13)^3$\allowbreak $(x-15)^{12}$\allowbreak $(x^3-45x^2+663x-3203)$ & One-vertex Farkas infeasibility (Lean \texttt{decide}) \\
121 & 7 & 72 & $(x-11)$\allowbreak $(x-13)^5$\allowbreak $(x-15)^{10}$\allowbreak $(x-17)$\allowbreak $(x^2-32x+251)$ & One-vertex Farkas infeasibility (Lean \texttt{decide}) \\
122 & 8 & 84 & $(x-13)^3$\allowbreak $(x-15)^{11}$\allowbreak $(x^5-71x^4+1998x^3-27858x^2+192465x-527143)$ & One-vertex Farkas infeasibility (Lean \texttt{decide}) \\
123 & 9 & 21 & $(x-13)^4$\allowbreak $(x-15)^9$\allowbreak $(x-17)$\allowbreak $(x^2-28x+191)$\allowbreak $(x^3-43x^2+607x-2813)$ & One-vertex Farkas infeasibility (Lean \texttt{decide}) \\
124 & 7 & 41 & $(x-13)^3$\allowbreak $(x-15)^{12}$\allowbreak $(x^4-56x^3+1158x^2-10472x+34921)$ & One-vertex Farkas infeasibility (Lean \texttt{decide}) \\
125 & 8 & 25 & $(x-13)^3$\allowbreak $(x-15)^{11}$\allowbreak $(x-17)$\allowbreak $(x^4-54x^3+1080x^2-9482x+30839)$ & One-vertex Farkas infeasibility (Lean \texttt{decide}) \\
126 & 6 & 27 & $(x-13)^4$\allowbreak $(x-15)^{12}$\allowbreak $(x^3-43x^2+599x-2677)$ & One-vertex Farkas infeasibility (Lean \texttt{decide}) \\
127 & 8 & 94 & $(x-13)^2$\allowbreak $(x-15)^{12}$\allowbreak $(x^5-69x^4+1886x^3-25518x^2+170849x-452653)$ & One-vertex Farkas infeasibility (Lean \texttt{decide}) \\
128 & 8 & 91 & $(x-13)^2$\allowbreak $(x-15)^{12}$\allowbreak $(x^2-26x+161)$\allowbreak $(x^3-43x^2+607x-2813)$ & One-vertex Farkas infeasibility (Lean \texttt{decide}) \\
129 & 8 & 35 & $(x-13)^3$\allowbreak $(x-15)^{11}$\allowbreak $(x-17)$\allowbreak $(x^2-26x+161)$\allowbreak $(x^2-28x+191)$ & One-vertex Farkas infeasibility (Lean \texttt{decide}) \\
130 & 8 & 79 & $(x-13)^2$\allowbreak $(x-15)^{12}$\allowbreak $(x-17)$\allowbreak $(x^2-24x+139)$\allowbreak $(x^2-28x+191)$ & One-vertex Farkas infeasibility (Lean \texttt{decide}) \\
131 & 7 & 7 & $(x-13)$\allowbreak $(x-15)^{12}$\allowbreak $(x^2-26x+161)$\allowbreak $(x^2-28x+191)^2$ & One-vertex Farkas infeasibility (Lean \texttt{decide}) \\
132 & 6 & 21 & $(x-13)^2$\allowbreak $(x-15)^{14}$\allowbreak $(x^3-39x^2+491x-1997)$ & One-vertex Farkas infeasibility (Lean \texttt{decide}) \\
133 & 7 & 28 & $(x-13)$\allowbreak $(x-15)^{14}$\allowbreak $(x^4-52x^3+998x^2-8380x+25977)$ & One-vertex Farkas infeasibility (Lean \texttt{decide}) \\
134 & 5 & 18 & $(x-13)^3$\allowbreak $(x-15)^{14}$\allowbreak $(x-17)$\allowbreak $(x-9)$ & Exact $2$-vertex Jacobi ($2^{72}$ + interlacing $\to 0$; Prop.~\ref{prop:other_19_74}) \\
135 & 7 & 5 & $(x-15)^{14}$\allowbreak $(x^2-28x+191)$\allowbreak $(x^3-37x^2+447x-1763)$ & One-vertex Farkas infeasibility (Lean \texttt{decide}) \\
136 & 4 & 2 & $(x-11)$\allowbreak $(x-15)^{17}$\allowbreak $(x-9)$ & $\Leftrightarrow \SRG(75,32,10,16)$ (Thm.~\ref{thm:srg136};~\cite{AzarijaMarc2018}) \\
137 & 6 & 24 & $(x-13)^5$\allowbreak $(x-15)^{11}$\allowbreak $(x^3-45x^2+659x-3143)$ & One-vertex Farkas infeasibility (Lean \texttt{decide}) \\
138 & 7 & 34 & $(x-13)^6$\allowbreak $(x-15)^9$\allowbreak $(x-17)$\allowbreak $(x^3-45x^2+663x-3203)$ & One-vertex Farkas infeasibility (Lean \texttt{decide}) \\
139 & 7 & 25 & $(x-13)^4$\allowbreak $(x-15)^{11}$\allowbreak $(x^2-28x+191)$\allowbreak $(x^2-30x+213)$ & One-vertex Farkas infeasibility (Lean \texttt{decide}) \\
140 & 8 & 19 & $(x-13)^4$\allowbreak $(x-15)^{10}$\allowbreak $(x^2-30x+217)$\allowbreak $(x^3-43x^2+607x-2813)$ & One-vertex Farkas infeasibility (Lean \texttt{decide}) \\
141 & 7 & 50 & $(x-13)^4$\allowbreak $(x-15)^{11}$\allowbreak $(x^4-58x^3+1244x^2-11686x+40531)$ & One-vertex Farkas infeasibility (Lean \texttt{decide}) \\
142 & 8 & 43 & $(x-13)^5$\allowbreak $(x-15)^9$\allowbreak $(x-17)$\allowbreak $(x^4-58x^3+1248x^2-11798x+41311)$ & One-vertex Farkas infeasibility (Lean \texttt{decide}) \\
143 & 6 & 35 & $(x-13)^3$\allowbreak $(x-15)^{13}$\allowbreak $(x^3-41x^2+543x-2335)$ & One-vertex Farkas infeasibility (Lean \texttt{decide}) \\
144 & 6 & 13 & $(x-11)^2$\allowbreak $(x-13)^2$\allowbreak $(x-15)^{13}$\allowbreak $(x^2-32x+251)$ & One-vertex Farkas infeasibility (Lean \texttt{decide}) \\
145 & 8 & 84 & $(x-13)^3$\allowbreak $(x-15)^{11}$\allowbreak $(x^5-71x^4+1998x^3-27842x^2+192017x-524023)$ & One-vertex Farkas infeasibility (Lean \texttt{decide}) \\
146 & 8 & 31 & $(x-11)$\allowbreak $(x-13)^2$\allowbreak $(x-15)^{12}$\allowbreak $(x^4-58x^3+1248x^2-11798x+41311)$ & One-vertex Farkas infeasibility (Lean \texttt{decide}) \\
147 & 8 & 7 & $(x-13)$\allowbreak $(x-15)^{11}$\allowbreak $(x^2-28x+191)^2$\allowbreak $(x^3-41x^2+551x-2423)$ & One-vertex Farkas infeasibility (Lean \texttt{decide}) \\
148 & 6 & 25 & $(x-13)^2$\allowbreak $(x-15)^{14}$\allowbreak $(x^3-39x^2+491x-2005)$ & One-vertex Farkas infeasibility (Lean \texttt{decide}) \\
149 & 7 & 32 & $(x-11)$\allowbreak $(x-13)$\allowbreak $(x-15)^{14}$\allowbreak $(x^3-41x^2+547x-2371)$ & Thm.~\ref{thm:pfaffian} ($n_v=73$) + LP (Farkas) \\
150 & 8 & 11 & $(x-11)$\allowbreak $(x-15)^{13}$\allowbreak $(x^2-28x+191)$\allowbreak $(x^3-41x^2+551x-2423)$ & One-vertex Farkas infeasibility (Lean \texttt{decide}) \\
151 & 6 & 6 & $(x-15)^{13}$\allowbreak $(x^2-24x+139)$\allowbreak $(x^2-28x+191)^2$ & One-vertex Farkas infeasibility (Lean \texttt{decide}) \\
152 & 5 & 4 & $(x-13)$\allowbreak $(x-15)^{14}$\allowbreak $(x^2-26x+161)^2$ & One-vertex Farkas infeasibility (Lean \texttt{decide}) \\
153 & 5 & 5 & $(x-13)^2$\allowbreak $(x-15)^{15}$\allowbreak $(x^2-24x+131)$ & One-vertex Farkas infeasibility (Lean \texttt{decide}) \\
154 & 5 & 4 & $(x-13)^7$\allowbreak $(x-15)^{10}$\allowbreak $(x^2-34x+281)$ & One-vertex Farkas infeasibility (Lean \texttt{decide}) \\
155 & 7 & 18 & $(x-13)^6$\allowbreak $(x-15)^9$\allowbreak $(x^4-62x^3+1428x^2-14482x+54571)$ & One-vertex Farkas infeasibility (Lean \texttt{decide}) \\
156 & 8 & 15 & $(x-13)^5$\allowbreak $(x-15)^9$\allowbreak $(x^2-32x+251)$\allowbreak $(x^3-43x^2+607x-2813)$ & One-vertex Farkas infeasibility (Lean \texttt{decide}) \\
157 & 7 & 67 & $(x-11)$\allowbreak $(x-13)^4$\allowbreak $(x-15)^{11}$\allowbreak $(x^3-47x^2+727x-3697)$ & One-vertex Farkas infeasibility (Lean \texttt{decide}) \\
158 & 6 & 7 & $(x-13)^7$\allowbreak $(x-15)^7$\allowbreak $(x-17)^3$\allowbreak $(x^2-28x+191)$ & One-vertex Farkas infeasibility (Lean \texttt{decide}) \\
159 & 7 & 29 & $(x-13)^3$\allowbreak $(x-15)^{12}$\allowbreak $(x^4-56x^3+1158x^2-10488x+35129)$ & One-vertex Farkas infeasibility (Lean \texttt{decide}) \\
160 & 6 & 3 & $(x-13)^3$\allowbreak $(x-15)^{10}$\allowbreak $(x^3-43x^2+607x-2813)^2$ & One-vertex Farkas infeasibility (Lean \texttt{decide}) \\
161 & 8 & 23 & $(x-11)$\allowbreak $(x-13)^2$\allowbreak $(x-15)^{12}$\allowbreak $(x^2-28x+191)$\allowbreak $(x^2-30x+217)$ & One-vertex Farkas infeasibility (Lean \texttt{decide}) \\
162 & 7 & 20 & $(x-11)$\allowbreak $(x-13)^4$\allowbreak $(x-15)^{10}$\allowbreak $(x-17)^2$\allowbreak $(x^2-28x+191)$ & One-vertex Farkas infeasibility (Lean \texttt{decide}) \\
163 & 6 & 6 & $(x-11)^2$\allowbreak $(x-15)^{14}$\allowbreak $(x^3-43x^2+607x-2813)$ & One-vertex Farkas infeasibility (Lean \texttt{decide}) \\
164 & 7 & 34 & $(x-13)^2$\allowbreak $(x-15)^{13}$\allowbreak $(x^4-54x^3+1076x^2-9370x+30059)$ & One-vertex Farkas infeasibility (Lean \texttt{decide}) \\
165 & 7 & 35 & $(x-13)^3$\allowbreak $(x-15)^{12}$\allowbreak $(x-17)$\allowbreak $(x^3-39x^2+495x-2033)$ & One-vertex Farkas infeasibility (Lean \texttt{decide}) \\
166 & 7 & 17 & $(x-13)^2$\allowbreak $(x-15)^{13}$\allowbreak $(x^2-26x+157)$\allowbreak $(x^2-28x+191)$ & One-vertex Farkas infeasibility (Lean \texttt{decide}) \\
167 & 7 & 6 & $(x-15)^{14}$\allowbreak $(x^2-24x+139)$\allowbreak $(x^3-41x^2+551x-2423)$ & One-vertex Farkas infeasibility (Lean \texttt{decide}) \\
168 & 6 & 9 & $(x-13)^2$\allowbreak $(x-15)^{14}$\allowbreak $(x^3-39x^2+491x-1973)$ & One-vertex Farkas infeasibility (Lean \texttt{decide}) \\
\end{longtable}
}

% ============================================================================
\section{Remaining candidate spectra for \texorpdfstring{$(18, 59)$}{(18, 59)} and literature calibration tables}\label{app:census_18_59}
% ============================================================================

\subsection{Classification of all \texorpdfstring{$1{,}214$}{1,214} admissible characteristic polynomials for \texorpdfstring{$(18, 59)$}{(18, 59)}}
Table~\ref{tab:18_59_by_s} cross-tabulates all $1{,}214$ admissible characteristic polynomials $q_0, \dots, q_{1213}$ of $(18, 59)$ compatible with Corollary~\ref{cor:cor43_full} (\nolinkurl{certificates/stage1/table_18_59_cor43.json}, total interlacing deck cardinality $781{,}998$) by distinct-eigenvalue count $s \in \{4, \dots, 11\}$ and exact elimination mechanism. All $1{,}128$ one-vertex compatibility eliminations ($572$ \texttt{1-vertex\_farkas}, $552$ \texttt{kx\_rref\_exact}, and $4$ \texttt{prop61\_exact}) are archived with explicit exact duals and cone proofs in \nolinkurl{certificates/phase2/phase2_18_59.json}; Table~\ref{tab:surv_18_59} then records every one of the $86$ remaining candidate spectra ($66$ ruled out by Eulerian gauge and Galois orthogonality, and $20$ by higher-order obstructions).

\begin{table}[htbp]
\centering
\caption{Classification of all $1{,}214$ admissible characteristic polynomials of $(18, 59)$ by $s = \deg H_x$ and elimination mechanism.}\label{tab:18_59_by_s}
\scriptsize
\setlength{\tabcolsep}{4.5pt}
\begin{tabular}{@{}crrrrrr@{}}
\toprule
$s$ & Admissible & LP (Farkas) & $K^\times$ RREF & Prop.~\ref{prop:prop61} & Eulerian / Gal.\ ($66$) & Higher-order ($20$) \\
\midrule
$4$ & $1$ & $1$ & $0$ & $0$ & $0$ & $0$ \\
$5$ & $28$ & $9$ & $10$ & $0$ & $3$ & $6$ \\
$6$ & $103$ & $33$ & $48$ & $2$ & $14$ & $6$ \\
$7$ & $308$ & $90$ & $183$ & $1$ & $26$ & $8$ \\
$8$ & $351$ & $145$ & $188$ & $1$ & $17$ & $0$ \\
$9$ & $330$ & $213$ & $111$ & $0$ & $6$ & $0$ \\
$10$ & $86$ & $74$ & $12$ & $0$ & $0$ & $0$ \\
$11$ & $7$ & $7$ & $0$ & $0$ & $0$ & $0$ \\
\midrule
\textbf{Total} & $\mathbf{1{,}214}$ & $\mathbf{572}$ & $\mathbf{552}$ & $\mathbf{4}$ & $\mathbf{66}$ & $\mathbf{20}$ \\
\bottomrule
\end{tabular}
\end{table}

{\scriptsize
\setlength{\tabcolsep}{3.5pt}
\begin{longtable}{@{}rrr>{\raggedright\arraybackslash}p{2.45in}>{\raggedright\arraybackslash}p{2.35in}@{}}
\caption{Complete disposition of the $86$ remaining candidate spectra for $(18, 59)$ ($\chi_S(x) = (x+5)^{41} Q(x)$; $66$ Eulerian/Galois + $20$ higher-order $\to 0$).}\label{tab:surv_18_59}\\
\toprule
Index & $s$ & $|\Deck|$ & Cofactor $Q(x)$ & Exact Elimination Mechanism \\
\midrule
\endfirsthead
\multicolumn{5}{@{}l}{\scriptsize\emph{Table~\ref{tab:surv_18_59} (continued)}}\\
\toprule
Index & $s$ & $|\Deck|$ & Cofactor $Q(x)$ & Exact Elimination Mechanism \\
\midrule
\endhead
\midrule
\multicolumn{5}{r@{}}{\scriptsize\emph{Continued on next page}}\\
\endfoot
\bottomrule
\endlastfoot
31 & 7 & 1339 & $(x-11)^{13}$\allowbreak $(x-15)$\allowbreak $(x^4-47x^3+805x^2-5957x+16102)$ & Thm.~\ref{thm:galois} Galois orthogonality (Farkas) \\
38 & 6 & 388 & $(x-11)^{12}$\allowbreak $(x-13)^3$\allowbreak $(x^3-34x^2+359x-1198)$ & Thm.~\ref{thm:galois} Galois orthogonality (Farkas) \\
41 & 7 & 174 & $(x-11)^{11}$\allowbreak $(x-13)^2$\allowbreak $(x-15)$\allowbreak $(x-9)^2$\allowbreak $(x^2-25x+142)$ & Thm.~\ref{thm:eulerian} Eulerian $k_u$ / $2$-vertex Jacobi ($6$ sols $\to 0$) \\
48 & 7 & 507 & $(x-11)^{11}$\allowbreak $(x^2-24x+139)^2$\allowbreak $(x^3-36x^2+413x-1502)$ & Thm.~\ref{thm:galois} Galois orthogonality (Farkas) \\
66 & 7 & 106 & $(x-11)^{11}$\allowbreak $(x-13)$\allowbreak $(x^2-22x+113)^2$\allowbreak $(x^2-27x+174)$ & Thm.~\ref{thm:eulerian} Eulerian $k_u$ / $2$-vertex Jacobi ($16$ sols $\to 0$) \\
103 & 6 & 243 & $(x-11)^9$\allowbreak $(x-13)^5$\allowbreak $(x-9)^2$\allowbreak $(x^2-23x+114)$ & $8$ coupled duals ($15\to 7$) + $21$ $2$-vertex Farkas \\
116 & 7 & 1124 & $(x-11)^{11}$\allowbreak $(x-13)^3$\allowbreak $(x^4-45x^3+733x^2-5099x+12778)$ & Thm.~\ref{thm:galois} Galois orthogonality (Farkas) \\
140 & 5 & 93 & $(x-11)^{15}$\allowbreak $(x^3-40x^2+509x-2026)$ & Thm.~\ref{thm:galois} Galois orthogonality (Farkas) \\
142 & 8 & 169 & $(x-11)^{12}$\allowbreak $(x-9)$\allowbreak $(x^5-64x^4+1616x^3-20126x^2+123663x-300034)$ & Thm.~\ref{thm:galois} Galois orthogonality (Farkas) \\
151 & 7 & 1039 & $(x-11)^{13}$\allowbreak $(x-13)$\allowbreak $(x-17)$\allowbreak $(x-9)$\allowbreak $(x^2-23x+122)$ & $2{,}832$ Galois sols $\to 0$ (Prop.~\ref{prop:prop61}) \\
179 & 6 & 139 & $(x-11)^{10}$\allowbreak $(x-13)^4$\allowbreak $(x-9)^2$\allowbreak $(x^2-25x+138)$ & Thm.~\ref{thm:eulerian} Eulerian $k_u$ / $2$-vertex Jacobi ($28$ sols $\to 0$) \\
188 & 6 & 28 & $(x-11)^8$\allowbreak $(x-13)^4$\allowbreak $(x-9)^4$\allowbreak $(x^2-29x+206)$ & Thm.~\ref{thm:eulerian} Eulerian $k_u$ / $2$-vertex Jacobi ($14$ sols $\to 0$) \\
228 & 6 & 480 & $(x-11)^{11}$\allowbreak $(x-13)^4$\allowbreak $(x-9)$\allowbreak $(x^2-23x+110)$ & $12$ coupled + $1$ Prop.~\ref{prop:prop61} + $497$ $2$-vertex + $1$ $S_v$-Galois \\
245 & 7 & 306 & $(x-11)^9$\allowbreak $(x-13)^5$\allowbreak $(x-7)$\allowbreak $(x^3-34x^2+371x-1306)$ & Thm.~\ref{thm:eulerian} Eulerian $k_u$ / $2$-vertex Jacobi ($23$ sols $\to 0$) \\
286 & 8 & 492 & $(x-11)^{10}$\allowbreak $(x-13)^2$\allowbreak $(x-9)^2$\allowbreak $(x^4-51x^3+957x^2-7813x+23354)$ & Thm.~\ref{thm:galois} Galois orthogonality (Farkas) \\
290 & 7 & 446 & $(x-11)^{13}$\allowbreak $(x-15)$\allowbreak $(x^4-47x^3+805x^2-5941x+15958)$ & Thm.~\ref{thm:galois} Galois orthogonality (Farkas) \\
293 & 6 & 304 & $(x-11)^{14}$\allowbreak $(x-13)$\allowbreak $(x^3-38x^2+455x-1658)$ & Thm.~\ref{thm:eulerian} Eulerian $k_u$ / $2$-vertex Jacobi ($1$ sol $\to 0$) \\
316 & 6 & 17 & $(x-10)$\allowbreak $(x-11)^7$\allowbreak $(x-13)^4$\allowbreak $(x-15)^2$\allowbreak $(x-9)^4$ & $948 \to 917$ (Prop.~\ref{prop:prop61}) $\to 468$ $2$-vertex + $449$ $S_v$-Galois \\
319 & 6 & 409 & $(x-11)^{11}$\allowbreak $(x-13)^4$\allowbreak $(x-7)$\allowbreak $(x^2-25x+142)$ & $32$ Galois sols $\to 0$ ($2$-vertex Farkas) \\
341 & 7 & 274 & $(x-11)^{11}$\allowbreak $(x-13)^2$\allowbreak $(x-15)^2$\allowbreak $(x-9)$\allowbreak $(x^2-19x+82)$ & $12$ Galois sols $\to 0$ (Prop.~\ref{prop:prop61} / $2$-vertex) \\
345 & 7 & 1130 & $(x-11)^{11}$\allowbreak $(x-13)^3$\allowbreak $(x-14)$\allowbreak $(x-9)$\allowbreak $(x^2-22x+101)$ & Thm.~\ref{thm:eulerian} Eulerian $k_u$ / $2$-vertex Jacobi ($26$ sols $\to 0$) \\
364 & 5 & 81 & $(x-11)^{12}$\allowbreak $(x-13)^4$\allowbreak $(x^2-21x+86)$ & $7$ Galois sols $\to 0$ ($2$-vertex Farkas) \\
372 & 6 & 49 & $(x-11)^{13}$\allowbreak $(x-17)$\allowbreak $(x-9)^2$\allowbreak $(x^2-27x+178)$ & Thm.~\ref{thm:galois} Galois orthogonality (Farkas) \\
396 & 7 & 1351 & $(x-11)^{13}$\allowbreak $(x-13)$\allowbreak $(x^4-49x^3+873x^2-6675x+18458)$ & Thm.~\ref{thm:galois} Galois orthogonality (Farkas) \\
425 & 8 & 1708 & $(x-11)^{11}$\allowbreak $(x-13)^2$\allowbreak $(x-15)$\allowbreak $(x-9)$\allowbreak $(x^3-34x^2+367x-1262)$ & Thm.~\ref{thm:galois} Galois cone ($0$ $\mathbb{Z}_{\ge 0}$ sols) \\
440 & 8 & 1507 & $(x-11)^{12}$\allowbreak $(x-15)$\allowbreak $(x^5-58x^4+1322x^3-14780x^2+80925x-173378)$ & Thm.~\ref{thm:galois} Galois orthogonality (Farkas) \\
441 & 5 & 19 & $(x-10)$\allowbreak $(x-11)^{13}$\allowbreak $(x-15)^3$\allowbreak $(x-7)$ & $23\to 1$ ($2$-vertex); $E_{15}$ $[18,16,15]$ violates $E_{15}v_{10}=0$ (Prop.~\ref{prop:kill441}) \\
451 & 8 & 303 & $(x-11)^9$\allowbreak $(x-13)^4$\allowbreak $(x^2-20x+95)$\allowbreak $(x^3-34x^2+367x-1274)$ & Thm.~\ref{thm:galois} Galois orthogonality (Farkas) \\
459 & 8 & 337 & $(x-11)^{11}$\allowbreak $(x-13)$\allowbreak $(x-15)^2$\allowbreak $(x-9)$\allowbreak $(x^3-32x^2+329x-1082)$ & Thm.~\ref{thm:eulerian} Eulerian $k_u$ / $2$-vertex Jacobi ($6$ sols $\to 0$) \\
488 & 7 & 186 & $(x-10)$\allowbreak $(x-11)^7$\allowbreak $(x-13)^6$\allowbreak $(x-15)$\allowbreak $(x-7)$\allowbreak $(x-9)^2$ & $468\to 2$ ($2$-vertex); $\rk E_{11}[B,B]=10 > 7$ ($\frac{25}{4}\le\frac{35}{8}$; Prop.~\ref{prop:kill488}) \\
500 & 7 & 543 & $(x-11)^{11}$\allowbreak $(x-13)^3$\allowbreak $(x-15)$\allowbreak $(x^3-30x^2+283x-838)$ & Thm.~\ref{thm:eulerian} Eulerian $k_u$ / $2$-vertex Jacobi ($1$ sol $\to 0$) \\
513 & 6 & 211 & $(x-11)^{13}$\allowbreak $(x-13)^2$\allowbreak $(x^3-36x^2+405x-1446)$ & Thm.~\ref{thm:galois} Galois orthogonality (Farkas) \\
546 & 8 & 1387 & $(x-11)^{11}$\allowbreak $(x-13)^2$\allowbreak $(x-9)$\allowbreak $(x^4-49x^3+877x^2-6779x+19102)$ & Thm.~\ref{thm:galois} Galois orthogonality (Farkas) \\
552 & 7 & 858 & $(x-11)^{13}$\allowbreak $(x-15)$\allowbreak $(x^4-47x^3+805x^2-5937x+15866)$ & Thm.~\ref{thm:galois} Galois orthogonality (Farkas) \\
586 & 9 & 663 & $(x-11)^{10}$\allowbreak $(x-13)^2$\allowbreak $(x-15)$\allowbreak $(x-9)$\allowbreak $(x^4-45x^3+741x^2-5267x+13562)$ & Thm.~\ref{thm:galois} Galois orthogonality (Farkas) \\
587 & 8 & 630 & $(x-11)^9$\allowbreak $(x-13)^4$\allowbreak $(x^2-22x+113)$\allowbreak $(x^3-32x^2+325x-1058)$ & Thm.~\ref{thm:galois} Galois orthogonality (Farkas) \\
592 & 7 & 413 & $(x-10)$\allowbreak $(x-11)^9$\allowbreak $(x-13)^5$\allowbreak $(x-9)$\allowbreak $(x^2-22x+101)$ & Thm.~\ref{thm:galois} Galois orthogonality (Farkas) \\
594 & 7 & 103 & $(x-11)^{10}$\allowbreak $(x-13)^2$\allowbreak $(x^2-21x+94)$\allowbreak $(x^2-24x+139)^2$ & Thm.~\ref{thm:galois} Galois cone ($0$ $\mathbb{Z}_{\ge 0}$ sols) \\
603 & 7 & 139 & $(x-11)^8$\allowbreak $(x-13)^6$\allowbreak $(x-7)$\allowbreak $(x^3-32x^2+329x-1094)$ & $68$ Galois sols $\to 0$ ($2$-vertex Farkas) \\
608 & 6 & 74 & $(x-11)^9$\allowbreak $(x-13)^6$\allowbreak $(x^3-28x^2+245x-682)$ & Thm.~\ref{thm:galois} Galois orthogonality (Farkas) \\
636 & 6 & 56 & $(x-10)$\allowbreak $(x-11)^9$\allowbreak $(x-13)^4$\allowbreak $(x-17)$\allowbreak $(x-9)^3$ & Thm.~\ref{thm:eulerian} Eulerian $k_u$ / $2$-vertex Jacobi ($7$ sols $\to 0$) \\
651 & 5 & 23 & $(x-11)^{13}$\allowbreak $(x-14)$\allowbreak $(x^2-24x+131)^2$ & $5{,}970 \to 5{,}820$ $2$-vertex Farkas + $150$ $S_v$-Galois \\
657 & 8 & 1659 & $(x-11)^{11}$\allowbreak $(x-13)^2$\allowbreak $(x^2-26x+157)$\allowbreak $(x^3-32x^2+329x-1086)$ & Thm.~\ref{thm:galois} Galois orthogonality (Farkas) \\
666 & 7 & 1784 & $(x-11)^{12}$\allowbreak $(x-13)^2$\allowbreak $(x-9)$\allowbreak $(x^3-38x^2+459x-1718)$ & Thm.~\ref{thm:eulerian} Eulerian $k_u$ / $2$-vertex Jacobi ($79$ sols $\to 0$) \\
684 & 9 & 527 & $(x-11)^{11}$\allowbreak $(x-13)$\allowbreak $(x-15)$\allowbreak $(x^5-56x^4+1232x^3-13302x^2+70463x-146434)$ & Thm.~\ref{thm:galois} Galois orthogonality (Farkas) \\
692 & 5 & 21 & $(x-11)^{14}$\allowbreak $(x-15)^2$\allowbreak $(x^2-21x+94)$ & $H=4C_4+P_3$ De Caen--Jacobi $\mathbb{F}_2^4$ ($\sum w_3=0\neq 16$; Prop.~\ref{prop:kill692}) \\
697 & 6 & 161 & $(x-11)^{13}$\allowbreak $(x-13)^2$\allowbreak $(x^3-36x^2+405x-1366)$ & Thm.~\ref{thm:galois} Galois orthogonality (Farkas) \\
705 & 7 & 381 & $(x-11)^9$\allowbreak $(x-13)^4$\allowbreak $(x-14)$\allowbreak $(x-15)$\allowbreak $(x-7)$\allowbreak $(x-9)^2$ & Cone~$0$: $19{,}456$ $2$-vertex + $519$ $S_v$-Galois; cones $1$--$5$: coupled + Prop.~\ref{prop:prop61} \\
716 & 6 & 298 & $(x-11)^{10}$\allowbreak $(x-13)^5$\allowbreak $(x^3-30x^2+279x-818)$ & Thm.~\ref{thm:eulerian} Eulerian $k_u$ / $2$-vertex Jacobi ($2$ sols $\to 0$) \\
724 & 8 & 110 & $(x-11)^7$\allowbreak $(x-13)^6$\allowbreak $(x^2-22x+113)$\allowbreak $(x^3-28x^2+253x-734)$ & Thm.~\ref{thm:galois} Galois orthogonality (Farkas) \\
745 & 8 & 144 & $(x-11)^{11}$\allowbreak $(x-13)$\allowbreak $(x-9)^2$\allowbreak $(x^4-53x^3+1037x^2-8863x+27878)$ & Thm.~\ref{thm:galois} Galois orthogonality (Farkas) \\
778 & 7 & 1465 & $(x-11)^{12}$\allowbreak $(x-13)^2$\allowbreak $(x^4-47x^3+801x^2-5853x+15506)$ & Thm.~\ref{thm:eulerian} Eulerian $k_u$ / $2$-vertex Jacobi ($30$ sols $\to 0$) \\
800 & 7 & 614 & $(x-11)^{11}$\allowbreak $(x-13)$\allowbreak $(x^2-23x+114)$\allowbreak $(x^2-24x+139)^2$ & $41$ coupled $1$-/$2$-vertex Farkas duals ($16\to 0,\; 25\to 0$) \\
825 & 9 & 867 & $(x-11)^9$\allowbreak $(x-13)^3$\allowbreak $(x-9)$\allowbreak $(x^5-58x^4+1322x^3-14784x^2+81045x-174214)$ & Thm.~\ref{thm:galois} Galois orthogonality (Farkas) \\
834 & 6 & 99 & $(x-11)^7$\allowbreak $(x-13)^6$\allowbreak $(x-9)^3$\allowbreak $(x^2-23x+118)$ & Thm.~\ref{thm:eulerian} Eulerian $k_u$ / $2$-vertex Jacobi ($1$ sol $\to 0$) \\
852 & 7 & 169 & $(x-11)^{10}$\allowbreak $(x-13)^4$\allowbreak $(x-15)$\allowbreak $(x-7)$\allowbreak $(x^2-21x+102)$ & Thm.~\ref{thm:galois} Galois orthogonality (Farkas) \\
861 & 6 & 12 & $(x-11)^5$\allowbreak $(x-13)^8$\allowbreak $(x-9)^3$\allowbreak $(x^2-19x+82)$ & $[30,1,28]$, $XY+YZ=8Y \Rightarrow 840 \le 648$ (Prop.~\ref{prop:kill861}; Lean) \\
864 & 5 & 10 & $(x-11)^7$\allowbreak $(x-13)^8$\allowbreak $(x-6)$\allowbreak $(x-9)^2$ & $8\to 3$ ($2$-vertex); $E_9$ $(7,7,7)$ + $63E_{13}$ Cauchy--Schwarz (Prop.~\ref{prop:kill864}) \\
898 & 6 & 179 & $(x-11)^{13}$\allowbreak $(x-13)^2$\allowbreak $(x-17)$\allowbreak $(x^2-19x+82)$ & Thm.~\ref{thm:galois} Galois orthogonality (Farkas) \\
901 & 7 & 1453 & $(x-11)^{12}$\allowbreak $(x-13)^2$\allowbreak $(x^4-47x^3+801x^2-5849x+15430)$ & Thm.~\ref{thm:galois} Galois orthogonality (Farkas) \\
906 & 9 & 3068 & $(x-11)^{10}$\allowbreak $(x-13)^2$\allowbreak $(x-9)$\allowbreak $(x^5-60x^4+1416x^3-16414x^2+93367x-208294)$ & Thm.~\ref{thm:galois} Galois orthogonality (Farkas) \\
909 & 9 & 1552 & $(x-11)^{11}$\allowbreak $(x-13)$\allowbreak $(x-15)$\allowbreak $(x-9)$\allowbreak $(x^4-47x^3+809x^2-6033x+16438)$ & Thm.~\ref{thm:galois} Galois orthogonality (Farkas) \\
953 & 7 & 501 & $(x-11)^{10}$\allowbreak $(x-13)^4$\allowbreak $(x-7)$\allowbreak $(x^3-36x^2+417x-1546)$ & Thm.~\ref{thm:galois} Galois orthogonality (Farkas) \\
974 & 7 & 225 & $(x-11)^9$\allowbreak $(x-13)^5$\allowbreak $(x^4-41x^3+609x^2-3871x+8886)$ & Thm.~\ref{thm:galois} Galois orthogonality (Farkas) \\
978 & 5 & 49 & $(x-11)^{13}$\allowbreak $(x-13)^2$\allowbreak $(x-18)$\allowbreak $(x-9)^2$ & $67 \to 60$ (Prop.~\ref{prop:prop61}) $\to 60$ $2$-vertex Farkas \\
985 & 7 & 804 & $(x-11)^{13}$\allowbreak $(x-15)$\allowbreak $(x^4-47x^3+805x^2-5969x+16186)$ & Thm.~\ref{thm:galois} Galois orthogonality (Farkas) \\
986 & 5 & 67 & $(x-11)^{15}$\allowbreak $(x-17)$\allowbreak $(x^2-23x+118)$ & Thm.~\ref{thm:eulerian} Eulerian $k_u$ / $2$-vertex Jacobi ($1$ sol $\to 0$) \\
997 & 6 & 61 & $(x-11)^{11}$\allowbreak $(x-13)^2$\allowbreak $(x-14)$\allowbreak $(x-17)$\allowbreak $(x-9)^3$ & $58$ Galois sols $\to 0$ (Prop.~\ref{prop:prop61}) \\
998 & 6 & 182 & $(x-11)^{12}$\allowbreak $(x^2-24x+139)^2$\allowbreak $(x^2-25x+138)$ & Thm.~\ref{thm:eulerian} Eulerian $k_u$ / $2$-vertex Jacobi ($11$ sols $\to 0$) \\
1025 & 7 & 200 & $(x-11)^{10}$\allowbreak $(x-13)^4$\allowbreak $(x^4-43x^3+669x^2-4497x+11078)$ & Thm.~\ref{thm:galois} Galois orthogonality (Farkas) \\
1044 & 7 & 1106 & $(x-11)^{12}$\allowbreak $(x-13)^2$\allowbreak $(x-15)$\allowbreak $(x-7)$\allowbreak $(x^2-25x+146)$ & Thm.~\ref{thm:eulerian} Eulerian $k_u$ / $2$-vertex Jacobi ($108$ sols $\to 0$) \\
1059 & 8 & 1893 & $(x-11)^{10}$\allowbreak $(x-13)^3$\allowbreak $(x-9)$\allowbreak $(x^4-47x^3+805x^2-5929x+15794)$ & Thm.~\ref{thm:galois} Galois orthogonality (Farkas) \\
1071 & 8 & 693 & $(x-11)^{10}$\allowbreak $(x-13)^3$\allowbreak $(x^5-56x^4+1228x^3-13158x^2+68771x-139922)$ & Thm.~\ref{thm:galois} Galois orthogonality (Farkas) \\
1092 & 6 & 84 & $(x-11)^{13}$\allowbreak $(x-15)$\allowbreak $(x-9)^2$\allowbreak $(x^2-29x+202)$ & Thm.~\ref{thm:galois} Galois orthogonality (Farkas) \\
1096 & 7 & 226 & $(x-11)^{12}$\allowbreak $(x-13)$\allowbreak $(x-9)^2$\allowbreak $(x^3-42x^2+575x-2558)$ & Thm.~\ref{thm:galois} Galois orthogonality (Farkas) \\
1105 & 8 & 655 & $(x-11)^{12}$\allowbreak $(x-13)$\allowbreak $(x-9)$\allowbreak $(x^4-51x^3+953x^2-7717x+22846)$ & Thm.~\ref{thm:galois} Galois orthogonality (Farkas) \\
1113 & 7 & 1699 & $(x-11)^{13}$\allowbreak $(x-13)$\allowbreak $(x-15)$\allowbreak $(x^3-34x^2+363x-1234)$ & $24$ coupled $1$-/$2$-vertex Farkas duals ($24\to 0$) \\
1116 & 7 & 998 & $(x-11)^{13}$\allowbreak $(x-15)$\allowbreak $(x-9)$\allowbreak $(x^3-38x^2+463x-1786)$ & Thm.~\ref{thm:eulerian} Eulerian $k_u$ / $2$-vertex Jacobi ($18$ sols $\to 0$) \\
1124 & 8 & 1462 & $(x-11)^{11}$\allowbreak $(x-13)^2$\allowbreak $(x-9)$\allowbreak $(x^4-49x^3+877x^2-6783x+19138)$ & Thm.~\ref{thm:galois} Galois orthogonality (Farkas) \\
1158 & 8 & 1979 & $(x-11)^{11}$\allowbreak $(x-13)^2$\allowbreak $(x^5-58x^4+1318x^3-14652x^2+79625x-169162)$ & Thm.~\ref{thm:galois} Galois orthogonality (Farkas) \\
1164 & 7 & 174 & $(x-11)^{10}$\allowbreak $(x-13)^2$\allowbreak $(x^2-22x+113)^2$\allowbreak $(x^2-25x+146)$ & $12$ coupled duals ($18\to 6$) + $1$ $1$-vertex Farkas \\
1179 & 9 & 406 & $(x-11)^{10}$\allowbreak $(x-13)^2$\allowbreak $(x-15)$\allowbreak $(x^2-22x+113)$\allowbreak $(x^3-32x^2+329x-1086)$ & Thm.~\ref{thm:galois} Galois orthogonality (Farkas) \\
1183 & 8 & 482 & $(x-11)^9$\allowbreak $(x-13)^4$\allowbreak $(x-7)$\allowbreak $(x^4-47x^3+813x^2-6133x+17038)$ & Thm.~\ref{thm:galois} Galois orthogonality (Farkas) \\
1202 & 7 & 458 & $(x-11)^{11}$\allowbreak $(x-13)^3$\allowbreak $(x^4-45x^3+733x^2-5079x+12486)$ & Thm.~\ref{thm:galois} Galois orthogonality (Farkas) \\
1211 & 7 & 66 & $(x-11)^7$\allowbreak $(x-13)^6$\allowbreak $(x-9)^2$\allowbreak $(x^3-32x^2+325x-1018)$ & Thm.~\ref{thm:galois} Galois orthogonality (Farkas) \\
1213 & 5 & 31 & $(x-11)^{11}$\allowbreak $(x-13)^5$\allowbreak $(x-14)$\allowbreak $(x-5)$ & Thm.~\ref{thm:eulerian} Eulerian $k_u$ / $2$-vertex Jacobi ($1$ sol $\to 0$) \\
\end{longtable}
}

\subsection{Calibration candidate tables for \texorpdfstring{$(19, 75)$, $(18, 60)$, and $(17, 49)$}{(19, 75), (18, 60), and (17, 49)}}
Table~\ref{tab:cal_benchmarks} records the $8$ $\mathcal{P}_{75,7}$-admissible candidates of $(19, 75)$ (matching~\cite[Prop.~3.3]{GSY2021}), the $5$ candidates of $(18, 60)$ not ruled out by one-vertex LP Farkas duality (out of $44$, matching~\cite{GS2024}), and the $9$ candidates of $(17, 49)$ remaining after one-vertex compatibility (out of $109$ post-Theorem~\ref{thm:pfaffian}, matching~\cite{GSY2023}).

\begin{table}[htbp]
\centering
\caption{Calibration candidates and spectral dispositions for $(19, 75)$, $(18, 60)$, and $(17, 49)$.}\label{tab:cal_benchmarks}
\scriptsize
\setlength{\tabcolsep}{3.5pt}
\begin{tabular}{@{}llr>{\raggedright\arraybackslash}p{2.40in}>{\raggedright\arraybackslash}p{1.90in}@{}}
\toprule
Case & Index & $s$ & Cofactor $Q(x)$ ($\chi_S = (x+5)^{n-d}Q(x)$) & Disposition \\
\midrule
$(19,75)$ & $9$ (unfiltered) & $3$ & $(x-15)^{18}(x-10)$ & $\Leftrightarrow \SRG(75,32,10,16)$~\cite{GSY2021,AzarijaMarc2018} \\
$(19,75)$ & $1,2,5,6,8,10,13$ & $4$--$5$ & $7$ remaining $\mathcal{P}_{75,7}$ polynomials & $6$ $1$-vertex Farkas + $1$ warranted (\cite{GSY2021}) \\
\midrule
$(18,60)$ & $2$ & $5$ & $(x-15)(x-13)^6(x-11)^9(x-9)^2$ & Thm.~\ref{thm:pfaffian} ($n_v=59$) + LP (Farkas) \\
$(18,60)$ & $10$ & $4$ & $(x-13)^9(x-11)^6(x-9)^3$ & Prop.~\ref{prop:prop61} submatrix extraction \\
$(18,60)$ & $17$ & $5$ & $(x-13)^6(x-11)^{10}(x^2-22x+109)$ & $K^\times/(K^\times)^2$ exact RREF \\
$(18,60)$ & $11$ & $4$ & $(x-17)(x-13)^3(x-11)^{14}$ & $60\to 59$ match to $q_{898}$ of $(18,59)$ \\
$(18,60)$ & $34$ & $3$ & $(x-15)^3(x-11)^{15}$ & Regular graph / $3\times 5C_4$ ($\omega = 7 < 9$~\cite{GS2024}) \\
\midrule
$(17,49)$ & $9$ & $3$ & $(x-16)(x-9)^{16}$ & $\Leftrightarrow \SRG(49,16,3,6)$ (\cite{BHMW1989,GSY2023}) \\
$(17,49)$ & $27$ & $6$ & $(x-15)(x-11)^2(x-9)^{12}(x-8)(x-7)$ & Thm.~\ref{thm:eulerian} Eulerian $k_u$ infeasible \\
$(17,49)$ & $43$ & $5$ & $(x-13)^2(x-9)^{13}(x^2-17x+64)$ & $48\times 48$ submatrix elimination~\cite{GSY2023} \\
$(17,49)$ & $55$ & $5$ & $(x-11)^3(x-9)^{12}(x^2-19x+72)$ & $48\times 48$ submatrix elimination~\cite{GSY2023} \\
$(17,49)$ & $58$ & $5$ & $(x-15)(x-12)(x-9)^{14}(x-7)$ & Thm.~\ref{thm:eulerian} Eulerian $k_u$ infeasible \\
$(17,49)$ & $66$ & $6$ & $(x-11)^4(x-9)^{10}(x-7)(x^2-19x+76)$ & Thm.~\ref{thm:galois} Galois orthogonality (Farkas) \\
$(17,49)$ & $70$ & $4$ & $(x-11)^6(x-9)^{10}(x-4)$ & Two-vertex $49\to 48\to 47$ linear Farkas \\
$(17,49)$ & $72$ & $5$ & $(x-11)^2(x-9)^{13}(x^2-21x+92)$ & Exact $\mathrm{diff}=1$ $2$-vertex Jacobi ($2$ sols $\to 0$) \\
$(17,49)$ & $101$ & $6$ & $(x-13)^2(x-11)^2(x-9)^{10}(x-8)(x-7)^2$ & Two-vertex / block elimination~\cite{GSY2023} \\
\bottomrule
\end{tabular}
\end{table}

% ============================================================================
\section{Canonical \texorpdfstring{\texttt{certificates/}}{certificates/} archive manifest}\label{app:cert_manifest}
% ============================================================================

The read-only verifier \nolinkurl{certificates/verify_all.py} validates every JSON bundle in \nolinkurl{certificates/} bit-for-bit against \nolinkurl{certificates/SHA256SUMS.txt} (master SHA-256 of \nolinkurl{SHA256SUMS.txt}: \texttt{d9c57ae07182ef86\allowbreak 1e53093990aa8542\allowbreak 355a3c750098dfe6\allowbreak cde958a64a046d33}; \nolinkurl{verify_all.py} SHA-256: \texttt{92a54996934f37cc\allowbreak 22ad7235bb1ca8a0\allowbreak bf533f9b47943bd5\allowbreak 1648124537ef9774}), confirms zero machine-local path strings, and verifies all integer and rational Farkas duals, $K^\times$ square-class cones, Proposition~\ref{prop:prop61} bounds, Eulerian and Galois-orthogonality certificates, higher-order spectral certificates, and all $12$ calibrated Seidel matrices. Table~\ref{tab:cert_sha256} lists all $56$ JSON artifacts across the five subdirectories.

{\tiny
\setlength{\tabcolsep}{3.5pt}
\begin{longtable}{@{}>{\raggedright\arraybackslash}p{1.90in}r>{\raggedright\arraybackslash}p{3.35in}@{}}
\caption{Complete inventory of \nolinkurl{certificates/} JSON bundles, exact byte sizes, and SHA-256 digests.}\label{tab:cert_sha256}\\
\toprule
Relative Path in \texttt{certificates/} & Size (bytes) & SHA-256 Digest \\
\midrule
\endfirsthead
\multicolumn{3}{@{}l}{\scriptsize\emph{Table~\ref{tab:cert_sha256} (continued)}}\\
\toprule
Relative Path in \texttt{certificates/} & Size (bytes) & SHA-256 Digest \\
\midrule
\endhead
\midrule
\multicolumn{3}{r@{}}{\scriptsize\emph{Continued on next page}}\\
\endfoot
\bottomrule
\endlastfoot
\nolinkurl{calibration/cal_17_48_a5_3_a1_4.json} & 23,193 & \nolinkurl{3d02938df00f687a1cfd74a06ae3403418e76ffa4e817effac8998e548b80c6e} \\
\nolinkurl{calibration/cal_17_48_netto.json} & 23,184 & \nolinkurl{155df64870d8cd43a7323f4a3e7ad6ab761d232803a83c7120a32430eb42406e} \\
\nolinkurl{calibration/cal_18_57_a5_3_a1_4_spec1_111.json} & 32,373 & \nolinkurl{b1aa0587f85f2c97efc0e6ad619e83b356aa8360c2ad8bb669cf9699ea78e08f} \\
\nolinkurl{calibration/cal_18_57_a5_3_a1_4_spec2_11m1.json} & 32,472 & \nolinkurl{bb3d9787d4bff020cc60a9834805ef520a2e6a47d00b6fb9d615a90cb3768a37} \\
\nolinkurl{calibration/cal_18_57_a9_2_a1.json} & 32,598 & \nolinkurl{3a5b42936b4bee9f0d27600cd10a663dd5485870b81aedaaef56259ee4da4359} \\
\nolinkurl{calibration/cal_18_57_a9_2_a1_6ev_first.json} & 32,392 & \nolinkurl{155e46e8827b3007e40ca5d2f547ab7927cde266c9574d78c9c8e65b565b2939} \\
\nolinkurl{calibration/cal_18_57_a9_2_a1_6ev_last.json} & 32,338 & \nolinkurl{713ee14dede48c59d00402b92208ec462710ff0fe1e4b40541ab1c38319eb242} \\
\nolinkurl{calibration/cal_18_57_gsy21_f1.json} & 32,297 & \nolinkurl{7e836d91b15d7d2b9b7f4792a3ccbca89d0b1bd1552f9b132328c5a7284887f2} \\
\nolinkurl{calibration/cal_18_57_gsy21_f2.json} & 32,193 & \nolinkurl{fef0ca0c9721e8d66ffefd65ae2db98f6d52ca0b642b52baf45e4eff240aa957} \\
\nolinkurl{calibration/cal_18_57_gsy21_f3.json} & 32,294 & \nolinkurl{5136f8bb0d1faec4bb32b60ec3f4cf613fb7460b34fc5d6a399dd2c91052fb9b} \\
\nolinkurl{calibration/cal_18_57_gsy21_f4.json} & 32,285 & \nolinkurl{e3fa9a6c566ce0f76aaff3489c27e085009d470702ca6df096260ed2dc827925} \\
\nolinkurl{calibration/cal_19_72_asche.json} & 50,039 & \nolinkurl{ab6279fa145d9a0852008014d90739671ce4d45652687eec3faaf72eb4bf1dcd} \\
\nolinkurl{phase2/deck_17_49.json} & 215,343 & \nolinkurl{e867fe581457ff48c43acd2b38202871f932d6ef1288f869bdc15393317ff61c} \\
\nolinkurl{phase2/deck_18_59.json} & 30,935,619 & \nolinkurl{73ab0cff393b2c27b11251869ac423ad7cafad02ce4681035ea736088a515a98} \\
\nolinkurl{phase2/deck_18_60.json} & 20,975 & \nolinkurl{a70c980245404ef85f926d30159539eff9dc42f08cab8e9157a0863d01267cc4} \\
\nolinkurl{phase2/deck_19_74.json} & 208,029 & \nolinkurl{7c8491bad38eb624647c37e37dfc04e6c0f4f5617fa43c5f3a8e964adf8849e0} \\
\nolinkurl{phase2/phase2_17_49.json} & 266,475 & \nolinkurl{8358a3b6a6c8c7beb2e3b6044228aef476bf56da284f92bb1bf62a4ee10fb858} \\
\nolinkurl{phase2/phase2_18_59.json} & 207,320,936 & \nolinkurl{98ea8a2795da993cfaa78f747aa14149d899953ad0e85f89ffc4dab3dc7f5796} \\
\nolinkurl{phase2/phase2_18_60.json} & 16,047 & \nolinkurl{ead4d6600a508a1fbabc8b36e7bba921f4965676b724d6fc79ac600d6dcbd694} \\
\nolinkurl{phase2/phase2_19_74.json} & 91,554 & \nolinkurl{14f2ac302173880ea563516c2e70c8bed52a87ee582695abe279b3a29b770444} \\
\nolinkurl{phase3/idx13_exact/cert_s1_t2.json} & 4,039 & \nolinkurl{bfe35b52c34af82b9bb51c721993e34abea6fa125526e4f281424c785f982a6a} \\
\nolinkurl{phase3/idx13_exact/cert_s3_s2.json} & 10,627 & \nolinkurl{fb6e927c51786e0331806b43261067773fc0ace7b50feb7dedb1d2d648919996} \\
\nolinkurl{phase3/idx13_exact/s1_4_4_4_4_certs.json} & 122,943,806 & \nolinkurl{a8a7b4af4680ac04c2d315503c3328727dcd2ced03700cd238adcd1cdeae94fd} \\
\nolinkurl{phase3/idx13_exact/s1_6_6_4_certs.json} & 757,222 & \nolinkurl{37eead1b0a40be0762f0ad0abce67662ac3e082d515f2fa313c0f1a178c75235} \\
\nolinkurl{phase3/idx13_exact/s1_8_4_4_certs.json} & 1,461,203 & \nolinkurl{db8d04850b425638a5a82ce7e7e9b9a8496335b9da674cc76b5c99de5f50b383} \\
\nolinkurl{phase3/phase3_18_59_1164.json} & 15,646 & \nolinkurl{be61e8a611387b2a16cfa8cfa70b274a728199492242416deb7f962087275031} \\
\nolinkurl{phase3/phase3_18_59_316.json} & 430,633 & \nolinkurl{ed54aa93bf9bbf947ca2121acbc77e85ba074ad75a60c7dc4feb35a0330c3dd1} \\
\nolinkurl{phase3/phase3_18_59_441_488_692.json} & 1,878 & \nolinkurl{0f1d9263d26013a5fbc7351a5de7a61cf8a700a6ff4ea8ca32cc856c487f4f79} \\
\nolinkurl{phase3/phase3_18_59_705_cone0.json} & 29,506,870 & \nolinkurl{593238c7746544d2958991b8f3beecc7c081449864ca89482165a314defc8684} \\
\nolinkurl{phase3/phase3_18_59_705_cones1_5.json} & 25,761 & \nolinkurl{61b1aa2cf5687b9e78c319fa79b4a63a71c8c94ae5c7f1ee8504fe92960391f6} \\
\nolinkurl{phase3/phase3_18_59_864.json} & 1,930 & \nolinkurl{ed6804c38883b52e2f9f639a5c5c636f50c72abad6dac119dfab50bb511cc085} \\
\nolinkurl{phase3/phase3_18_59_clean_a.json} & 5,517,249 & \nolinkurl{75bed8a2a841bbdc33fc0aa7fb66cdb12e58645f23a3bb601c9042fe5df5eda4} \\
\nolinkurl{phase3/phase3_18_59_clean_b.json} & 4,061,302 & \nolinkurl{a1cd616995d4b58a10c01fa24612a1baf5141d80a55b1b816a151d6b8a18b776} \\
\nolinkurl{phase3/phase3_19_74_47_13_6_24.json} & 39,118 & \nolinkurl{b80f45ec02749c521050a80549f9ed11961ac8c348da7c04ff4446c1ef6ec356} \\
\nolinkurl{phase3/phase3_kills_861_and_47.json} & 803 & \nolinkurl{357974186d63a4bfc772627f50c61deeb225d94d40be24b32115000293ef1d9d} \\
\nolinkurl{sec7/MANIFEST.json} & 660 & \nolinkurl{43ea15f38c13f8712218d1520825e2f5f7e1493faca7acab9644b3d16c39fa77} \\
\nolinkurl{sec7/sec7_17_49.json} & 11,708 & \nolinkurl{c2855c8d72a9422b8b7304110f1a43785a5230bd7efde80e64dfe00faeae8b06} \\
\nolinkurl{sec7/sec7_18_59.json} & 471,596 & \nolinkurl{24a23f128baf4d52546af30d41fe22cf51661c0565028a06409290ad6f3a101b} \\
\nolinkurl{sec7/sec7_18_59_phase3.json} & 1,236,887 & \nolinkurl{d2093372bf91f14e43dea6d757cbf4859c61629266c5390cc54d0c9fbd7e5cb0} \\
\nolinkurl{sec7/sec7_18_60.json} & 385 & \nolinkurl{cb0016c330ee9368bc794bb7325369278b8be428eef003bd422303c0855b5817} \\
\nolinkurl{sec7/sec7_19_74.json} & 1,036,369 & \nolinkurl{61bdb5ee5491530c7658fbb1f3ecb116d86dfad300d96f3e159c16a959ae6079} \\
\nolinkurl{stage1/cor43_17_49.json} & 7,326 & \nolinkurl{50d7234dbffe6b0367725a9021baca1e0af0f58b38e6a293013b6fff5daa34af} \\
\nolinkurl{stage1/cor43_18_59.json} & 86,781 & \nolinkurl{c90cc8f168e378436e5262de0d53ce115dd336a488d8894d898e0b91b07cc37f} \\
\nolinkurl{stage1/mod128_17_49.json} & 12,508 & \nolinkurl{802d641f2fa7a74b173bd801e164aec999989449367e33024e01bcdcc72f0261} \\
\nolinkurl{stage1/mod128_18_60.json} & 3,191 & \nolinkurl{a1f56045845aa83248b4b525da26f751b1e28f8c4e10c5fd938453fcad3ebe1d} \\
\nolinkurl{stage1/mod128_19_74.json} & 12,173 & \nolinkurl{6d1e947b0a3c692c7b7b2cb724c87911c24864eef0894a4026ddbd673ee6595c} \\
\nolinkurl{stage1/mod128_19_75.json} & 620 & \nolinkurl{a5b7d7ed696a1013fd0c58ae36c6de5cd156c5bc4cd7faa72954cf445ae6b1c6} \\
\nolinkurl{stage1/raw_17_49.json} & 46,428 & \nolinkurl{e9409fd07d1caed09a8309007992b04937aaf27bf1b9e8faed400baa54c40bd3} \\
\nolinkurl{stage1/raw_18_60.json} & 3,266 & \nolinkurl{59c54f363ad2eed9417c338429065ff51839ab323dff83dea1a72e77e1b00bc8} \\
\nolinkurl{stage1/raw_19_74.json} & 12,248 & \nolinkurl{ae247d60816db681d0bb6d77e3e76c5c36860b81973303f6e0adcea56adec178} \\
\nolinkurl{stage1/raw_19_75.json} & 1,101 & \nolinkurl{e80f9e8ee01402316f2e630e13b9a9c78579f8393146ebd9be3875271ec5d4c2} \\
\nolinkurl{stage1/table_17_49_cor43.json} & 137,766 & \nolinkurl{9a1dddf3707ed6545bcc69d014bae06b961a8b61470ed44d4a17139146b0ee48} \\
\nolinkurl{stage1/table_18_59_cor43.json} & 1,694,473 & \nolinkurl{f47503a1ac6dc5598a74a9412553281c269ce11a849e791227359d1dfaa1da78} \\
\nolinkurl{stage1/table_18_60.json} & 60,473 & \nolinkurl{67317b5039b10bb45e87c026ed2f74143dfff0253563dd50211869bec509e874} \\
\nolinkurl{stage1/table_19_74.json} & 254,037 & \nolinkurl{40431034d7ca3c3b69f9f0a5316a31b81f6e53be4158af2303f86a0a5aa74216} \\
\nolinkurl{stage1/table_19_75_raw.json} & 20,833 & \nolinkurl{e662939cd2d2325102eb2816d316c4b772bb8fcc259b135b37f94f13720983ed} \\
\end{longtable}
}

% ============================================================================
\section{Complete \texorpdfstring{$127$}{127}-declaration Lean 4 \texorpdfstring{\texttt{\#print axioms}}{print axioms} audit}\label{app:lean_audit}
% ============================================================================

Executing \texttt{lake env lean EquiangularLines/Audit.lean} under \nolinkurl{leanprover/lean4:v4.36.0} verifies all $127$ declarations listed in Table~\ref{tab:lean_127_audit} with zero \texttt{sorry}, zero \nolinkurl{native_decide}, and zero custom axioms.

{\scriptsize
\setlength{\tabcolsep}{3.5pt}
\begin{longtable}{@{}r>{\raggedright\arraybackslash}p{1.2in}>{\raggedright\arraybackslash}p{2.55in}>{\raggedright\arraybackslash}p{1.85in}@{}}
\caption{All $127$ exported Lean~4 declarations verified by \nolinkurl{EquiangularLines/Audit.lean} and their foundational axioms.}\label{tab:lean_127_audit}\\
\toprule
\# & Module & Declaration (\texttt{EquiangularLines.*}) & Verified Axioms (\texttt{\#print axioms}) \\
\midrule
\endfirsthead
\multicolumn{4}{@{}l}{\scriptsize\emph{Table~\ref{tab:lean_127_audit} (continued)}}\\
\toprule
\# & Module & Declaration (\texttt{EquiangularLines.*}) & Verified Axioms (\texttt{\#print axioms}) \\
\midrule
\endhead
\midrule
\multicolumn{4}{r@{}}{\scriptsize\emph{Continued on next page}}\\
\endfoot
\bottomrule
\endlastfoot
1 & \nolinkurl{Basic.lean} & \nolinkurl{isEquiangularSystem_iff_ne_pm} & \texttt{[propext, Classical.choice, Quot.sound]} \\
2 & \nolinkurl{Basic.lean} & \nolinkurl{isEquiangularSystem_iff_distinct_lines} & \texttt{[propext, Classical.choice, Quot.sound]} \\
3 & \nolinkurl{Basic.lean} & \nolinkurl{card_le_dim_of_alpha_zero} & \texttt{[propext, Classical.choice, Quot.sound]} \\
4 & \nolinkurl{Basic.lean} & \nolinkurl{alpha_pos_of_dim_lt} & \texttt{[propext, Classical.choice, Quot.sound]} \\
5 & \nolinkurl{Basic.lean} & \nolinkurl{not_isEquiangularSystem_zero_dim} & \texttt{[propext, Classical.choice, Quot.sound]} \\
6 & \nolinkurl{Basic.lean} & \nolinkurl{IsEquiangularSystem.comp_embedding} & \texttt{[propext, Classical.choice, Quot.sound]} \\
7 & \nolinkurl{Basic.lean} & \nolinkurl{IsEquiangularSystem.sign_flip} & \texttt{[propext, Classical.choice, Quot.sound]} \\
8 & \nolinkurl{Seidel.lean} & \nolinkurl{IsSeidelMatrix.trace_eq_zero} & \texttt{[propext, Classical.choice, Quot.sound]} \\
9 & \nolinkurl{Seidel.lean} & \nolinkurl{IsSeidelMatrix.sq_diag} & \texttt{[propext, Classical.choice, Quot.sound]} \\
10 & \nolinkurl{Seidel.lean} & \nolinkurl{IsSeidelMatrix.trace_sq} & \texttt{[propext, Classical.choice, Quot.sound]} \\
11 & \nolinkurl{Seidel.lean} & \nolinkurl{IsSeidelMatrix.submatrix} & \texttt{[propext, Quot.sound]} \\
12 & \nolinkurl{Seidel.lean} & \nolinkurl{IsSeidelMatrix.switch_isSeidelMatrix} & \texttt{[propext, Classical.choice, Quot.sound]} \\
13 & \nolinkurl{Seidel.lean} & \nolinkurl{IsSeidelMatrix.switch_charpoly} & \texttt{[propext, Classical.choice, Quot.sound]} \\
14 & \nolinkurl{Seidel.lean} & \nolinkurl{seidelGram_isHermitian} & \texttt{[propext, Classical.choice, Quot.sound]} \\
15 & \nolinkurl{Seidel.lean} & \nolinkurl{gramMatrix_eq_transpose_mul_self} & \texttt{[propext, Classical.choice, Quot.sound]} \\
16 & \nolinkurl{Seidel.lean} & \nolinkurl{gramMatrix_posSemidef} & \texttt{[propext, Classical.choice, Quot.sound]} \\
17 & \nolinkurl{Seidel.lean} & \nolinkurl{gramMatrix_rank_le} & \texttt{[propext, Classical.choice, Quot.sound]} \\
18 & \nolinkurl{Seidel.lean} & \nolinkurl{equiangularSystem_to_seidel} & \texttt{[propext, Classical.choice, Quot.sound]} \\
19 & \nolinkurl{Seidel.lean} & \nolinkurl{seidel_to_equiangularSystem} & \texttt{[propext, Classical.choice, Quot.sound]} \\
20 & \nolinkurl{Seidel.lean} & \nolinkurl{hasEquiangularLinesWithAngle_iff_seidel} & \texttt{[propext, Classical.choice, Quot.sound]} \\
21 & \nolinkurl{Seidel.lean} & \nolinkurl{seidelGram_charpoly} & \texttt{[propext, Classical.choice, Quot.sound]} \\
22 & \nolinkurl{Seidel.lean} & \nolinkurl{pow_sub_rank_dvd_charpoly} & \texttt{[propext, Classical.choice, Quot.sound]} \\
23 & \nolinkurl{Seidel.lean} & \nolinkurl{seidelGram_dvd_charpoly_real} & \texttt{[propext, Classical.choice, Quot.sound]} \\
24 & \nolinkurl{Seidel.lean} & \nolinkurl{seidelGram_dvd_charpoly_int} & \texttt{[propext, Classical.choice, Quot.sound]} \\
25 & \nolinkurl{LemmensSeidel.lean} & \nolinkurl{lemmens_seidel_bound} & \texttt{[propext, Classical.choice, Quot.sound]} \\
26 & \nolinkurl{Bounds.lean} & \nolinkurl{gerzon_bound} & \texttt{[propext, Classical.choice, Quot.sound]} \\
27 & \nolinkurl{Bounds.lean} & \nolinkurl{relative_bound_mul} & \texttt{[propext, Classical.choice, Quot.sound]} \\
28 & \nolinkurl{Bounds.lean} & \nolinkurl{relative_bound} & \texttt{[propext, Classical.choice, Quot.sound]} \\
29 & \nolinkurl{Bounds.lean} & \nolinkurl{neumann_theorem} & \texttt{[propext, Classical.choice, Quot.sound]} \\
30 & \nolinkurl{Bounds.lean} & \nolinkurl{phase1_frontier_forces_one_fifth} & \texttt{[propext, Classical.choice, Quot.sound]} \\
31 & \nolinkurl{Spectral.lean} & \nolinkurl{seidel_charpoly_factorization} & \texttt{[propext, Classical.choice, Quot.sound]} \\
32 & \nolinkurl{Spectral.lean} & \nolinkurl{seidel_five_trace_int} & \texttt{[propext, Classical.choice, Quot.sound]} \\
33 & \nolinkurl{Spectral.lean} & \nolinkurl{seidel_five_trace_sq_int} & \texttt{[propext, Classical.choice, Quot.sound]} \\
34 & \nolinkurl{Spectral.lean} & \nolinkurl{seidel_five_trace} & \texttt{[propext, Classical.choice, Quot.sound]} \\
35 & \nolinkurl{Spectral.lean} & \nolinkurl{seidel_five_trace_sq} & \texttt{[propext, Classical.choice, Quot.sound]} \\
36 & \nolinkurl{Spectral.lean} & \nolinkurl{seidel_charpoly_coeff_sub_one} & \texttt{[propext, Classical.choice, Quot.sound]} \\
37 & \nolinkurl{Spectral.lean} & \nolinkurl{two_mul_seidel_charpoly_coeff_sub_two} & \texttt{[propext, Classical.choice, Quot.sound]} \\
38 & \nolinkurl{Spectral.lean} & \nolinkurl{seidel_charpoly_coeff_sub_two} & \texttt{[propext, Classical.choice, Quot.sound]} \\
39 & \nolinkurl{Spectral.lean} & \nolinkurl{shifted_factor_nonneg_roots} & \texttt{[propext, Classical.choice, Quot.sound]} \\
40 & \nolinkurl{Spectral.lean} & \nolinkurl{shifted_factor_trace_coeffs} & \texttt{[propext, Classical.choice, Quot.sound]} \\
41 & \nolinkurl{Spectral.lean} & \nolinkurl{det_add_vecMulVec} & \texttt{[propext, Classical.choice, Quot.sound]} \\
42 & \nolinkurl{Spectral.lean} & \nolinkurl{det_sub_vecMulVec} & \texttt{[propext, Classical.choice, Quot.sound]} \\
43 & \nolinkurl{Spectral.lean} & \nolinkurl{two_dvd_det_of_symm_diag_zero_odd} & \texttt{[propext, Classical.choice, Quot.sound]} \\
44 & \nolinkurl{Spectral.lean} & \nolinkurl{seidel_shifted_charpoly_two_adic_aux} & \texttt{[propext, Classical.choice, Quot.sound]} \\
45 & \nolinkurl{Spectral.lean} & \nolinkurl{seidel_shifted_charpoly_two_adic} & \texttt{[propext, Classical.choice, Quot.sound]} \\
46 & \nolinkurl{Spectral.lean} & \nolinkurl{seidel_shifted_charpoly_two_adic_even} & \texttt{[propext, Classical.choice, Quot.sound]} \\
47 & \nolinkurl{Spectral.lean} & \nolinkurl{two_pow_sub_one_dvd_charpoly_two_mul_sub_one} & \texttt{[propext, Classical.choice, Quot.sound]} \\
48 & \nolinkurl{Spectral.lean} & \nolinkurl{two_pow_factor_dvd_cofactor_comp} & \texttt{[propext, Classical.choice, Quot.sound]} \\
49 & \nolinkurl{Spectral.lean} & \nolinkurl{exists_compl_equiv_pair} & \texttt{[propext, Classical.choice, Quot.sound]} \\
50 & \nolinkurl{Spectral.lean} & \nolinkurl{jacobi_complementary_minor} & \texttt{[propext, Classical.choice, Quot.sound]} \\
51 & \nolinkurl{Spectral.lean} & \nolinkurl{charmatrix_adjugate_diag} & \texttt{[propext, Classical.choice, Quot.sound]} \\
52 & \nolinkurl{Spectral.lean} & \nolinkurl{derivative_charpoly_eq_sum_adjugate_diag} & \texttt{[propext, Classical.choice, Quot.sound]} \\
53 & \nolinkurl{Spectral.lean} & \nolinkurl{derivative_charpoly_eq_sum_vertex_deletions} & \texttt{[propext, Classical.choice, Quot.sound]} \\
54 & \nolinkurl{Spectral.lean} & \nolinkurl{jacobi_two_vertex_charpoly} & \texttt{[propext, Classical.choice, Quot.sound]} \\
55 & \nolinkurl{Spectral.lean} & \nolinkurl{cauchy_interlacing} & \texttt{[propext, Classical.choice, Quot.sound]} \\
56 & \nolinkurl{TwoAdic} & \nolinkurl{even_order_two_pow_dvd_charpoly} & \texttt{[propext, Classical.choice, Quot.sound]} \\
57 & \nolinkurl{TwoAdic} & \nolinkurl{skew_det_eq_sq} & \texttt{[propext, Classical.choice, Quot.sound]} \\
58 & \nolinkurl{TwoAdic} & \nolinkurl{odd_order_two_adic_congruence} & \texttt{[propext, Classical.choice, Quot.sound]} \\
59 & \nolinkurl{Multiplicity} & \nolinkurl{even_order_cofactor_two_pow_dvd} & \texttt{[propext, Classical.choice, Quot.sound]} \\
60 & \nolinkurl{Multiplicity} & \nolinkurl{cofactor_comp_two_mul_add_odd_even} & \texttt{[propext, Classical.choice, Quot.sound]} \\
61 & \nolinkurl{Multiplicity} & \nolinkurl{cofactor_comp_two_mul_add_odd_odd} & \texttt{[propext, Classical.choice, Quot.sound]} \\
62 & \nolinkurl{Multiplicity} & \nolinkurl{odd_order_cofactor_two_adic_congruence} & \texttt{[propext, Classical.choice, Quot.sound]} \\
63 & \nolinkurl{Multiplicity} & \nolinkurl{seidel_cofactor_trace_diff_eq} & \texttt{[propext, Classical.choice, Quot.sound]} \\
64 & \nolinkurl{Multiplicity} & \nolinkurl{gsy_multiplicity_bound_even} & \texttt{[propext, Classical.choice, Quot.sound]} \\
65 & \nolinkurl{Multiplicity} & \nolinkurl{gsy_multiplicity_bound_odd} & \texttt{[propext, Classical.choice, Quot.sound]} \\
66 & \nolinkurl{Multiplicity} & \nolinkurl{frontier_18_59_mult_11} & \texttt{[propext, Classical.choice, Quot.sound]} \\
67 & \nolinkurl{Multiplicity} & \nolinkurl{frontier_19_74_mult_15} & \texttt{[propext, Classical.choice, Quot.sound]} \\
68 & \nolinkurl{Multiplicity} & \nolinkurl{frontier_17_49_mult_9} & \texttt{[propext, Classical.choice, Quot.sound]} \\
69 & \nolinkurl{Multiplicity} & \nolinkurl{frontier_18_60_mult_11} & \texttt{[propext, Classical.choice, Quot.sound]} \\
70 & \nolinkurl{SRG136} & \nolinkurl{isSRG_of_matrix_eq} & \texttt{[propext, Classical.choice, Quot.sound]} \\
71 & \nolinkurl{SRG136} & \nolinkurl{isSRG_matrix_eq} & \texttt{[propext, Classical.choice, Quot.sound]} \\
72 & \nolinkurl{SRG136} & \nolinkurl{idx136_implies_SRG_75_32_10_16} & \texttt{[propext, Classical.choice, Quot.sound]} \\
73 & \nolinkurl{SRG136} & \nolinkurl{SRG_75_32_10_16_implies_idx136} & \texttt{[propext, Classical.choice, Quot.sound]} \\
74 & \nolinkurl{SRG136} & \nolinkurl{idx136_eliminated_of_not_existsSRG} & \texttt{[propext, Classical.choice, Quot.sound]} \\
75 & \nolinkurl{Obstructions} & \nolinkurl{even_order_seidel_zmod2_charpoly} & \texttt{[propext, Classical.choice, Quot.sound]} \\
76 & \nolinkurl{Obstructions} & \nolinkurl{odd_order_seidel_zmod2_charpoly} & \texttt{[propext, Classical.choice, Quot.sound]} \\
77 & \nolinkurl{Obstructions} & \nolinkurl{even_eig_mult_even_order_eq_zero} & \texttt{[propext, Classical.choice, Quot.sound]} \\
78 & \nolinkurl{Obstructions} & \nolinkurl{even_eig_mult_odd_order_le_one} & \texttt{[propext, Classical.choice, Quot.sound]} \\
79 & \nolinkurl{Obstructions} & \nolinkurl{exists_hensel_root_mod_two_pow} & \texttt{[propext, Classical.choice, Quot.sound]} \\
80 & \nolinkurl{Obstructions} & \nolinkurl{even_root_mod_four_eq_sub_one} & \texttt{[propext, Classical.choice, Quot.sound]} \\
81 & \nolinkurl{Obstructions} & \nolinkurl{strict_type2_eval_mod_eight} & \texttt{[propext, Classical.choice, Quot.sound]} \\
82 & \nolinkurl{Obstructions} & \nolinkurl{weapon_b_vertex_deletion_mod_eight} & \texttt{[propext, Classical.choice, Quot.sound]} \\
83 & \nolinkurl{Obstructions} & \nolinkurl{weapon_b_vertex_deletion_odd} & \texttt{[propext, Classical.choice, Quot.sound]} \\
84 & \nolinkurl{Obstructions} & \nolinkurl{weapon_b_padicVal_zero} & \texttt{[propext, Classical.choice, Quot.sound]} \\
85 & \nolinkurl{Obstructions} & \nolinkurl{weapon_b_derivative_mod_eight} & \texttt{[propext, Classical.choice, Quot.sound]} \\
86 & \nolinkurl{Obstructions} & \nolinkurl{jacobi_square_class_mod_charpoly} & \texttt{[propext, Classical.choice, Quot.sound]} \\
87 & \nolinkurl{Obstructions} & \nolinkurl{psd_submatrix_kernel_extends} & \texttt{[propext, Classical.choice, Quot.sound]} \\
88 & \nolinkurl{Obstructions} & \nolinkurl{ker_submatrix_vanish_of_nullity_gt} & \texttt{[propext, Classical.choice, Quot.sound]} \\
89 & \nolinkurl{GaloisObstructions} & \nolinkurl{exists_unique_eulerian_switch} & \texttt{[propext, Classical.choice, Quot.sound]} \\
90 & \nolinkurl{GaloisObstructions} & \nolinkurl{eulerian_switch_sq_mod_four} & \texttt{[propext, Classical.choice, Quot.sound]} \\
91 & \nolinkurl{GaloisObstructions} & \nolinkurl{simple_eig_projector_adjugate} & \texttt{[propext, Classical.choice, Quot.sound]} \\
92 & \nolinkurl{GaloisObstructions} & \nolinkurl{simple_eig_projector_diag} & \texttt{[propext, Classical.choice, Quot.sound]} \\
93 & \nolinkurl{GaloisObstructions} & \nolinkurl{galois_orthogonality_of_orth_vectors} & \texttt{[propext, Classical.choice, Quot.sound]} \\
94 & \nolinkurl{GaloisObstructions} & \nolinkurl{galois_orthogonality_trace_constraint} & \texttt{[propext, Classical.choice, Quot.sound]} \\
95 & \nolinkurl{GaloisObstructions} & \nolinkurl{ker_submatrix_preserves_nullity} & \texttt{[propext, Classical.choice, Quot.sound]} \\
96 & \nolinkurl{GaloisObstructions} & \nolinkurl{submatrix_nullity_ge_sub} & \texttt{[propext, Classical.choice, Quot.sound]} \\
97 & \nolinkurl{GaloisObstructions} & \nolinkurl{submatrix_nullity_ge_of_vertex_deletions} & \texttt{[propext, Classical.choice, Quot.sound]} \\
98 & \nolinkurl{GaloisObstructions} & \nolinkurl{submatrix_eigenvalue_count_eq_finrank_ker} & \texttt{[propext, Classical.choice, Quot.sound]} \\
99 & \nolinkurl{GaloisObstructions} & \nolinkurl{submatrix_trace_cauchy_schwarz_obstruction} & \texttt{[propext, Classical.choice, Quot.sound]} \\
100 & \nolinkurl{Phase3Algebra} & \nolinkurl{block_commutator_sim_eig} & \texttt{[propext, Classical.choice, Quot.sound]} \\
101 & \nolinkurl{Phase3Algebra} & \nolinkurl{symm_block_quadratic_eigenvalues} & \texttt{[propext, Classical.choice, Quot.sound]} \\
102 & \nolinkurl{Phase3Algebra} & \nolinkurl{idx47_eliminated} & \texttt{[propext, Classical.choice, Quot.sound]} \\
103 & \nolinkurl{Phase3Algebra} & \nolinkurl{idx861_eliminated} & \texttt{[propext, Classical.choice, Quot.sound]} \\
104 & \nolinkurl{Phase3Algebra} & \nolinkurl{jacobi_k_complementary_minor} & \texttt{[propext, Classical.choice, Quot.sound]} \\
105 & \nolinkurl{Phase3Algebra} & \nolinkurl{idx13_minimal_poly_of_charpoly} & \texttt{[propext, Classical.choice, Quot.sound]} \\
106 & \nolinkurl{Phase3Algebra} & \nolinkurl{idx13_resolvent_identity} & \texttt{[propext, Classical.choice, Quot.sound]} \\
107 & \nolinkurl{Phase3Algebra} & \nolinkurl{idx13_resolvent_at_15} & \texttt{[propext, Classical.choice, Quot.sound]} \\
108 & \nolinkurl{Phase3Algebra} & \nolinkurl{idx13_pole_cancellation_at_15} & \texttt{[propext, Classical.choice, Quot.sound]} \\
109 & \nolinkurl{Phase3Algebra} & \nolinkurl{pairwise_orth_linearIndependent} & \texttt{[propext, Classical.choice, Quot.sound]} \\
110 & \nolinkurl{Phase3Algebra} & \nolinkurl{disjoint_eigenvectors_rank_ge} & \texttt{[propext, Classical.choice, Quot.sound]} \\
111 & \nolinkurl{Phase3Algebra} & \nolinkurl{hermitian_trace_sq_le_rank_mul_trace_mul} & \texttt{[propext, Classical.choice, Quot.sound]} \\
112 & \nolinkurl{Phase3Algebra} & \nolinkurl{idx488_rank_obstruction} & \texttt{[propext, Classical.choice, Quot.sound]} \\
113 & \nolinkurl{Certs / Certificates} & \nolinkurl{Certs.Certs19_74A_all_ok} & \texttt{[propext]} \\
114 & \nolinkurl{Certs / Certificates} & \nolinkurl{Certs.Certs19_74B_all_ok} & \texttt{[propext]} \\
115 & \nolinkurl{Certs / Certificates} & \nolinkurl{Certs.Certs19_74C_all_ok} & \texttt{[propext]} \\
116 & \nolinkurl{Certs / Certificates} & \nolinkurl{farkas_infeasible_nonneg} & \texttt{[propext, Classical.choice, Quot.sound]} \\
117 & \nolinkurl{Certs / Certificates} & \nolinkurl{farkas_infeasible_box} & \texttt{[propext, Classical.choice, Quot.sound]} \\
118 & \nolinkurl{Certs / Certificates} & \nolinkurl{polyOfDescList_coeff} & \texttt{[propext, Classical.choice, Quot.sound]} \\
119 & \nolinkurl{Certs / Certificates} & \nolinkurl{checkCertZ_sound} & \texttt{[propext, Classical.choice, Quot.sound]} \\
120 & \nolinkurl{Certs / Certificates} & \nolinkurl{charpoly_eliminated_of_deck_farkas} & \texttt{[propext, Classical.choice, Quot.sound]} \\
121 & \nolinkurl{Certs / Certificates} & \nolinkurl{certs19_74_all_147_verified} & \texttt{[propext, Quot.sound]} \\
122 & \nolinkurl{Main} & \nolinkurl{frontier_19_74_reduction} & \texttt{[propext, Classical.choice, Quot.sound]} \\
123 & \nolinkurl{Main} & \nolinkurl{frontier_18_59_reduction} & \texttt{[propext, Classical.choice, Quot.sound]} \\
124 & \nolinkurl{Main} & \nolinkurl{frontier_19_74_core_cases_eliminated} & \texttt{[propext, Classical.choice, Quot.sound]} \\
125 & \nolinkurl{Main} & \nolinkurl{frontier_18_59_core_cases_eliminated} & \texttt{[propext, Classical.choice, Quot.sound]} \\
126 & \nolinkurl{Main} & \nolinkurl{frontier_19_74_no_lines_of_srg_and_cert_exhaustion} & \texttt{[propext, Classical.choice, Quot.sound]} \\
127 & \nolinkurl{Main} & \nolinkurl{frontier_18_59_no_lines_of_cert_exhaustion} & \texttt{[propext, Classical.choice, Quot.sound]} \\
\end{longtable}
}

\section*{Acknowledgments}
The author thanks Sri Chaitanya Educational Institutions for institutional support, and gratefully acknowledges the open-source contributors to the Lean~4 theorem prover, the Mathlib library, and the C++ and Python scientific software ecosystem.

\section*{Data and code availability}
The Lean~4 formalization, the exact C++20 and Python verification programs, and the frozen $56$-file exact certificate archive (with canonical manifest \nolinkurl{certificates/SHA256SUMS.txt}, SHA-256 \nolinkurl{d9c57ae07182ef861e53093990aa8542355a3c750098dfe6cde958a64a046d33}) are openly archived at \url{https://github.com/Mehardeep-Singh-GG/equiangular-lines-18-19} (release \texttt{v1.0.0}). Large JSON certificate bundles are stored compressed (\texttt{.json.gz}) in the public release, while \nolinkurl{certificates/verify_all.py} verifies their decompressed contents bit-for-bit against the canonical \nolinkurl{SHA256SUMS.txt} digests.

\bibliographystyle{amsplain}
\bibliography{refs}

\end{document}
